 \documentclass{article}
 \usepackage{amsmath}
 \usepackage{amsthm}
 \usepackage{amssymb}
 \usepackage{pgfplots}
 \usepackage{amsfonts}
 \usepackage[margin=2.5cm]{geometry}
 \usepackage{graphicx}
 \usepackage[export]{adjustbox}
 \usepackage{fancyhdr}
 \usepackage[portuguese]{babel}
 \usepackage{hyperref}
 \usepackage{lastpage}
 \usepackage{mathtools}
 \usetikzlibrary{patterns}

 \pagestyle{fancy}
 \fancyhf{}

 \pgfplotsset{compat = 1.18}

 \hypersetup{
     colorlinks,
     citecolor=black,
     filecolor=black,
     linkcolor=black,
     urlcolor=black
 }
 \newtheorem*{def*}{\underline{Defini\c c\~ao}}
 \newtheorem*{theorem*}{\underline{Teorema}}
 \newtheorem*{lemma*}{\underline{Lema}}
 \newtheorem*{prop*}{\underline{Proposi\c c\~ao}}
 \newtheorem{example}{\underline{Exemplo}}
 \newtheorem*{proof*}{\underline{Prova}}
 \newtheorem*{crl*}{\underline{Corol\'ario}}
 \renewcommand\qedsymbol{$\blacksquare$}
 \newcommand{\Lin}[1]{Lin_{\mathbb{K}}({#1})}

 \rfoot{P\'agina \thepage \hspace{1pt} de \pageref{LastPage}}

 \title{Notas de An\'alise}
 \author{Renan Wenzel}
 \date{\today}

 \begin{document}
 \begin{figure}[ht]
	\minipage{0.76\textwidth}
		\includegraphics[width=4cm]{../icmc.png}
		\hspace{7cm}
		\includegraphics[height=4.9cm,width=4cm]{../brasao_usp_cor.jpg}
	\endminipage	
\end{figure}

\begin{center}
	\vspace{1cm}
	\LARGE
	UNIVERSIDADE DE S\~AO PAULO

	\vspace{1.3cm}
	\LARGE
	INSTITUTO DE CI\^ENCIAS MATEM\'ATICAS E COMPUTACIONAIS - ICMC

	\vspace{1.7cm}
	\Large
	\textbf{Notas de Aula de An\'alise}

	\vspace{1.3cm}
	\large
	\textbf{Renan Wenzel - 11169472}

	\vspace{1.3cm}
	\large
	\textbf{Alexandre Nolasco de Carvalho - andcarva@icmc.usp.br}

	\vspace{1.3cm}
	\today
\end{center}

 \newpage

 \tableofcontents

 \newpage

\section{Aula 01 - 13/03/2023}
\subsection{Motiva\c c\~ao}
\begin{itemize}
  \item Relembrar sistemas b\'asicos da matem\'atica;
  \item Relembrar propriedades b\'asicas das principais estruturas ($\mathbb{N}, \mathbb{Z}, \mathbb{Q}$).
\end{itemize}

\subsection{Os N\'umeros Naturais}
  Os n\'umeros naturais s\~ao os que utilizamos para contar objetos, e s\~ao caracterizados pelos Axiomas de Peano:
 \begin{itemize}
   \item[1)] Todo n\'umero natural tem um \'unico sucessor; 
   \item[2)] N\'umeros naturais diferentes t\^em sucessores diferentes;
   \item[3)] Existe um \'unico n\'umero natural, zero (0), que n\~ao \'e sucessor de nenhum n\'umero natural.
   \item[4)] Seja $X \subseteq{\mathbb{N}}$ tal que $0\in{X}$ e, se n pertence a X, seu sucessor n+1 tamb\'em pertence 
a X. Ent\~ao, X = $\mathbb{N}.$ (Propriedade de Indu\c c\~ao).
 \end{itemize}

\begin{def*}
  Definimos a adi\c c\~ao por: $n + 0 = n, n\in \mathbb{N},\text{ e }n+(p+1) = (n+p)+1, p\in{\mathbb{N}}$. Al\'em disso,
a multiplica\c c\~ao \'e dada por: $n.0 = 0, n.(p+1) = n.p + n, n, p\in\mathbb{N}.$ Ou seja, sabendo somar ou multiplicar um n\'umero,
sabemos somar e multiplicar seu sucessor.
\end{def*}
  Com rela\c c\~ao ao quarto axioma, ele leva este nome porque um dos m\'etodos de demonstra\c c\~ao, conhecido como
prova por indu\c c\~ao. Nele, mostramos um caso base, o caso 0, e utilizamos a segunda parte para provar que, se um
resultado vale para o caso n, ele vale para n+1, portanto sendo verdadeiro para todos os naturais.

\begin{lemma*}
  Para todo n natural, 1 + n = n + 1.
\end{lemma*}
\begin{proof*}
  Note que o resultado \'e verdadeiro para n = 0. Suponha que o resultado seja v\'alido para n = k e mostremos que 
vale tamb\'em para n = k+1. Com efeito, segue pela propriedade de indu\c c\~ao e pela defini\c c\~ao de soma que
 $$
    1 + (k + 1) = (1 + k) + 1 =  (k + 1) + 1. 
 $$
 Segue que o resultado vale para todo n natural. \qedsymbol
\end{proof*}
  A seguir, mostramos a associatividade e a comutatividade, respectivamente, das opera\c c\~oes nos naturais.
\begin{lemma*}
  Para todo n, p, r naturais, (n + p) + r = n + (p + r).
\end{lemma*}
\begin{proof*}
  Note que o resultado \'e v\'alido trivialmente para r = 0 e r = 1. Suponha que o resultado seja v\'alido para
r = k e mostremos que vale tamb\'em para r = k + 1. Com efeito, pela hip\'otese de indu\c c\~ao e defini\c c\~ao de adi\c c\~ao,
  $$
    n + (p + (k + 1)) = n + ((p + k) + 1) = (n + (p + k)) + 1 = ((n + p) + k) + 1 = (n + p) + (k + 1).
  $$
  Segue o resultado por indu\c c\~ao. \qedsymbol
\end{proof*}
\begin{lemma*}
  Para todo n, p naturais, n + p = p + n. 
\end{lemma*}
\begin{proof*}
  Observe que j\'a mostramos o caso em que p = 1. Suponha que o resultado vale para p = k e vamos mostrar o caso
p = k + 1. De fato, pela hip\'otese de indu\c c\~ao e defini\c c\~ao de adi\c c\~ao, junto do lema de associatividade,
temos  
  $$
    n + (k + 1) = (n + k) + 1 = (k + n) + 1 = 1 + (k + n) = (1 + k) + n = (k + 1) + n.
  $$
  Por indu\c c\~ao, segue que isso vale para todo natural n. \qedsymbol
\end{proof*}

\begin{def*}
  Definimos uma ordem em $\mathbb{N}$ colocando que $m\leq{n}$ se existe p natural tal que $n = m + p. \square$
\end{def*}
  A rela\c c\~ao de ordem possui as seguintes propriedades:
 \begin{itemize}
   \item[i)] Reflexiva: Para todo n natural, $n\leq{n};$
   \item[ii)] Antissim\'etrica: Se $m\leq n$ e $n\leq m,$ ent\~ao $m = n;$
   \item[iii)] Transitiva: Se $m \leq n$ e $n \leq p$, ent\~ao $m\leq p;$
   \item[i] Dados m, n naturais, temos ou $m \leq n$, ou $n \leq m;$
   \item[v] Se $m \leq n$ e p \'e um natural, ent\~ao $n + p\leq n\text{ e } mp\leq np$
 \end{itemize}

 \subsection{N\'umeros Inteiros e Racionais}
  Usualmente, construimos os inteiros a partir dos naturais tomando os pares ordenados de n\'umeros naturais
com a seguinte identifica\c c\~ao (a, b) $\mathtt{\sim}$ (c, d) se a + d = b + c. Assim, podemos representar
  $$
  \mathbb{N} = \{(0, 0), (1, 0), (2, 0), (3, 0), \ldots\}, \quad -\mathbb{N}^* = \{\cdots, (0, 3), (0, 2), (0, 1)\}.
  $$
  Tomar o sucessor ser\'a somar 1 \`a primeira coordenada e, para os inteiros negativos, voltar a identificar (1, n) com (0, n-1).

  Os n\'umeros racionais s\~ao constru\'idos tomando o conjunto $\mathbb{Z}\times{\mathbb{Z}^*}$ e identificando os pares $(a, b)\mathtt{\sim}(c, d)$
para os quais ad = bc. Representamos um par (a, b) neste conjunto por $\displaystyle \frac{a}{b}.$ A soma e o produto em $\mathbb{Q}$
s\~ao dados, respectivamente, por: 
 \begin{align*}
   &\frac{a}{b} + \frac{c}{d} \coloneqq \frac{ad + bc}{bd} \\
   &\frac{a}{b}\cdot\frac{c}{d} \coloneqq \frac{ac}{bd}.
 \end{align*}
 Chamamos a adi\c c\~ao a opera\c c\~ao que a cada par $(x, y)\in \mathbb{Q}\times{\mathbb{Q}}$ associa sua soma $x+y\in \mathbb{Q}$
e chamamos multiplica\c c\~ao a opera\c c\~ao que a cada par $(x, y)\in \mathbb{Q}\times \mathbb{Q}$ associa seu produto $x.y\in \mathbb{Q}.$
A terna $(\mathbb{Q}, +, \cdot)$ satisfaz as propriedades de um corpo, i.e., 
 \begin{align*}
   &(A1) (x + y) + z = x + (y + z), \quad\forall x, y, z\in \mathbb{Q}\\
   &(A2) x + y = y + x, \quad\forall x, y\in \mathbb{Q}\\
   &(A3) \exists 0\in \mathbb{Q}: x + 0 = x, \quad\forall x\in \mathbb{Q}\\
   &(A4) \forall x\in \mathbb{Q}, \exists y\in \mathbb{Q} (y = -x): x + y = 0\\
   &(M1) (xy)z = x(yz), \quad\forall x, y, z\in \mathbb{Q}\\
   &(M2) xy = yx, \quad x, y\in \mathbb{Q}\\
   &(M3) \exists 1\in \mathbb{Q}: 1.x = x.1 = x, \quad\forall x\in \mathbb{Q}\\
   &(M4) \forall x\in \mathbb{Q}^*, \exists y = \frac{1}{x}\in \mathbb{Q}: x.y = 1\\
   &(D) x(y+z) = xy + xz,\quad\forall x, y, z\in \mathbb{Q}.
 \end{align*}
 \newpage

\section{Aula 02 - 15/03/2023}
\subsection{Motiva\c c\~oes}
\begin{itemize}
  \item Propriedades b\'asicas dos racionais;
  \item Constru\c c\~ao do corpo dos reais a partir dos racionais;
  \item Cortes de Dedekind.
\end{itemize}
\subsection{Propriedades de $\mathbb{Q}$ e sua Ordem}
  Com as 9 propriedades de corpo, conseguimos obter novas regras nos racionais, como a famosa lei do cancelamento:
 \begin{prop*}
   Em $\mathbb{Q},$ vale
   $$
    x + z = y + z \Rightarrow x = y
   $$
   e, se $z\neq{0}$,
   $$
    xz = yz \Rightarrow x = y 
   $$
 \end{prop*}
 \begin{proof*}
   \begin{align*}
   &x = x + 0 = x + (z + (-z)) = (x+z) + (-z) = (y + z) + (-z) = y + (z + (-z)) = y +0 = y\\
   &x = x.1 = x(z.\frac{1}{z}) = (xz)\frac{1}{z} = (yz)\frac{1}{z} = y (z \frac{1}{z}) = y.1 = y. \text{ \qedsymbol}
 \end{align*}
 \end{proof*}
 \begin{prop*}
  Os elementos neutros da adi\c c\~ao e multiplica\c c\~ao s\~ao \'unicos. Os elementos oposto e inverso tamb\'em o s\~ao.
 \end{prop*}
 \begin{prop*}
   Para todo x racional, x.0 = 0.
 \end{prop*}
 \begin{prop*}
   Para todo x racional, -x = (-1)x.
 \end{prop*}
 A maioria desses resultados acima seguem diretamente da lei do cancelamento. Suas demonstra\c c\~oes ficam como exerc\'icio.
\begin{def*}
  Diremos que 
  $$
    \frac{a}{b}\in \mathbb{Q} = \left\{\begin{array}{ll}
        \text{ n\~ao-negativo, } \quad ab\in \mathbb{N}\\
        \text{ positivo, } \quad ab\in \mathbb{N}, a\neq 0
      \end{array}\right.
  $$
  e diremos que 
  $$
    \frac{a}{b}\in \mathbb{Q} = \left\{\begin{array}{ll}
        \text{ n\~ao-positivo, } \quad \frac{a}{b} \text{ n\~ao for postivo}\\
        \text{ negativo, } \quad \frac{a}{b} \text{ n\~ao for n\~ao-negativo.}
      \end{array}\right.\square
  $$
\end{def*}
\begin{def*}
  Sejam x, y racionais. Diremos que x \'e menor e que y e escrevemos ``$x < y$'' se existir t racional positivo tal que 
  $$
    y = x + t.
  $$
  Neste mesmo caso, podemos dizer que y \'e maior que x, escrevendo ``$x > y$''. Em particular, temos $x > 0$ se x for positivo e 
$x < 0$ se x for negativo. 

  Ademais, se $x < y$ ou x = y, escrevemos ``$x \leq{y}$'' se existir racional t n\~ao-negativo tal que 
  $$
    y = x + t
  $$
  e, se $x > y$ ou x = y, escrevemos ``$x \geq{y}$'' caso exista racional t n\~ao-positivo com
  $$
    y = x + t. \square
  $$
\end{def*}
  A qu\'adrupla ($\mathbb{Q}, +, \cdot, \leq{}$) satisfaz as propriedades de um corpo ordenado, i.e., 
 \begin{align*}
   &(O1) x \leq{x}\forall x\in \mathbb{Q};\\
   &(O2) x\leq{y} \text{ e } y \leq{x}\Rightarrow x = y \forall x, y\in \mathbb{Q};\\
   &(O3) x \leq{y}, y \leq{z}\Rightarrow x \leq{z}\forall x, y, z\in \mathbb{Q};\\
   &(O4)\forall x, y \in \mathbb{Q}, x \leq{y} \text{ ou } y \leq{x};\\
   &(OA) x \leq{y}\Rightarrow x + z \leq{y + z};\\
   &(OM) x \leq{y} \text{ e } z \geq{0}\Rightarrow xz \leq{yz}. 
 \end{align*}
 \begin{prop*}
  Para quaisquer x, y, z, w no corpo ordenado dos racionais, valem 
 \begin{align*}
   &i.) x < y\Longleftrightarrow x + z < y + z \\
   &ii.) z > 0\Longleftrightarrow \frac{1}{z} > 0 \\
   &iii.) z > 0\Longleftrightarrow -z < 0 \\
   &iv.) z > 0\Rightarrow x < y\Longleftrightarrow xz < yz \\
   &v.) z < 0\Rightarrow x < y\Longleftrightarrow xz > yz \\
   &vi.) xz < yw\Longleftrightarrow \left\{\begin{array}{ll}
       0 \leq{x} < y \\
       0 \leq{z} < w
     \end{array}\right. \\
   &vii.) 0 < x < y\Longleftrightarrow 0 < \frac{1}{y} < \frac{1}{x}\\
   &viii.) x < y \text{ ou } x =y \text{ ou } x > y \\
   &ix.) xy = 0\Longleftrightarrow x = 0\text{ ou }y = 0. \\
   &x.) \left.\begin{array}{ll}
       x \leq{y} \\
       z \leq{w}
     \end{array}\right\}\Rightarrow x + z \leq{y + w}\\
   &xi.) \left.\begin{array}{ll}
       0 \leq{x} \leq{y}\\
       0 \leq{z} \leq{w}
   \end{array}\right\}\Rightarrow xz \leq{yw}.
 \end{align*}
 \end{prop*}
 \subsection{Incompletude de $\mathbb{Q}$}
  Os n\'umeros racionais podem ser representados por pontos em uma reta horizontal ordenada, chamada reta real. Se P for a representa\c c\~ao
de um n\'umero racional x, diremos que x \'e a abscissa de P. Note que nem todo ponto da reta real \'e racional. Para isso, considere
um quadrado de lado 1 e diagonal d. Pelo Teorema de Pit\'agoras, $d^{2} = 1^2 + 1^2 = 2.$ Agora, seja P a intersec\c c\~ao do eixo
x com a circunfer\^encia de centro em 0 e raio d. Mostremos que P \'e um ponto da reta com abscissa $x\not\in \mathbb{Q}.$
\begin{prop*}
  Seja a um inteiro. Ent\~ao, se a for \'impar, seu quadrado tamb\'em ser\'a. Al\'em disso, se a for par, seu quadrado tamb\'em \'e par.
\end{prop*}
\begin{prop*}
  A equa\c c\~ao $x^2 = 2$ n\~ao admite solu\c c\~ao racional.
\end{prop*}
  A ideia da prova \'e escrever um x na forma de fra\c c\~ao e chegar na contradi\c c\~ao de que tanto o numerador quanto o denominador
ser\~ao n\'umeros pares. Com isso, conclui-se que n\~ao existe racional irredut\'ivel com quadrado igual a 2, portanto n\~ao existe racional
satisfazendo a equa\c c\~ao. 

  Essa discuss\~ao mostra que existem v\~aos na ``reta'' dos racionais, requerindo a ado\c c\~ao de um novo corpo. Essa \'e a principal
motiva\c c\~ao por tr\'as dos n\'umeros reais, "preencher" os buracos deixados pelos racionais.
\begin{prop*}
  (Exerc\'icio.) Sejam $p_{1}, \ldots, p_{n}$ n\'umeros primos distintos. Ent\~ao, a equa\c c\~ao $x^{2} = p_{1}p_2\cdots p_{n}$ n\~ao
tem solu\c c\~ao racional.
\end{prop*}
  Vimos que os n\'umeros racionais com a sua adi\c c\~ao, multiplica\c c\~ao e rela\c c\~ao de ordem \'e um corpo ordenado. Nos interessamos,
tamb\'em, pelo corpo dos reais e dos racionais ($\mathbb{R}, \mathbb{C}$). De forma abstrata, um corpo \'e um conjunto n\~ao-vazio
 $\mathbb{F}$ em que est\~ao definidas duas opera\c c\~oes bin\'arias
 $$
    +:\mathbb{F}\times \mathbb{F}\rightarrow \mathbb{F}, \quad (x, y)\mapsto x + y
 $$
 e 
 $$
    \cdot: \mathbb{F}\times \mathbb{F}\rightarrow \mathbb{F}, \quad (x, y)\mapsto xy
 $$
 em que valem as oito propriedades vistas previamente para a defini\c c\~ao das opera\c c\~oes em $\mathbb{Q}$
 Se, ainda por cima, no corpo $\mathbb{F}$ est\'a definida uma ordem com propriedades an\'alogas \`as vistas para a qu\'adrupla
($\mathbb{Q}, +, \cdot, \leq{}$), diremos que ($\mathbb{F}, +, \cdot, \leq{}$) \'e um corpo ordenado.
\begin{def*}
  Diremos que um subconjunto A de um corpo $\mathbb{F}$ ordenado \'e limitado superiormente se existe um L neste corpo tal que $a \leq{L}$ para todo
a de A. 

  Definimos para um subconjunto limitado superiormente um n\'umero $\sup(A)\in \mathbb{F}$ como o menor limitante superior de
A, i.e., se $a \leq{\sup(A)}$ para todo a de A e se existe $f\in \mathbb{F}$ com $f < \sup(A),$ ent\~ao existe um a em A com $ f < a.$

  Por fim, diremos que um corpo ordenado \'e completo se todo subconjunto limitado superiormente possui supremo. $\square$
\end{def*}
  Nem todo subconjunto limitado superiormente de $\mathbb{Q}$ tem supremo, ou seja, $\mathbb{Q}$ n\~ao \'e completo.

\subsection{Os N\'umeros Reais ($\mathbb{R}$)}
  A ideia que iremos usar para construir o conjunto dos reais \'e que o conjunto dos n\'umeros reais preenche toda a reta real. Os Elementos
de $\mathbb{R}$ ser\~ao os subconjuntos de $\mathbb{Q}$ \`a esquerda de um ponto da reta real e ser\~ao chamados de cortes.
 \begin{def*}
   Um corte \'e um subconjunto $\alpha\subsetneq \mathbb{Q}$ com as seguintes propriedades:
  \begin{itemize}
    \item[i)] $\alpha\neq \emptyset$ e $\alpha \neq \mathbb{Q};$
    \item[ii)] Se $p\in \alpha$ e q \'e um racional com $q < p$, ent\~ao $q\in \alpha$ (todos os racionais \`a esquerda de um elemento
      de $\alpha$ est\~ao em $\alpha$);
    \item[iii)] Se $p\in \alpha$, existe um $r\in \alpha$ com $p < r$ ($\alpha$ n\~ao tem um maior elemento). $\square$
  \end{itemize}
 \end{def*}
  Essa ideia foi proposta inicialmente por Julius Wilhelm Richard Dedekind, um matem\'atico alem\~ao, em 1872, com o objetivo de encontrar uma explica\c c\~ao
e constru\c c\~ao elementar para os n\'umeros reais.
\begin{example}
  Se q \'e um racional, definimos $q^{*} = \{r\in \mathbb{Q}: r < q\}$.Ent\~ao, $q^{*}$ \'e um corte que chamamos de racional. Os 
cortes que n\~ao s\~ao desse tipo se chamam cortes irracionais.
\end{example}
\begin{example}
  $\sqrt{2} = \{q\in \mathbb{Q}: q^{2} < 2\}\cup \{q\in \mathbb{Q}: q < 0\}$ \'e um corte irracional.
\end{example}
  Observe que se $\alpha$ \'e um corte, p \'e um ponto dele e q n\~ao \'e, ent\~ao $p < q$. Al\'em disso, se r pertence a $\alpha$
e $r < s,$ ent\~ao s n\~ao pertence ao corte.
\begin{def*}
  Diremos que $\alpha < \beta$, em que $\alpha$ e $\beta$ s\~ao cortes, se $\alpha\subsetneq \beta.\square$
\end{def*}
\begin{prop*}
  Se $\alpha, \beta, \gamma$ s\~ao cortes,
 \begin{itemize}
   \item[i)] $\alpha < \beta$ e $\beta < \gamma$ implica que $\alpha < \gamma$;
   \item[ii)] Exatamente uma das seguintes rela\c c\~oes \'e v\'alida: $\alpha < \beta$ ou $\alpha = \beta$ ou $\beta < \alpha$
   \item[iii)] Todo subconjunto n\~ao-vazio e limitado superiormente de $\mathbb{R}$ tem supremo.
 \end{itemize}
\end{prop*}
\newpage

\section{Aula 03 - 17/03/2023}
\subsection{Motiva\c c\~oes}
\begin{itemize}
  \item Finalizar a constru\c c\~ao de $\mathbb{R}$ por cortes;
  \item Definir um corpo ordenado com base nos cortes;
\end{itemize}
\subsection{Cortes - Soma e Ordem}
  Coloquemos, para fins de conveni\^encia, $\mathbb{R}$ como a uni\~ao de todos os cortes.

  Vamos mostrar que os cortes racionais s\~ao, de fato, cortes. Considere, dado um racional q, $q^{*} = \{p\in \mathbb{Q}: p < q\}.$
Ele n\~ao pode completar todos os racionais, pois q + 1 n\~ao pertence a $q^{*}$. Al\'em disso, ele \'e n\~ao vazio, visto que
q-1 pertence a ele, mostrando a primeira propriedade dos cortes. 

  Ademais, se r pertence a $q^{*}$ e p \'e um racional menor que r, segue da transitividade da ordem que p \'e menor que
q, j\'a que r tamb\'em \'e. Assim, p pertence a $q^{*}$, mostrando a segunda propriedade dos cortes. 

  Por fim, dado um r em $q^{*},$ seja $s = \displaystyle \frac{r + q}{2}$. Ent\~ao, 
  $$
    r - \frac{r+q}{2} = \frac{r - q}{2} < 0,
  $$
tal que s \'e menor que r e, logo, pertence a $q^{*}$. Portanto, $q^{*}$ forma um corte.

  Daremos continuidade \`as atividades da aula anterior demonstrando a \'ultima proposi\c c\~ao vista.
\begin{prop*}
  Se $\alpha, \beta, \gamma$ s\~ao cortes,
 \begin{itemize}
   \item[i)] $\alpha < \beta$ e $\beta < \gamma$ implica que $\alpha < \gamma$;
   \item[ii)] Exatamente uma das seguintes rela\c c\~oes \'e v\'alida: $\alpha < \beta$ ou $\alpha = \beta$ ou $\beta < \alpha$
   \item[iii)] Todo subconjunto n\~ao-vazio e limitado superiormente de $\mathbb{R}$ tem supremo.
 \end{itemize}
\end{prop*}
  
 \begin{proof*}
   As duas primeiras partes seguem automaticamente da forma que definimos a ordem $\leq{}$ para os cortes. Resta mostrar a \'ultima.

   Vamos exibir o supremo explicitamente. Com efeito, seja $\mathcal{A}\subseteq{\mathbb{R}}$ um cole\c c\~ao de cortes
 limitada superiormente, i.e., existe um l em $\mathbb{R}$ tal que $\alpha \leq{l}$ para todo $\alpha$ em $\mathcal{A}.$
 Defina $\mathcal{S} = \bigcup_{\alpha\in \mathcal{A}}\alpha$. Mostremos que $\mathcal{S}$ \'e um corte. Com efeito,
 que $\mathcal{S}$ \'e n\~ao-vazio e diferente de $\mathbb{Q}$ \'e autom\'atico. Al\'em disso, dado q em $\mathcal{S}$ e $r < q,$
segue que $r\in \alpha_{0}$ para algum $\alpha_{0}$ em $\mathcal{A}.$

  Para ver que $\mathcal{S}$ \'e o supremo, suponha que $\mathcal{S}' < \mathcal{S}.$ Ent\~ao, existe r em $\mathcal{S}/\mathcal{S}'$.
Como r pertence a $\mathcal{S}$, r pertence a $\alpha_{0}$ para algum $\alpha_{0}\in \mathcal{A}.$ Logo, $\alpha_{0} > \mathcal{S}'.$ Portanto,
 $\mathcal{S}$ \'e o menor limitante superior de $\mathcal{A},$ ou seja, seu supremo. \qedsymbol
 \end{proof*}
 \begin{def*}
   Se $\alpha, \beta$ s\~ao cortes, definimos $\alpha + \beta$ como o conjunto de todos os racionais da forma $r + s$, com r em $\alpha$
 e s em $\beta$. Ademais, tome $0^* = \{s\in \mathbb{Q}: s < 0\}.\square$
 \end{def*}
  Vamos conferir a defini\c c\~ao, i.e., que $\alpha + \beta$ \'e um corte. Com efeito, $\alpha + \beta\neq\emptyset$, pois $\alpha\neq\emptyset$
e $\beta\neq\emptyset$. Al\'em disso, se p n\~ao pertence a $\alpha$ e q n\~ao pertence a $\beta$, mas r pertence a $\alpha$ e s a $\beta$,
ent\~ao $r + s < p + q$, tal que $p + q$ n\~ao pertence a $\alpha + \beta.$

  Al\'em disso, tome $r + s$ em $\alpha + \beta$ e $p < r + s$. Escreva $p = r' + s' = \underbrace{p - r}_{\in \beta} + \underbrace{r}_{\in \alpha}.$ Assim,
p pertence a $\alpha + \beta.$

  Por fim, tome $r + s$ em $\alpha + \beta$ e seja $r' > r$ (ambos em $\alpha$). Logo, $\underbrace{r' + s}_{\in \alpha + \beta} > r + s$. Portanto, 
$\alpha + \beta$ \'e um corte. 

  Fica de exerc\'icio mostrar que $0^*$ \'e um corte. Agora, mostremos os axiomas de corpo.

  A comutatividade e associatividade da adi\c c\~ao s\~ao triviais. Al\'em disso, dado r em $\alpha$ e s em $0^*,$
  $$
    r + s < r + 0 = r\Rightarrow r + s \in \alpha.
  $$
  Logo, $\alpha + 0^* \subseteq{\alpha}$. Por outro lado, dado r em $\alpha,$ existe $r'\text{ em } \alpha$ tal que $r' > r.$
  Assim, $r = \underbrace{r'}_{\alpha} + \underbrace{(r - r')}_{\in 0^*}$, pois $r - r' < 0.$ Portanto, $\alpha \subseteq{\alpha + 0^*}$ e
 $\alpha = \alpha + 0^*.$
 \begin{prop*}
   Dado um corte $\alpha$, existe um \'unico corte $\beta$ tal que $\alpha + \beta = 0^*$, em que 
   $$
    \beta = \{-p\in \mathbb{Q}: p - r\not\in \alpha \text{ para algum } r\in \mathbb{Q}, r > 0\}
   $$
   e \'e denotado por $-\alpha$. 
 \end{prop*}
\begin{proof*}
  Come\c camos mostrando que $\beta$ \'e um corte. Feito isso, vamos mostrar que $\beta + \alpha = 0^*.$

  Com efeito, dado -p em $\beta,$ segue que p n\~ao pertence a $\beta$. Caso $s = p + r$, -s pertence a $\beta$, tal que 
 $\beta$ \'e n\~ao-vazio. Ademais, se $p\in \alpha, -p\not\in \beta$, tal que $\beta$ \'e diferente de $\mathbb{Q}.$

  Al\'em disso, se $-q < -p$ e $-p\in \beta$, ent\~ao $-q \in \beta$. Por fim, se -p pertencer a $\beta$, $\displaystyle -p + \frac{r}{2}\in \beta$.
Portanto, $\beta$ \'e um corte. 

  Agora, vamos conferir o outro item. De fato, se r pertence a $\alpha$ e s a $-\alpha,$ ent\~ao $-s\not\in \alpha$ e $r < -s$, i.e.,
 $r + s < 0.$ Segue que $\alpha + (-\alpha) \subseteq{0^*}.$ Por outro lado, se $-2r\in 0^*$ com $r > 0,$ existe um inteiro n tal que
  $nr\in \alpha$ e $(n+1)r\not\in \alpha$. Escolha $p = -(n+2)r\in -\alpha$ e escreva $-2r = nr + p.$ Portanto, $0^*\subseteq{\alpha + (-\alpha)}$ e os
  conjuntos s\~ao iguais. \qedsymbol
\end{proof*}

\section{Aula 04 - 20/02/2023}
\subsection{Motiva\c c\~oes}
\begin{itemize}
  \item Definir multiplica\c c\~ao de cortes;
  \item Definir conceito de dist\^ancia entre n\'umeros de $\mathbb{R}$ 
\end{itemize}
\subsection{Cortes - Multiplica\c c\~ao}
\begin{def*}
  Se $\alpha, \beta$ s\~ao cortes,
  $$
    \alpha\beta = \left\{\begin{array}{ll}
      \alpha0^*, \quad\forall \alpha\in \mathbb{R}\\
      \{p\in\mathbb{Q}: \exists0 < r\in\alpha \text{ e }0 < s\in\alpha: p \leq rs\},\quad \alpha, \beta > 0^*\\
      (-\alpha)(-\beta), \quad \alpha, \beta < 0^*\\
      -[(-\alpha)\beta], \quad \alpha < 0^* e \beta > 0^*\\
      -[\alpha(-\beta)], \quad \alpha > 0^* e \beta < 0^*
      \end{array}\right.
  $$
  Definimos, tamb\'em, $1^*\{s\in\mathbb{Q}: s < 1\}$.
\end{def*}
\newpage

\subsection{$\mathbb{R}$ Como Corpo Ordenado Completo}
  Temos $\mathbb{Q}\subseteq{\mathbb{R}}$ e diremos que todo n\'umero que n\~ao \'e real \'e irracional. 
 \begin{theorem*}
   A qu\'adrupla $(\mathbb{R}, +, \cdot, \leq)$ satisfaz as condi\c c\~oes de corpo ordenado, de corpo e \'e completo.
 \end{theorem*}
\begin{def*}
  Seja $x\in \mathbb{R}.$ O m\'odulo, ou valor absoluto de x, \'e dado por
  $$
    |x| = \left\{\begin{array}{ll}
        x, \quad x \geq 0\\
        -x, x < 0
      \end{array}\right.
  $$
  Disto segue que $|x|\geq 0$ e $-|x|\leq x\leq |x|$ para todo x real.
\end{def*}
\begin{example}
  Mostre que $|x|^2 = x^2$, ou seja, o quadrado de um n\'umero real n\~ao muda quando se troca seu sinal.
\end{example}
\begin{example}
  A equa\c c\~ao $|x| = r$, com r maior que 0, tem como solu\c c\~oes apenas r e -r.
\end{example}
  Sejam P e Q dois pontos da reta real de abscissas x e y. Ent\~ao, a dist\^ancia de P a Q \'e definida por $|x-y|$. Assim,
 $|x-y|$ \'e a medida do segmento PQ. Em particular, como $|x|=|x-0|, |x|$ \'e a dist\^ancia de x a 0.
\begin{example}
  Seja r maior que 0. Ent\~ao, $|x| < r$ se, e somente se, $-r < x < r.$ Logo, o intervalo (-r, r) \'e o conjunto dos pontos reais
cuja dis\^ancia de 0 \'e menor que r.
\end{example}
\begin{example}
  Para quaisquer x, y reais, vale
  $$
    |xy| = |x||y|.
  $$
\end{example}
\begin{example}
  Para quaisquer x, y reais, temos
  $$
    |x+y| \leq |x| + |y|.
  $$
  Com efeito, somando $-|x|\leq{x}\leq{|x|}$ e $-|y|\leq{y}\leq{|y|},$ obtemos $-|x|-|y|\leq{x + y}\leq{|x|+|y|.}$\qedsymbol
\end{example}
\begin{def*}
  Um intervalo em $\mathbb{R}$ \'e um subconjunto de $\mathbb{R}$ que tem uma das seguintes formas:
 \begin{align*}
   &[a, b] = \{x\in \mathbb{R}: a\leq{x}\leq{b}\},\quad \text{ (Intervalo fechado.) }\\
   &(a, b) = \{x\in \mathbb{R}: a < x < b\},\quad \text{ (Intervalo aberto.) }\\
   &[a, b) = \{x\in \mathbb{R}: a \leq{x} < b\}\\
   &(a, b] = \{x\in \mathbb{R}: a < x \leq{b}\}\\
   &(-\infty, b] = \{x\in \mathbb{R}: x\leq{b}\}\\
   &(-\infty, b) = \{x\in \mathbb{R}: x > b\}\\
   &[a, +\infty) = \{x\in \mathbb{R}: x\geq{a}\}\\
   &(a, +\infty) = \{x\in \mathbb{R}: a < x\}\\
   &(-\infty, +\infty) = \mathbb{R}.
 \end{align*}
\end{def*}
\begin{def*}
  Um conjunto A de $\mathbb{R}$ \'e dito limitado se existir L positivo tal que $|x| \leq L$ para todo x em A.
\end{def*}
\begin{prop*}
  Um conjunto A de $\mathbb{R}$ \'e limitado se, e s\'o se, existir L positivo, tal que A est\'a contido em $[-L, L]$
\end{prop*}
\begin{example}
 \begin{itemize}
   \item[a)] $A = [0, 1]$ \'e limitado;
   \item[b)] $\mathbb{N}$ n\~ao \'e limitado;
   \item[c)] $B = \biggl\{\displaystyle \frac{2^n-1}{2^n}: n\in \mathbb{N}\biggr\}$ \'e limitado; 
   \item[d)] $C = \biggl\{\displaystyle \frac{2^n-1}{2^n}: n\in \mathbb{N}^{*}\biggr\}$ \'e limitado. 
 \end{itemize}
\end{example}
\begin{def*}
  Seja $A\subseteq{\mathbb{R}}$.
 \begin{itemize}
   \item A ser\'a dito limitado superiormente se existir um L real tal que $x\leq L$ para todo x de A. Diremos que L \'e o limitante superior de A.;
   \item A ser\'a dito limitado inferiormente se existir um L real tal que $x\geq L$ para todo x de A. Diremos que L \'e o limitante inferior de A.;
 \end{itemize}
 Caso ambos ocorram, diremos que A \'e limitado.
\end{def*}
\begin{def*}
  Seja A um subconjunto dos reais limitado superiormente e n\~ao-vazio. Diremos que $\overline{L}$ \'e o supremo de A se for um limitante superior
  e para qualquer outro limitante superior L de A, tivermos $\overline{L}\leq L$. Quando o supremo pertencer ao conjunto, chamaremos ele de m\'aximo.
\end{def*}
  Vimos que todo subconjunto n\~ao-vazio e limitado superiormente de $\mathbb{R}$ tem supremo.
\begin{def*}
  Seja A um subconjunto dos reais limitado inferiormente e n\~ao-vazio. Diremos que $\overline{l}$ \'e o \'infimo de A se for um limitante inferior
  e para qualquer outro limitante inferior l de A, tivermos $\overline{l}\geq l$. Quando o \'infimo pertencer ao conjunto, chamaremos ele de m\'inimo.
\end{def*}
\begin{prop*}
  Dado um subconjunto A dos reais n\~ao-vazio e limitado superiormente, $L = \sup{A}$ se, e somente se, 
 \begin{itemize}
   \item[a)] L for limitante superior de A;
   \item[b)] para todo $\epsilon > 0$, existe $a\in A$ tal que $a > L - \epsilon.$ 
 \end{itemize}
\end{prop*}
\begin{theorem*}
  O conjunto $A=\{nx: n\in\mathbb{N}\}$ ser\'a ilimitado para todo x n\~ao-nulo.
\end{theorem*}
\begin{proof*}
  Se $x > 0$, suponhamos, por absurdo, que A seja limitado e seja L seu supremo. Como $x > 0$, deve existir um natural m tal que
  $$
    L - x < mx \quad\text{ e } L = \sup{A} < (m+1)x.
  $$
  Mas isso \'e uma contradi\c c\~ao. 

  A prova para $x < 0$ \'e an\'aloga e ser\'a deixada como exerc\'icio. \qedsymbol 
\end{proof*}
 \begin{example}
  \begin{itemize}
    \item[a)] Considere $A = [0, 1).$ Ent\~ao, -2 e 0 s\~ao limitantes inferiores de A enquanto $1, \pi, 101$ s\~ao limitantes
  superiores de A.
    \item[b)] $\mathbb{N}$ n\~ao \'e limitado, mas \'e limitado inferiormente por 0, visto que $0\leq{x}$ para todo x natural.
    \item[c)] $B=\{x\in \mathbb{Q}: x\leq{\sqrt{2}}\}$ n\~ao \'e limitado, mas \'e limitado superiormente por L, em que $L\geq{2}.$ \qedsymbol
  \end{itemize}
 \end{example}
\begin{crl*}
  Para todo $\epsilon > 0$, existe um n natural tal que 
  $$
  \frac{1}{n} < \epsilon, \quad \frac{1}{n\sqrt{2}}<\epsilon, \quad 2^{-n} < \epsilon.
  $$
\end{crl*}
  J\'a sabemos, por constru\c c\~ao, que entre dois n\'umeros reais distintos existe um n\'umero racional. O mesmo vale para irracionais.
De fato, sejam a e b n\'umeros reais distintos. Se $a < b$ e $\epsilon = b - a > 0$, do corol\'ario, tome um natural n tal que
$\displaystyle \frac{1}{n\sqrt{2}} < \frac{1}{n} < \epsilon.$ Se a \'e racional, $r = \displaystyle a + \frac{1}{n\sqrt{2}}$ \'e irracional e
 $a < r < b.$ Por outro lado, se a \'e irracional, $r =\displaystyle a + \frac{1}{n}$ tamb\'em \'e, tal que $a < r < b.$ Portanto,
 dados dois n\'umeros reais quaisquer, existe um n\'umero irracional.
\begin{crl*}
  Qualquer intervalo aberto e n\~ao-vazio cont\'em infinitos n\'umeros racionais e infinitos irracionais.
\end{crl*}
\begin{crl*}
  Se $A = \biggr\{\displaystyle \frac{1}{n}: n\in \mathbb{N}^*\biggl\}$, ent\~ao $\inf A = 0.$
\end{crl*}
\begin{example}
 \begin{align*}
   &(a) \text{ Seja }A = (0, 1]. \text{ Ent\~ao, } \inf{A} = 0, \max{A} = 1; \\
   &(b) \sqrt{2} = \{r\in\mathbb{Q}: r \leq 0\}\cup \{r\in\mathbb{Q}: r^2 < 2\}\text{ \'e um corte.} 
   &(c) C = \{x\in\mathbb{Q}: x^2 < 2\} \Rightarrow \sqrt{2}=\sup{C}\text{ e }\inf{C} = -\sqrt{2}.
 \end{align*}
 Vamos analisar mais cautelosamente o item b e prov\'a-lo. De fato, se $0 < r\in \mathbb{Q}$ e $r^2 < 2,$ existe n natural tal que 
 $[2r + 1]\frac{1}{n} < 2 - r^2$ e $(r + \frac{1}{n})^2 < 2.$ As outras propriedades de cortes s\~ao triviais. 

 Olhando tamb\'em para o item C, como todos seus elementos s\~ao racionais saitsfazendo $x^2 < 2, \sqrt{2}$\'e um limitante superior de C.
 Agora, se $0 < L < \sqrt{2}$, existe um racional $r\in(L, \sqrt{2})$ e $L^2 < r^2 < 2.$ Logo, r pertence a C e L n\~ao \'e limitante superior para C,
 provando o resultado.
\end{example}
\begin{prop*}
  Se A \'e um subconjunto n\~ao-vazio e limitado inferiormente, ent\~ao $-A = \{-x: x\in A\}$ ser\'a limitado superiormente e 
  $\inf{A} = -\sup{(-A)}$. Analogamente, se for limitado superiormente, o conjunto -A ser\'a limitado inferiormente, e $\sup{A}=-\inf{(-A)}$
\end{prop*}
\begin{proof*}
  Se A for limitado inferiormente, $\inf{(A)} \leq{x}$ para todo x de A e, dado $\epsilon > 0$, deve existir a em A tal que
  $a < \inf{(A)} + \epsilon$, ou, trocando o sinal, $-\inf{(A)} \geq{-x}$ para todo -x de -A e, dado $\epsilon > 0,$ deve existir
  $b = -a$ em -A tal que $-a > -\inf{(A)} - \epsilon.$

  Com isso, segue que -A ser\'a limitado superiormente, e $\sup{(-A)} = -\inf{(A)}.$ A outra prova fica como exerc\'icio. \qedsymbol
\end{proof*}
\begin{crl*}
  Todo conjunto A n\~ao-vazio e limitado inferiormente de $\mathbb{R}$ tem \'infimo.
\end{crl*}
\begin{crl*}
  Todo conjunto A n\~ao-vazio e limitado de $\mathbb{R}$ tem \'infimo e supremo.
\end{crl*}
\begin{def*}
  Uma vizinhan\c ca de um n\'umero real a \'e qualquer intervalo aberto da reta contendo a.
\end{def*}
\begin{example}
  Se $\delta > 0, V_{\delta}(a)\coloneqq(a - \delta, a + \delta)$ \'e uma vizinhan\c ca de a que ser\'a chamada de $\delta-$vizinhan\c ca de a. 
\end{example}
\begin{def*}
  Sejam A um subconjunto de $\mathbb{R}$ e b um n\'umero real. Se, para todo $\delta > 0$, existir $a\in V_{\delta}(b)\cap{A}, a\neq b,$ 
ent\~ao b ser\'a dito um ponto de acumula\c c\~ao de A.
\end{def*}
\begin{example}
 \begin{itemize}
   \item[a)]O conjunto dos pontos de acumula\c c\~ao de (a, b) \'e [a, b];
    \item[b)]Seja $B = \mathbb{Z}.$ Ent\~ao, B n\~ao tem pontos de acumula\c c\~ao; 
      \item[c)] Subconjuntos finitos de $\mathbb{R}$ n\~ao t\^em pontos de acumula\c c\~ao;
        \item[d)] O conjunto dos pontos de acumula\c c\~ao de $\mathbb{Q}$ \'e $\mathbb{R}$.
 \end{itemize}
\end{example}
\begin{def*}
  Seja $B\subseteq{\mathbb{R}}$. Um ponto b de B ser\'a dito um ponto isolado de B, se existir $\delta > 0$ tal que $V_{\delta}(b)$
n\~ao cont\'em pontos de B distintos de b. $\square$
\end{def*}
\begin{example}
  Seja $B=\{1, \frac{1}{2}, \frac{1}{3}, \cdots\}$. Ent\~ao, o conjunto dos pontos de acumula\c c\~ao de B \'e $\{0\}$ e o conjunto dos pontos
isolados de B \'e o pr\'oprio conjunto B.
\end{example}
  Observe que existem conjuntos infinitos sem pontos de acumula\c c\~ao, tal como $\mathbb{Z}.$ Por outro lado, todo conjunto infinito e limitado possui
pelo menos um ponto de acumula\c c\~ao.
\begin{theorem*}
  Se A \'e um subconjunto infinito e limitado de $\mathbb{R},$ ent\~ao A possui pelo menos um ponto de acumula\c c\~ao.
\end{theorem*}
\begin{proof*}
  Se $A\subseteq{[-L, L]}$ e $[a_{n}, b_{n}], n\in \mathbb{N}$ s\~ao escolhidos tais que $[a_{n+1}, b_{n+1}]\subseteq{[a_{n}, b_{n}]}, b_{0} = -a_{0} = L,
  b_{n} - a_{n} = \frac{2L}{2^{n}}, n\in \mathbb{N}^*$ e $[a_{n}, b_{n}]$ cont\'em infinitos elementos de A. Seja $a = \sup{\{a_{n}: n\in \mathbb{N}\}}$.

  Note que $[a_{n}, b_{n}]\subseteq{a_{j}, b_{j}}, j\leq{n}$ e $[a_{j}, b_{j}]\subseteq{[a_{n}, b_{n}]}, j>n.$ Em qualquer um dos casos, $a_{n}\leq{b_{j}}$
para todo $j\in \mathbb{N}$. Logo, $a \leq{b_{j}}, j\in \mathbb{N}.$ Segue que $a_{n}\leq{a=\sup \{s_{n}: n\in \mathbb{N}\}}\leq{b_{n}}$ para todo
$n\in \mathbb{N}$ e $a\in\displaystyle\bigcap_{n\geq{1}}[a_{n}, b_{n}]$. Dado $\delta > 0,$ escolha $n\in \mathbb{N}$ tal que $\frac{2L}{2^{n}}<\delta.$ Segue
que $a\in[a_{n}, b_{n}]\subseteq{(a-\delta, a+\delta) = V_{\delta}(a)}$ e a \'e ponto de acumula\c c\~ao de A. \qedsymbol
\end{proof*}
\newpage

\section{Aula 05 - 22/03/2023}
\subsection{Motiva\c c\~oes}
\begin{itemize}
  \item Sequ\^encias de N\'umeros Reais;
  \item Converg\^encia de Sequ\^encias;
\end{itemize}
\subsection{Sequ\^encias de N\'umeros Reais}
\begin{def*}
  Uma sequ\^encia \'e uma fun\c c\~ao definida no conjunto dos n\'umeros reais que, para cada n natural, associa um n\'umero real $a_{n}$.
  \begin{align*}
    \mathbb{N}&=\{0, 1, 2, \cdots\}\\
              &f:\mathbb{N}\rightarrow \mathbb{R}\\
              &n\mapsto a_{n}.
  \end{align*}
  Denotamos a fun\c c\~ao por $\{a_{n}\}\square$
\end{def*}
 \begin{example}
   Sendo $a_{n}=\frac{f1}{n+1}$ para todo n natural, temos a sequ\^encia $\{1, \frac{1}{2}, \frac{1}{3}, \cdots\}$. 
 \end{example}
\begin{example}
  Sendo $a_{n} = 6$ para todo n natural, temos a sequ\^encia constante 
  $$
    \{6, 6, 6,\cdots\}.
  $$
\end{example}
\begin{example}
  Coloque $a_{2n+1} = 7, a_{2n}=4$ para todo n natural. Temos 
  $$
    \{4, 7, 4, 7, \cdots\}
  $$
\end{example}
  Consideremos as sequ\^encias
  $$
    \alpha_{n} = n, \quad \beta_{n} = (-1)^{n},\quad \text{ e } \gamma_{n} = \frac{1}{n}.
  $$
  Como fun\c c\~oes, elas podem ter os gr\'aficos tra\c cados, mas n\~ao s\~ao muito significativos, visto que consistem em
colet\^aneas de pontos discretos. Ademais, note que a sequ\^encia $(\alpha_{n})$ ``diverge'' para infinito, a sequ\^encia
$(\beta_{n})$ ``oscila'' e a sequ\^encia $(\gamma_{n})$ ``converge para 0''. Precisamente, 
 \begin{def*}
   A sequ\^encia $\{a_{n}\}$ \'e dita convergente com limite l se, para todo $\epsilon > 0$, existe um natural
N dependendo de $\epsilon (N = N(\epsilon)\in \mathbb{N})$ tal que $n > N$ implica em $|a_{n} - l|< \epsilon.$
  Ou seja, a partir de um certo N, os $a_{n}$ est\~ao no intervalo $(l-\epsilon, l+\epsilon)$ e, como $\epsilon$
\'e arbitr\'ario, os $a_{n}$ se juntam em torno de l. Disto, segue que a condi\c c\~ao exigida equivale a
  $$
  l - \epsilon < a_{n} < l + \epsilon, \quad n\geq{N}.
  $$
Denotamos esse fen\^omeno por $\displaystyle\lim_{n\to\infty}a_{n} = l$, ou $a_{n}\rightarrow l$, ou $a_{n}\overbracket[0pt]{\longrightarrow}^{n\to \infty}a.\square$.
\end{def*}
\begin{example}
  $\circ{}\frac{1}{n}\rightarrow0, n\rightarrow\infty$. De fato, dado $\epsilon > 0$, da propriedade arquimediana, segue que 
existe um N natural tal que $N\epsilon > 1.$ Logo, para todo $n\geq{N},$ temos 
  $$
  0 - \epsilon < \frac{1}{n}\leq{\frac{1}{N}} < 0 + \epsilon.
  $$

  $\circ \frac{n}{n+1}\rightarrow 1, n\rightarrow\infty$. Com efeito, dado $\epsilon > 0,$ queremos encontrar N natural n\~ao-nulo tal que
se n \'e maior que N, temos 
  $$
    \biggl|\frac{n}{n+1} - 1\biggr| < \epsilon.
  $$
  No entanto, $|\frac{n}{n+1}-1| = \frac{1}{n+1}$ e, da propriedade Archimediana, existe N em $\mathbb{N}^{\times}$ tal que
  $(N+1)\epsilon > 1$. Logo, se $n\geq{N},$
  $$
    1 - \epsilon < \frac{n}{n+1} < 1 + \epsilon.
  $$
\end{example}
\begin{def*}
  Uma sequ\^encia $\{a_{n}\}$ ser\'a divergente quando ela n\~ao for convergente.
 \begin{itemize}
    \item[I)] Sequ\^encia divergente para $+\infty:$ Este caso ocorre se dado $K > 0$, existe N natural tal que se $n > N,
a_{n} > K.$
    \item[II)] Sequ\^encia divergente para $-\infty:$ Acontece quando dado $K > 0$, existe N natural tal que se $n > N, 
a_{n} < -K.$
    \item[III)]Sequ\^encia oscilante: Por fim, ocorre quando a sequ\^encia diverge, mas nem para $+\infty$ e nem para $-\infty.\square$
 \end{itemize}
\end{def*}
  Note que, como sequ\^encias s\~ao fun\c c\~oes, podemos multiplic\'a-las por constante, somar, dividir e multiplicar por outras sequ\^encia. De fato,
 \begin{def*}
   Dadas sequ\^encias $\{a_{n}\}, \{b_{n}\}$ e um n\'umero real c, deifnimos 
  \begin{align*}
    &i) \{a_{n}\} + \{b_{n}\} = \{a_{n} + b_{n}\}\\
    &ii) c\{a_{n}\} = \{c \cdot a_{n}\}\\
    &iii) \{a_{n}\}\{b_{n}\} = \{a_{n}b_{n}\}\\
    &iv) \text{ Se }b_{n}\neq0\forall n\in \mathbb{N}, \frac{\{a_{n}\}}{\{b_{n}\}} = \biggl\{\frac{a_{n}}{b_{n}}\biggr\}\square
  \end{align*}
 \end{def*}
\begin{def*}
  Seja $\{a_{n}\}$ uma sequ\^encia de n\'umero reais. Diremos que $\{a_{n}\}$ \'e limitada se sua imagem for um subconjunto
limitado de $\mathbb{R}.\square$
\end{def*}
\begin{theorem*}
  Seja $\{a_{n}\}$ uma sequ\^encia de n\'umeros reais.
 \begin{itemize}
   \item[a)] $a_{n}\overbracket[0pt]{\longrightarrow}^{n\rightarrow\infty}a$ se, e somente, toda vizinhan\c ca de a cont\'em todo, exceto uma poss\'ivel quantidade
  finita de $a_{n}$'s.
   \item[b)] O limite \'e \'unico.
   \item[c)] Se $\{a_{n}\}$ \'e convergente, ent\~ao $\{a_{n}\}$ \'e limitada
   \item[d)] Se $a_{n}\overbracket[0pt]{\longrightarrow}^{n\rightarrow\infty}a$, exite N natural tal que $a_{n} > 0$ para todo $n\geq{N}.$
   \item[e)] Se $A\subseteq{\mathbb{R}}$ e a \'e um ponto de acumula\c c\~ao de A, ent\~ao existe uma sequ\^encia $\{a_{n}\}$ de elementos
    de A que converge para a.
 \end{itemize}
\end{theorem*}
\begin{proof*}
  O item a \'e trivial. Mostremos a unicidade do limite: Suponha que $a_{n}$ converge para a e para b, com a diferente de b. Ent\~ao,
dado $\epsilon > 0$, existem naturais $N_{1}, N_{2}$ tais que se $n\geq{N_{1}}, |a_{n}-a|<\epsilon$ e se $n\geq{N_{2}}, |a_{n} - b| < \epsilon.$
Tome $N = \min{N_{1}, N_{2}}$ e suponha que $n \geq{N}.$ Ent\~ao, temos 
 $$
  |b - a| \leq{|b - a_{n}| + |a_{n} - a|} = |b - a_{n}| + |a - a_{n}| < 2\epsilon.
 $$
 (\textit{P.S.}: pode ser boa pr\'atica tomar $\frac{\epsilon}{2}$ ao inv\'es de $\epsilon$, pois assim obtemos $|b-a|<\frac{2\epsilon}{2}=\epsilon.$)

 Como $\epsilon$ \'e abritr\'ario, podemos selecionar $\epsilon$ infinitamente pr\'oximo de 0. Portanto, b = a.

 Para o item c, suponha que $a_{n}$ converge para a, isto \'e, dado $\epsilon > 0, \epsilon = 1$ em particular, existe 
 $N\in \mathbb{N}$ tal que se $n \geq{N}, |a_{n} - a| < 1$. Logo, $a_{n}\in(a - 1, a + 1)$ para n maior que N suficientemente grande.
 Restam os N-1 primeiros elementos da sequ\^encia. Assim, tome $R = \max{\biggl\{|a_{1}|, \cdots, |a_{N-1}|, |a + 1|, |a - 1|\biggr\}}$. Deste modo,
 $a_{n}\in[-R, R]$ para todo n natural. 

 Com rela\c c\~ao ao item d, basta tomar $\epsilon = \frac{a}{2} > 0.$ 

 Por fim, quanto ao item e, suponha o que \'e dito no enunciado. Como a \'e ponto de acumula\c c\~ao, dado $\epsilon > 0,$ existe
 $a'\in{A}, a'\neq a$ tal que 
 $$
  a'\in V_{\epsilon}(a) = (a - \epsilon, a + \epsilon).
 $$
 Logo, tomadno $\epsilon = \frac{1}{n},$ podemos encontrar $a_{n}\in A, a_{n}\neq a$ tal que $a_{n}\in\biggl(a-\frac{1}{n}, a + \frac{1}{n}\biggr)$. A sequ\^encia
 $\{a_{n}\}$ converge para a. De fato, dado $\epsilon > 0$, tome N natural tal que $N\epsilon > 1.$ Assim, se $n\geq{N}, a_{n}\in(a-\frac{1}{n}, a+\frac{1}{n})\subseteq{a-\epsilon}, a+\epsilon)$.
 Portanto, $a_{n}\rightarrow a.$ \qedsymbol
\end{proof*}
 \begin{theorem*}
   Seja $a_{n}\overbracket[0pt]{\longrightarrow}^{n\to\infty}a, b_{n}\overbracket[0pt]{\rightarrow}^{n\to\infty}b$ e c um n\'umero real. Ent\~ao,
  \begin{align*}
    &a) a_{n} + b_{n}\overbracket[0pt]{\longrightarrow}^{n\to\infty} a + b.\\
    &b) ca_{n}\overbracket[0pt]{\longrightarrow}^{n\to\infty} ca\\
    &c) a_{n}b_{n}\overbracket[0pt]{\longrightarrow}^{n\to\infty} ab\\
    &d)\text{Se} b\neq0, b_{n}\neq0\forall n\in \mathbb{N}, \frac{a_{n}}{b_{n}}\overbracket[0pt]{\longrightarrow}^{n\to\infty}\frac{a}{b}.
  \end{align*}
 \end{theorem*}
\begin{proof*}
  Item c). Suponha $a_{n}\overbracket[0pt]{\longrightarrow}^{n\to \infty}a, b_{n}\overbracket[0pt]{\longrightarrow}^{n\to \infty}b$. Note que 
  $$
    |a_{n}b_{n} - ab| = a_{n}b_{n} - a_{n}b + a_{n}b - ab \leq{|a_{n}||b_{n}-b| + |b||a_{n}-a|}
  $$
  Como $\{a_{n}\}$ \'e convergente, ela \'e limitada pelo teorema anterior. Assim, existe $M > 0$ tal que $|a_{n}|\leq{M}$ para todo n natural, tal que
  Assim, 
  $$
    |a_{n}b_{n} - ab| \leq{|a_{n}||b_{n} - b| + |b||a_{n} - a|} \leq{M|b_{n} - b| + (|b| + 1)|a_{n} - a|}.
  $$
  Agora, dado $\epsilon > 0,$ existem naturais $N_{1}, N_{2}$ tais que 
 \begin{align*}
   &|a_{n}-a| < \frac{\epsilon}{2(|b|+1)},\quad\forall n\geq{N_{1}}\\
   &|b_{n}-b| < \frac{\epsilon}{2M},\quad\forall n\geq{N_{2}}.
 \end{align*}
 Logo, tomando $N = \max\{N_{1}, N_{2}\},$ se $n\geq{N},$
 $$
  |a_{n}b_{n}-ab| < \frac{\epsilon}{2} + \frac{\epsilon}{2} = \epsilon.
 $$
 Portanto, $a_{n}b_{n}\overbracket[0pt]{\longrightarrow}^{n\to \infty}ab.$ \qedsymbol
\end{proof*}
\begin{def*}
   Seja $\{a_{n}\}$ uma sequ\^encia. Diremos que $\{b_{n}\}$ \'e uma subsequ\^encia de $\{a_{n}\}$ se existir uma fun\c c\~ao
estritamente crescente $s:\mathbb{N}\rightarrow \mathbb{N}$ tal que $b_{k} = a_{s(k)}$ para todo k natural. $\square$
\end{def*}
\begin{def*}
  Seja $\{a_{n}\}$ uma sequ\^encia. Diremos que $\{a_{n}\}$ \'e de Cauchy se, dado $\epsilon > 0$, existe um natural
  $N = N(\epsilon)$ tal que $|a_{n}-a_{m}| < \epsilon$ para todo $n, m\geq{N}.\square$
\end{def*}
\begin{theorem*}
 \begin{itemize}
   \item[a)]Uma sequ\^encia \'e convergente se, e somente se, toda subsequ\^encia dela converge para o mesmo limite.
   \item[b)] Toda sequ\^encia convergente \'e de Cauchy;
   \item[c)] Toda sequ\^encia limitada tem subsequ\^encia convergente;
   \item[d)] Toda sequ\^encia de Cauchy \'e limitada;
   \item[e)] Toda sequ\^encia de Cauchy que tem subsequ\^encia convergente \'e convergente.
   \item[f)] Toda sequ\^encia de Cauchy \'e convergente;
   \item[g)] Toda sequ\^encia crescente e limitada \'e convergente;
   \item[h)] Toda sequ\^encia decrescente e limitada \'e convergente.
 \end{itemize}  
\end{theorem*}
\newpage

\section{Aula 06 - 24/03/2023}
\subsection{Motiva\c c\~oes}
 \begin{itemize}
   \item Provar o teorema da aula anterior;
   \item Exemplos.
 \end{itemize}
\subsection{Propriedades de Sequ\^encias}
  Recapitulemos o teorema da aula anterior:
\begin{theorem*}
 \begin{itemize}
   \item[a)]Uma sequ\^encia \'e convergente se, e somente se, toda subsequ\^encia dela converge para o mesmo limite.
   \item[b)] Toda sequ\^encia convergente \'e de Cauchy;
   \item[c)] Toda sequ\^encia limitada tem subsequ\^encia convergente;
   \item[d)] Toda sequ\^encia de Cauchy \'e limitada;
   \item[e)] Toda sequ\^encia de Cauchy que tem subsequ\^encia convergente \'e convergente.
   \item[f)] Toda sequ\^encia de Cauchy \'e convergente;
   \item[g)] Toda sequ\^encia crescente e limitada \'e convergente;
   \item[h)] Toda sequ\^encia decrescente e limitada \'e convergente.
   \end{itemize}
 \end{theorem*}
 \begin{proof*}
a.) $\Leftarrow)$ Se toda subsequ\^encia de $\{a_{n}\}$ converge, ent\~ao $\{a_{n}\}$ converge, pois ela \'e uma subsequ\^encia de si mesma (basta tomar $s:\mathbb{N}\rightarrow \mathbb{N}, s(n) = n.$)'

$\Rightarrow)$ Suponha que $a_{n}\overbracket[0pt]{\longrightarrow}^{n\to \infty}l$ e $\{b_{n}\}$ \'e uma subsequ\^encia de $\{a_{n}\}$, existe 
$s:\mathbb{N}\rightarrow \mathbb{N}$ estritamente crescente tal que $b_{k} = a_{s(k)}.$ Dado $\epsilon > 0$, seja N o natural tal que
 $|a_{n} - l| < \epsilon$ para todo $n\geq{N}.$ Note que $s(n)\geq{n},$ tal que se $n\geq{N},$ ent\~ao $s(n)\geq{N},$ de forma que 
 $|a_{s(n)} - l| < \epsilon$. Portanto, qualquer subsequ\^encia de $\{a_{n}\}$ \'e convergente.

 b.) Se $a_{n}\overbracket[0pt]{\longrightarrow}^{n\to \infty}l,$ ent\~ao dado $\epsilon > 0$, existe N natural tal que 
 $$
  |a_{n}-l|<\frac{\epsilon}{2},\quad\forall n\geq{N}.
 $$
 Logo, $|a_{n}-a_{m}| = |a_{n}-l + l-a_{m}| \leq{|a_{n}-l| + |l-a_{m}|} < \frac{\epsilon}{2} + \frac{\epsilon}{2} = \epsilon$ para todo $n, m \geq{N}.$

 c.) Suponha que $\{a_{n}\}$ \'e uma sequ\^encia limitada. Recorde que, do teorema de Bolzano-Weierstrass, todo conjunto inifinito e limitado
possui um ponto de acumula\c c\~ao. Segue que a imagem I da sequ\^encia \'e finita ou infinita.

No primeiro caso, se I \'e finito, um dos valores pertencentes a I \'e tal que $a_{n} = a$ para infinitos \'indices. 
Construiremos a sequ\^encia como segue - Coloque s(0) como o menor elemento do conjunto dos n's para os quais 
$a_{n} = a, i.e.,\{n\in \mathbb{N}: a_{n} = a\} = A. $ Al\'em disso, tome s(1) como  o menor elemento de A, com  exce\c c\~ao do
s(0). Repetindo esse processo, obtemos uma subsequ\^encia constante at\'e que se obtenha s(n) = a, ou seja, ela ser\'a convergente.

Agora, se I \'e infinito, segue de Bolzano-Weierstrass que I tem um ponto de acumula\c c\~ao, nomeie-o de a. Dado $\epsilon > 0, (a-\epsilon, a+\epsilon)$ 
tem infinitos elementos do conjunto I. Analogamente ao anterior, coloque N = s(0) como o menor elemento de $\{n\in \mathbb{N}: a_{n}\neq a, a_{n}\in(a-\epsilon, a+\epsilon)\}$ e coloque, tamb\'em,
$\epsilon_{1} = |a-a_{s(0)}|$. Em seguida, tome $s(1) = \{n\in \mathbb{N}: a_{n}\neq a, a_{n}\in(a-\frac{\epsilon}{2}, a+\frac{\epsilon}{2})\}$. Indutivamente,
$b = a_{s(n)}$ \'e convergente para a.

d.) Dado $\epsilon = 1,$ seja N um n\'umero natural tal que 
  $$
    |a_{n}-a_{m}| < 1,\quad\forall n\geq{N}.
  $$
  Considere $M = \{|a_{0}|, |a_{1}|, \cdots, |a_{N-1}|, |a_{N}+1|, |a_{N}-1|\}.$ Assim, $a_{n}\in[-M, M]$ para todo n natural.

e.) Seja $\{a_{n}\}$ de Cauchy e $\{a_{s(n)}\}$ convergente para l. Dado $\epsilon > 0$, existe um natural $N_{1}$ tal que 
  $$
    |a_{n}-a_{m}| < \frac{\epsilon}{2},\quad\forall n\geq{N_{1}}.
  $$
  Al\'em disso, existe $N_{2}$ natural tal que 
  $$
    |a_{s(n)} - l| < \frac{\epsilon}{2},\quad\forall s(n)\geq{N_{2}}.
  $$
  Seja $N=\max\{s(N_{2}), N_{1}\}$ e tome $n\geq{N}.$
  $$
    |a_{n}-l| = |a_{n} - a_{s(N_{2})} + a_{s(N_{2})} - l| \leq{|a_{n}-a_{s(N_{2})}| + |a_{s(N_{2})} - l|} < \frac{\epsilon}{2} + \frac{\epsilon}{2} = \epsilon.
  $$

f.) Segue os itens (e), (d) e (c), visto que toda subsequ\^encia de Cauchy ter\'a subsequ\^encia convergente pelos itens (d) e (c).

g.) Seja $\{a_{n}\}$ limitada e crescente, $l = \sup\{a_{n}:n\in \mathbb{N}\}.$ Ent\~ao, para todo $n \geq{N}$, em que N \'e tal que $a_{N}\in(l-\epsilon, l)$
 $$
  l-\epsilon < a_{N}\leq{a_{n}} \leq{l}.
 $$

h.) An\'aloga ao g.
\end{proof*}
 \begin{example}
   Mostre que 
 \begin{itemize}
   \item[i)] $\{a, a, a, \cdots\},a\in \mathbb{R}$ \'e convergente;
   \item[ii)] $\{0, 1, 0, 1\}$ n\~ao \'e convergente;
   \item[iii)] $\{n\}$ n\~ao \'e convergente.
 \end{itemize}
 \end{example}
\begin{example}
  Se a \'e um n\'umero real mais ou igual a zero, ent\~ao a sequ\^encia $\{a^{n}\}$ \'e convergente se $0\leq{a}\leq{1}$ e divergente
  se $a > 1$. Com efeito, se $a > 1, a = 1 + h, h > 0$. Ent\~ao, 
  $$
    a^{n} = (1+h)^{n} = \sum\limits_{k=0}^{n}\binom{n}{k}1^{n-k}h^{k} = 1 + nh + \cdots > 1 + nh.
  $$
  Mas, segue da Archimediana que $1 + nh$ sempre forma um conjunto ilimitado para n natural, ou seja, $a_{n}$ \'e ilimitada. Logo, a sequ\^encia
diverge.

  Por outro lado, suponha que a pertence a (0, 1). Ent\~ao, $a^{n+1} = a a^{n} < a^{n}$, ou seja, \'e uma sequ\^encia decrescente e limitada inferiormente.
Portanto $\{a_{n}\}$ \'e convergente.
\end{example}
\begin{example}
  Mostre que, se a \'e diferente de 1, 
  $$
    \sum\limits_{i=0}^{n}a^{i} = \frac{1-a^{n+1}}{1-a}
  $$
  e que a sequ\^encia $\biggl\{\frac{1-a^{n+1}}{1-a}\biggr\}$ \'e convergente se $0\leq{a}<1$ e divergente se $a > 1$.
\end{example}
\begin{example}
  Mostre que a sequ\^encia $\{a_{n}\}$, com $a_{n} = \displaystyle \sum\limits_{i=0}^{n}\frac{1}{i!}$ \'e convergente para todo n natural. (Crescente e limitada por 3.)
\end{example}
\begin{example}
  Mostre que as sequ\^encias $\biggl\{(1+\frac{1}{n}^{n})\biggr\}, \{n^{\frac{1}{n}}\}$ e $\{a^{\frac{1}{n}}\}$ com $a >0,$ s\~ao
convergentes.
 \begin{align*}
   &\circ (1+\frac{1}{n})^{n} = 1 + 1 + \frac{1}{2!}(1-\frac{1}{n}) + \cdots + \frac{1}{n!}(1-\frac{1}{n})(1-\frac{1}{n})(1-\frac{2}{n})\cdots(1-\frac{n-1}{n})\\
   &\circ n^{\frac{1}{n}} > (n+1)^{\frac{1}{n+1}}\Longleftrightarrow n^{n+1} > (n+1)^{n}\Longleftrightarrow n>(1+\frac{1}{n})^{n}\\
   &\circ x = a^{n} < 1\Rightarrow x < 1, x^{n} = a, x^{n+1} = a^{\frac{n+1}{n}},\text{ e } y^{n+1} = a \Rightarrow \biggl(\frac{x}{y}\biggr)^{n+1} = a^{\frac{1}{n}}.
 \end{align*}
\end{example}
\newpage

\section{Aula 07 - 27/03/2023}
\subsection{Motiva\c c\~oes}
 \begin{itemize}
   \item Exemplos de Sequ\^encias;
   \item Teorema da Compara\c c\~ao e do Sandu\'iche;
   \item Limites superior e inferior.
 \end{itemize}
\subsection{Exemplos de Sequ\^encias}
  Revisemos os exemplos da \'ultima aula, com um extra ao final.
  \begin{example}
   Mostre que 
 \begin{itemize}
   \item[i)] $\{a, a, a, \cdots\},a\in \mathbb{R}$ \'e convergente;
   \item[ii)] $\{0, 1, 0, 1\}$ n\~ao \'e convergente;
   \item[iii)] $\{n\}$ n\~ao \'e convergente.
 \end{itemize}
 \end{example}
\begin{example}
  Se a \'e um n\'umero real mais ou igual a zero, ent\~ao a sequ\^encia $\{a^{n}\}$ \'e convergente se $0\leq{a}\leq{1}$ e divergente
  se $a > 1$. Com efeito, se $a > 1, a = 1 + h, h > 0$. Ent\~ao, 
  $$
    a^{n} = (1+h)^{n} = \sum\limits_{k=0}^{n}\binom{n}{k}1^{n-k}h^{k} = 1 + nh + \cdots > 1 + nh.
  $$
  Mas, segue da Archimediana que $1 + nh$ sempre forma um conjunto ilimitado para n natural, ou seja, $a_{n}$ \'e ilimitada. Logo, a sequ\^encia
diverge.

  Por outro lado, suponha que a pertence a (0, 1). Ent\~ao, $a^{n+1} = a a^{n} < a^{n}$, ou seja, \'e uma sequ\^encia decrescente e limitada inferiormente.
Portanto $\{a_{n}\}$ \'e convergente.
\end{example}
\begin{example}
  Mostre que, se a \'e diferente de 1, 
  $$
    \sum\limits_{i=0}^{n}a^{i} = \frac{1-a^{n+1}}{1-a}
  $$
  e que a sequ\^encia $\biggl\{\frac{1-a^{n+1}}{1-a}\biggr\}$ \'e convergente se $0\leq{a}<1$ e divergente se $a > 1$.
\end{example}
\begin{example}
  Mostre que a sequ\^encia $\{a_{n}\}$, com $a_{n} = \displaystyle \sum\limits_{i=0}^{n}\frac{1}{i!}$ \'e convergente para todo n natural. (Crescente e limitada por 3.) 

  De fato, \'e claro que $\{a_{n}\} $ \'e crescente e que $\frac{1}{n!}\leq{\frac{1}{2^{n-1}}},$ para $n\geq{2}.$ Logo, 
  $$
  a_{n}\leq{1+\sum\limits_{k=0}^{n}\frac{1}{2^{k}}} = 
  1 + \frac{1-\frac{1}{2^{n+1}}}{1-\frac{1}{2}} < 3.
  $$
  Portanto, $\{a_{n}\}$ \'e convergente, e denotamos seu limite por e.
\end{example}
\begin{example}
  Mostre que as sequ\^encias $\biggl\{(1+\frac{1}{n}^{n})\biggr\}, \{n^{\frac{1}{n}}\}$ e $\{a^{\frac{1}{n}}\}$ com $a >0,$ s\~ao
convergentes.
 \begin{align*}
   &\circ (1+\frac{1}{n})^{n} = 1 + 1 + \frac{1}{2!}(1-\frac{1}{n}) + \cdots + \frac{1}{n!}(1-\frac{1}{n})(1-\frac{1}{n})(1-\frac{2}{n})\cdots(1-\frac{n-1}{n})\\
   &\circ n^{\frac{1}{n}} > (n+1)^{\frac{1}{n+1}}\Longleftrightarrow n^{n+1} > (n+1)^{n}\Longleftrightarrow n>(1+\frac{1}{n})^{n}\\
   &\circ x = a^{n} < 1\Rightarrow x < 1, x^{n} = a, x^{n+1} = a^{\frac{n+1}{n}},\text{ e } y^{n+1} = a \Rightarrow \biggl(\frac{x}{y}\biggr)^{n+1} = a^{\frac{1}{n}}.
 \end{align*}
   Ainda mais, uma delas t\^em como limite o n\'umero e definido no exemplo anterior. Para observar isso, considere o primeiro exemplo. Note que
  \begin{align*}
    b_{n} &= 1 + \binom{n}{1}n^{-1} + \binom{n}{2}n^{-2} + \cdots + \binom{n}{n-1}n^{-n+1} + \binom{n}{n}n^{-n} \\ 
          &= 1 + 1 + \frac{1}{2!}(1-\frac{1}{n}) + \cdots + \frac{1}{n!}(1-\frac{1}{n})(1-\frac{1}{n})(1-\frac{2}{n})\cdots(1-\frac{n-1}{n})\\
          &\leq{1 + 1 + \frac{1}{2!} + \cdots + \frac{1}{n!} = a_{n} < e}
 \end{align*}
 Como cada termo da soma que define $b_{n}$ \'e crescente, obtemos que $b_{n}$ \'e crescente, tal que ela converge com limite $l = \sup{\{b_{n}:n\in \mathbb{N}\}}$.
  
   Com rela\c c\~ao ao \'ultimo item, resta elaborar como ele converge para 1. Lembre-se que $a^{\frac{1}{n}}$ \'e o \'unico n\'umero real positivo
  x tal que $x^{n} = a.$ Logo, se $x = a^{n} $ e $y = a^{\frac{1}{n+1}}$, temos $x^{n+1} = y^{n+1}x$ e, deste modo,
  \begin{align*}
  & (a) 0 < a < 1 \Rightarrow x < 1 \text{ e } \biggl(\frac{x}{y}\biggr)^{n+1} = x < 1\text{ e, assim, } x < y.\\
  & (b) a > 1 \Rightarrow x > 1\text{ e }\biggl(\frac{x}{y}\biggr)^{n+1} = x > 1\text{ tal que } x > y.
  \end{align*}
  Logo, se $a < 1,\{a^{\frac{1}{n}}\} $ \'e crescente e limitada superiormente por 1, mostrando que ela \'e convergente. Al\'em disto,
se $a > 1,\{a^{\frac{1}{n}}\} $ \'e descrescente e limitada inferiormente por 1, tamb\'em sendo convergente. Por fim, segue de
$a^{\frac{1}{n(n+1)}} = \frac{a^{\frac{1}{n}}}{a^{\frac{1}{n+1}}}$. Portanto, do item (a) do teorema junto com a regra para quociente de
sequ\^encias, segue que l = 1 \'e o limite dela.
\end{example}\begin{example}
   Mostre que a sequ\^encia $\{c_{n}\}, c_{0} = 1, c_{n} = n^{\frac{1}{n}}, n\geq{1}$, \'e convergente. Com efeito, lembre-se que,
para $n\geq{3}, n > b_{n} = (1+\frac{1}{n})^{n}.$ Logo, para $n\geq{3}, n^{n+1}>(n+1)^{n}$ e, consequentemente, $n^{\frac{1}{n}} > (n+1)^{\frac{1}{n+1}}.$\
Disto segue de $\{n^{\frac{1}{n}}\} $ \'e limitada e, por (h), que $\{c_{n}\} $ \'e convergente com limite $l\geq{1}.$ Ainda mais,
 $(2n)^{\frac{1}{2n}}(2n)^{\frac{1}{2n}} = (2n)^{\frac{1}{n}} = 2^{\frac{1}{n}}n^{\frac{1}{n}}$ e, portanto, usamos (a) e o exemplo da \'ultima aula para
mostrar que $l^{2} = l = 1.$\qedsymbol
 \end{example}

\subsection{Teoremas da Compara\c c\~ao e do Sandu\'iche}
  Nota\c c\~ao: Se uma sequ\^encia tem limite 0, ela \'e chamada infinit\'esima.
\begin{theorem*}
  Se $\{a_{n}\}$ \'e limitada e $\{b_{n}\}$ \'e infinit\'esima, ent\~ao $\{a_{n}\cdot b_{n}\} $ \'e infinit\'esima.
\end{theorem*}
\begin{proof*}
  Como $\{a_{n}\}$ \'e limitada, seja $M > 0$ tal que $|a_{n}|\leq{M}$ para todo n natural. Como $\{b_{n}\} $ \'e infinit\'esima, dado $\epsilon > 0$,
seja N outro natural tal que $|b_{n}|<\frac{\epsilon}{M}$ para todo $n\geq{N}.$ Segue que
  $$
    |a_{n}b_{n}| \leq{M|b_{n}|} < M \frac{\epsilon}{M} = \epsilon,\quad \forall n\geq{N}.
  $$
  Portanto, $\{a_{n}b_{n}\}\overbracket[0pt]{\longrightarrow}^{n\to \infty}0$.\qedsymbol
\end{proof*}
 \begin{example}
   Mostre que $\biggl\{\frac{n+\cos{(n)}}{n+1}\biggr\}$ converge.
 \end{example}
 Os resultados a seguir s\~ao os dois mencionados previamente, o teorema da compara\c c\~ao e o do sandu\'iche, respectivamente.
\begin{theorem*}
  Se $a_{n}\overbracket[0pt]{\longrightarrow}^{n\to \infty}a, b_{n}\overbracket[0pt]{\longrightarrow}^{n\to \infty}b$ e existe N natural tal que, 
para todo $n\geq{N}, a_{n}\leq{b_{n}}$, ent\~ao $a\leq{b}.$
\end{theorem*}
 \begin{proof*}
   Dado $\epsilon > 0$, existe $N_{1}\leq{N}$ tal que, para todo $n\geq{N_{1}}$,
   $$
    a - \epsilon < a_{n} < a + \epsilon\quad\text{ e }\quad b-\epsilon < b_{n} < b + \epsilon.
   $$
   Logo, para todo $n\geq{N},$
   $$
    a-\epsilon < a_{n} \leq{b_{n}} < b + \epsilon.
   $$
   Desta forma, $a-b < \epsilon$ para todo $\epsilon > 0$ e, portanto, $a - b\leq{0}.$\qedsymbol
 \end{proof*}

\begin{theorem*}
  Se $a_{n}\overbracket[0pt]{\longrightarrow}^{n\to \infty}l, c_{n}\overbracket[0pt]{\longrightarrow}^{n\to \infty}l$ e existe um
N natural tal que, para todo $n\geq{N}, a_{n}\leq{b_{n}}\leq{c_{n}}$, ent\~ao $b_{n}\overbracket[0pt]{\longrightarrow}^{n\to \infty}l.$
\end{theorem*}
\begin{proof*}
  Dado $\epsilon>0$, existe $N_{1}\geq{N}$ tal que, para todo $n\geq{N_{1}},$
  $$
    l - \epsilon < a_{n} < l + \epsilon\quad\text{ e }\quad l-\epsilon < c_{n} < l+\epsilon.
  $$
  Logo, para todo $n\geq{N_{1}},$
  $$
    l - \epsilon < a_{n}\leq{b_{n}}\leq{c_{n}}<l + \epsilon.
  $$
  Disto segue que $|b_{n} - l|< \epsilon$ para todo $n\geq{N_{1}}$ e que, portanto, $\{b_{n}\} $ \'e convergente para l. \qedsymbol
\end{proof*}
\begin{example}
  Vamos mostrar que 
  $$
   e = \lim_{n\to\infty}\overbrace{(1 + \frac{1}{1!} + \frac{1}{2!} + \cdots + \frac{1}{n!})}^{a_{n}} = \lim_{n\to\infty}\overbrace{\biggl(1 + \frac{1}{n}\biggr)^{n}}^{b_{n}} = l.
  $$  
  De fato, como $a_{n} \geq{b_{n}}$ para todo n natural, segue do Teorema da Compara\c c\~ao que $e\geq{l}.$ Por outro lado, se 
 $n\geq{p}\geq{2},$
 $$
  b_{n} > 1 + 1 + \frac{1}{2!}(1-\frac{1}{n})+\cdots+\frac{1}{p!}(1-\frac{1}{n})(1-\frac{2}{n})\cdots(1-\frac{p-1}{n}).
 $$
 Agora, novamente pelo Teorema da Compara\c c\~ao, $l = \displaystyle \lim_{n\to\infty}b_{n}\geq{a_{p}}$ para todo natural p.
 Portanto, $l = \displaystyle \lim_{n\to\infty}b_{n}\geq{\sup\{a_{n}:n\in \mathbb{N}\}} = \lim_{n\to\infty}a_{n} = e.$ \qedsymbol
\end{example}
\begin{def*}
   Seja $\{a_{n}\} $ uma sequ\^encia. Um n\'umero real a \'e um valor de ader\^encia de $\{a_{n}\} $ se a sequ\^encia $\{a_{n}\}$ possui
uma subsequ\^encia convergente para a.$\quad\square$
\end{def*}
\begin{def*}
  Seja $\{a_{n}\} $  uma sequ\^encia limitada. Definimos o limite superior $\displaystyle\limsup_{n\to\infty}a_{n}(\text{ inferior}\liminf_{n\to\infty}a_{n})$ da 
sequ\^encia $\{a_{n}\} $ por 
  \begin{align*}
    &\limsup_{n\to\infty}a_{n} = \lim_{n\to\infty}\sup_{k\geq{n}}a_{k}\\
    &\liminf_{n\to\infty}a_{n} = \lim_{n\to\infty}\inf_{k\geq{n}}a_{k}\quad\square
  \end{align*}
\end{def*}
  Uma consequ\^encia direta do Teorema do Confronto que utiliza os conceitos acima nos permite dizer se uma sequ\^encia converge apenas
utilizando as ideias de limite superior e inferior:
 \begin{theorem*}
   Se a \'e um valor de ader\^encia da sequ\^encia $\{a_{n}\} $, ent\~ao
   $$
   \liminf_{n\to\infty}a_{n}\leq{a}\leq{\limsup_{n\to\infty}a_{n}}.
   $$
   Al\'em disso, uma sequ\^encia \'e convergente se, e somente se, $\liminf_{n\to\infty}a_{n} = \limsup_{n\to\infty}a_{n}.$
 \end{theorem*}
\newpage

\section{Aula 08 - 29/03/2023}
\begin{itemize}
  \item Rela\c c\~ao entre limite superior (e inferior), limites normais e valores de ader\^encia;
  \item Aproxima\c c\~oes sucessivas de valores.
\end{itemize}
\subsection{Limite Superior e Inferior}
\begin{def*}
   Seja $\{a_{n}\} $ uma sequ\^encia. Um n\'umero real a \'e um valor de ader\^encia de $\{a_{n}\} $ se a sequ\^encia $\{a_{n}\}$ possui
uma subsequ\^encia convergente para a.$\quad\square$
\end{def*}
\begin{def*}
  Seja $\{a_{n}\} $  uma sequ\^encia limitada. Definimos o limite superior $\displaystyle\limsup_{n\to\infty}a_{n}(\text{ inferior }\liminf_{n\to\infty}a_{n})$ da 
sequ\^encia $\{a_{n}\} $ por 
  \begin{align*}
    &\limsup_{n\to\infty}a_{n} = \lim_{n\to\infty}\sup_{k\geq{n}}a_{k} = \inf_{n\in \mathbb{N}}\sup_{k\geq{n}}a_{k}\\
    &\liminf_{n\to\infty}a_{n} = \lim_{n\to\infty}\inf_{k\geq{n}}a_{k} = \sup_{n\in \mathbb{N}}\inf_{k\geq{n}}a_{k}\quad\square
  \end{align*}
\end{def*}
  Uma consequ\^encia direta do Teorema do Confronto que utiliza os conceitos acima nos permite dizer se uma sequ\^encia converge apenas
utilizando as ideias de limite superior e inferior:
 \begin{theorem*}
   Se $\{a_{n}\} $ \'e uma sequ\^encia limitada, ent\~ao $a = \liminf_{n\to\infty}a_{n}$ e $b = \limsup_{n\to\infty}a_{n}$
   s\~ao valores de ader\^encia de $\{a_{n}\} $
 \end{theorem*}
\begin{proof*}
  A prova se baseia em verificar que, dada uma vizinhan\c ca $V_{a}$ de a, temos $a_{n}\in V_{a}$ para infinitos \'indices n. 
Dado $\epsilon > 0$, existe N natural tal que, colocando $a =\displaystyle \liminf_{n\to\infty} = \lim_{n\to\infty}\inf_{k\geq{n}}a_{k} = \lim_{n\to\infty}i_{n},$
  $$
    a - \epsilon < i_{n} < a + \epsilon \quad \forall n \geq{N}.
  $$
  Assim, existe um $a_{\overline{k}}, \overline{k}\geq{N}$ em $(a-\epsilon, a+\epsilon).$ Assim, existe $\overline{n} > \overline{k}$ tal que 
    $$
      a - \epsilon < i_{\overline{n}} < a + \epsilon.
    $$
    Repetindo o racioc\'icio, existe $a_{\overline{\overline{k}}}, \overline{\overline{k}} \geq{\overline{n}} > k$ em $(a - \epsilon, a + \epsilon).$
  Dando continuidade a este racioc\'inio ad infinitum, o teorema est\'a provado. \qedsymbol
\end{proof*}
 \begin{theorem*}
   Se a \'e um valor de ader\^encia da sequ\^encia $\{a_{n}\} $, ent\~ao
   $$
     \liminf_{n\to\infty}a_{n}\leq{a}\leq{\limsup_{n\to\infty}a_{n}}.
   $$
   Al\'em disso, uma sequ\^encia \'e convergente se, e somente se, $\liminf_{n\to\infty}a_{n} = \limsup_{n\to\infty}a_{n}.$
 \end{theorem*}
\begin{proof*}
  Defina $i_{n} = \inf_{k\geq{n}}a_{k}$. Segue que $i_{s(n)}\leq{a_{s(n)}}\overbracket[0pt]{\longrightarrow}^{\to }a$, pois o conjunto
  $\{a_{k}: k\geq{s(n)}\}$ cont\'em $a_{s(n)}$. Logo, como $i_{s(n)}$ converge para $\liminf_{n\to\infty}a_{n}$, segue do Teorema da compara\c c\~ao que 
    $$
      \liminf_{n\to\infty}a_{n} \leq{\lim_{n\to\infty}a_{s(n)} = a}.
    $$
    Analogamente, como $a_{s(n)}\leq{\sup_{k\geq{s(n)}}{(a_{k})}} = s_{s(n)}$ e $\sup_{k\geq{s(n)}}{(a_{k})}\overbracket[0pt]{\longrightarrow}^{n\to\infty}a_{n}$, 
  pelo teorema da compara\c c\~ao, chegamos novamente em 
    $$
      \lim_{n\to\infty}a_{s(n)} = a \leq{\limsup_{n\to\infty}a_{n}}.
    $$
    Portanto, juntando ambos, segue o resultado. \qedsymbol
\end{proof*}
  A seguir, mostraremos um m\'etodo para aproximar n\'umeros por meio de sequ\^encias de Cauchy.
\begin{theorem*}
  Se $\kappa\in{[0, 1)},\{a_{n}\} $ \'e uma sequ\^encia tal que, para todo $n\in \mathbb{N}, |a_{n+2}-a_{n+1}|\leq{\lambda|a_{n+1}-a_{n}|}$, 
  ent\~ao $\{a_{n}\} $ \'e de Cauchy.
\end{theorem*}
\begin{proof*}
  Se $m > n$ s\~ao naturais, m = n + p para algum natural n\~ao-nulo p. Assim, como 
    $$
      |a_{n+p} - a_{n}| = |a_{n+p} - a_{n+p-1} + a_{n+p-1}+\cdots + a_{n+1} - a_{n}|,
    $$
  segue da desigualdade triangular que
 \begin{align*}
   |a_{n+p}-a_{n}| &\leq{\kappa|a_{n+p-1}-a_{n+p-2}|+\kappa|a_{n+p-2}-a_{n+p-3}|+\cdots+\kappa|a_{n}-a_{n-1}|}\\ 
                   &\leq{\kappa^{n+p-1}|a_{n+p} - a_{n+p-1}|}+\cdots+{\kappa^{n}|a_{p}-a_{p-1}|}\\
                   &={\kappa^{n}\biggl[\kappa^{p-1}+\cdots+1\biggr]}|a_{1}-a_{0}| < \frac{\kappa^{n}}{1-\kappa}|a_{1}-a_{0}|.
 \end{align*}
 Dado $\epsilon > 0,$ escolha N natural tal que $\displaystyle \frac{\kappa^{n}}{1-\kappa}|a_{1}-a_{0}| < \epsilon.$ Assim, segue que,
 se $m, n\geq{N}, |a_{m}-a_{n}|<\epsilon$ e $\{a_{n}\} $ \'e de Cauchy. \qedsymbol
\end{proof*}
 \begin{example}
   Seja $a > 0$ e $\{a_{n}\}$ a sequ\^encia definida por $a_{0} = c > 0$ e $a_{n+1}=\displaystyle \frac{1}{2}\biggl(a_{n} + \frac{a}{a_{n}}\biggr)$. Mostre que 
   $\{a_{n}\}$ \'e convergente com limite $\sqrt{a}$. 

   Com efeito, observe que 
     $$
       a_{n+2} - a_{n+1} = \frac{1}{2}(a_{n+1}-a_{n}) + \frac{a}{2}\biggl(\frac{1}{a_{n+1}}-\frac{1}{a_{n}}\biggr) = \biggl(\frac{1}{2} - \frac{a}{2a_{n}a_{n+1}}\biggr)(a_{n+1}-a_{n})
     $$
     e note que, para todo $t > 0, \frac{1}{2}(t + \frac{a}{t}) > \sqrt{\frac{a}{2}}.$ Logo, $a_{n} > \sqrt{\frac{a}{2}}$ para todo n maior ou igual que 1.
Disto segue que $2a_{n}a_{n+1} > a$ e que 
  $$
    \biggl|\frac{1}{2} - \frac{a}{2a_{n}a_{n+1}}\biggr| < \frac{1}{2}.
  $$
  Portanto, segue do m\'etodo das aproxima\c c\~oes sucessivas que $\{a_{n}\}$ \'e convergente e seu limite l deve satisfazer 
 $l = \frac{1}{2}(l + \frac{a}{l})$, ou seja, $l^{2} = a.$ \qedsymbol
 \end{example}
\newpage

\section{Aula 09 - 31/03/2023}
\subsection{Motiva\c c\~oes}
  \begin{itemize}
    \item Limites Infinitos;
    \item Sequ\^encias divergentes.
  \end{itemize}

\subsection{Pontos de Ader\^encia e Limites Superiores/Inferiores.}
\begin{def*}
   Seja $\{a_{n}\} $ uma sequ\^encia. Um n\'umero real a \'e um valor de ader\^encia de $\{a_{n}\} $ se a sequ\^encia $\{a_{n}\}$ possui
uma subsequ\^encia convergente para a.$\quad\square$
\end{def*}
\begin{def*}
  Seja $\{a_{n}\} $  uma sequ\^encia limitada. Definimos o limite superior $\displaystyle\limsup_{n\to\infty}a_{n}(\text{ inferior }\liminf_{n\to\infty}a_{n})$ da 
sequ\^encia $\{a_{n}\} $ por 
  \begin{align*}
    &\limsup_{n\to\infty}a_{n} = \lim_{n\to\infty}\sup_{k\geq{n}}a_{k} = \inf_{n\in \mathbb{N}}\sup_{k\geq{n}}a_{k}\\
    &\liminf_{n\to\infty}a_{n} = \lim_{n\to\infty}\inf_{k\geq{n}}a_{k} = \sup_{n\in \mathbb{N}}\inf_{k\geq{n}}a_{k}\quad\square
  \end{align*}
\end{def*}
  Uma consequ\^encia direta do Teorema do Confronto que utiliza os conceitos acima nos permite dizer se uma sequ\^encia converge apenas
utilizando as ideias de limite superior e inferior:
 \begin{theorem*}
   Se $\{a_{n}\} $ \'e uma sequ\^encia limitada, ent\~ao $a = \liminf_{n\to\infty}a_{n}$ e $b = \limsup_{n\to\infty}a_{n}$
   s\~ao valores de ader\^encia de $\{a_{n}\} $
 \end{theorem*}
\begin{proof*}
  A prova se baseia em verificar que, dada uma vizinhan\c ca $V_{a}$ de a, temos $a_{n}\in V_{a}$ para infinitos \'indices n. 
Dado $\epsilon > 0$, existe N natural tal que, colocando $a =\displaystyle \liminf_{n\to\infty} = \lim_{n\to\infty}\inf_{k\geq{n}}a_{k} = \lim_{n\to\infty}i_{n},$
  $$
    a - \epsilon < i_{n} < a + \epsilon \quad \forall n \geq{N}.
  $$
  Assim, existe um $a_{\overline{k}}, \overline{k}\geq{N}$ em $(a-\epsilon, a+\epsilon).$ Assim, existe $\overline{n} > \overline{k}$ tal que 
    $$
      a - \epsilon < i_{\overline{n}} < a + \epsilon.
    $$
    Repetindo o racioc\'icio, existe $a_{\overline{\overline{k}}}, \overline{\overline{k}} \geq{\overline{n}} > k$ em $(a - \epsilon, a + \epsilon).$
  Dando continuidade a este racioc\'inio ad infinitum, segue que $V_{a}$ cont\'em $a_{n}$ para infinitos ind\'ices n, tal que o teorema est\'a provado. \qedsymbol
\end{proof*}
\begin{crl*}
  (Exerc\'icio) Nas condi\c c\~oes do teorema, a \'e um valor de ader\^encia de $\{a_{n}\}$ e $b = \limsup_{n\to\infty}a_{n}$ tamb\'em \'e um valor de
  ader\^encia de $\{a_{n}\}$
\end{crl*}
 \begin{theorem*}
   Se a \'e um valor de ader\^encia da sequ\^encia $\{a_{n}\} $, ent\~ao
   $$
     \liminf_{n\to\infty}a_{n}\leq{a}\leq{\limsup_{n\to\infty}a_{n}}.
   $$
   Al\'em disso, uma sequ\^encia \'e convergente se, e somente se, $\liminf_{n\to\infty}a_{n} = \limsup_{n\to\infty}a_{n}.$
 \end{theorem*}
\begin{proof*}
  Defina $i_{n} = \inf_{k\geq{n}}a_{k}$. Segue que $i_{s(n)}\leq{a_{s(n)}}\overbracket[0pt]{\longrightarrow}^{\to }a$, pois o conjunto
  $\{a_{k}: k\geq{s(n)}\}$ cont\'em $a_{s(n)}$. Logo, como $i_{s(n)}$ converge para $\liminf_{n\to\infty}a_{n}$, segue do Teorema da compara\c c\~ao que 
    $$
      \liminf_{n\to\infty}a_{n} \leq{\lim_{n\to\infty}a_{s(n)} = a}.
    $$
    Analogamente, como $a_{s(n)}\leq{\sup_{k\geq{s(n)}}{(a_{k})}} = s_{s(n)}$ e $\sup_{k\geq{s(n)}}{(a_{k})}\overbracket[0pt]{\longrightarrow}^{n\to\infty}a_{n}$, 
  pelo teorema da compara\c c\~ao, chegamos novamente em 
    $$
      \lim_{n\to\infty}a_{s(n)} = a \leq{\limsup_{n\to\infty}a_{n}}.
    $$
    Portanto, juntando ambos, segue o resultado. \qedsymbol
\end{proof*}

\subsection{Sequ\^encias Divergente para $\pm\infty$.}
 Recorde-se que 
\begin{def*}
  Diremos que uma sequ\^encia $\{a_{n}\}$ diverge para $+\infty(-\infty)$ se, dado $M >0,$ existe um natural N tal que $a_{n}\geq{M}(a_{n}\leq{-M})$
para todo $n\geq{N}.$ Escreveremos $\lim_{n\to\infty}a_{n} = +\infty(-\infty),$ ou $a_{n}\overbracket[0pt]{\longrightarrow}^{n\to \infty}+\infty(-\infty)\square.$
\end{def*}
 \begin{def*}
   Diremos que a sequ\^encia $\{a_{n}\}$ \'e eventualmente positiva (negativa) se existe um natural N tal que $a_{n} > 0 (a_{n} < 0)$
para todo $n\geq{N}.\square$
 \end{def*}
  Vejamos a seguir algumas das propriedades dessas sequ\^encias.
 \begin{theorem*}
  \begin{itemize}
    \item[a)] Se $a_{n}\overbracket[0pt]{\longrightarrow}^{n\to \infty}\infty$ e $\{b_{n}\}$ \'e limitada inferiormente, ent\~ao
  $a_{n}+b_{n}\overbracket[0pt]{\longrightarrow}^{n\to \infty}\infty.$
    \item[b)] Se $a_{n}\overbracket[0pt]{\longrightarrow}^{n\to \infty}\infty$ e $\liminf_{n\to\infty}b_{n} > 0,$ ent\~ao
  $\lim_{n\to\infty}a_{n}b_{m} = +\infty$
  \item[c)] Seja $\{a_{n}\}$ uma sequ\^encia com $a_{n}\neq0$ para todo n natural. $\{a_{n}\}$ \'e eventualmente negativa e
  $a_{n}\overbracket[0pt]{\longrightarrow}^{n\to \infty}0$ se, e somente se, $\displaystyle \frac{1}{a_{n}}\overbracket[0pt]{\longrightarrow}^{n\to \infty}\infty.$
  \item[d)] Sejam $\{a_{n}\},\{b_{n}\} $ sequ\^encias eventualmente positivas, $b_{n}\neq 0$ para todo n natural.
    \item[d.1)]Se $\liminf_{n\to\infty}a_{n} > 0$ e $b_{n}\overbracket[0pt]{\longrightarrow}^{n\to \infty}0$, ent\~ao $\displaystyle\lim_{n\to\infty}\frac{a_{n}}{b_{n}} = +\infty.$
    \item[d.2)] Se $\{a_{n}\}$ \'e limitada e $b_{n}\overbracket[0pt]{\longrightarrow}^{n\to \infty}\infty$, ent\~ao $\displaystyle \frac{a_{n}}{b_{n}}\overbracket[0pt]{\longrightarrow}^{n\to \infty}0.$
  \end{itemize}
 \end{theorem*}
\begin{proof*}
  (a) $\Rightarrow)$ Como $\{b_{n}\}$ \'e limitada inferiormente, existe um n\'umero real $l > 0$ tal que $b_{n} \geq{ -l}$ para todo 
n natural. Como $a_{n}\overbracket[0pt]{\longrightarrow}^{n\to \infty}\infty,$ dado M positivo, existe N natural tal que 
 $a_{n}\geq{M}+l$ para todo $n\geq{N}.$ Logo, 
   $$
     a_{n} + b_{n}\geq{M+l-l} = M,\quad \forall n\geq{N}.
   $$
   Portanto, $a_{n} + b_{n}\overbracket[0pt]{\longrightarrow}^{n\to \infty}\infty.$

  (b) $\Rightarrow)$ Como $a_{n}\overbracket[0pt]{\longrightarrow}^{n\to \infty}\infty$ e $\liminf_{n\to\infty}b_{n} = r > 0,$
  existe $N_{1}$ natural tal que $\liminf_{k\geq{n}}a_{n} \geq{\frac{r}{2}}$ para todo $n\geq{N_{1}.}$ Como $a_{n}\overbracket[0pt]{\longrightarrow}^{n\to \infty}\infty,$
  dado M positivo, existe $N_{2}$ natural tal que $a_{n} > \frac{2M}{r}$ para todo $n\geq{N_{2}}.$ Disto segue que, para $n\geq{N = \max\{N_{1}, N_{2}\}}$, 
    $$
      a_{n}b_{n}\geq{\frac{2M}{r}\frac{r}{2}} = M,\quad \forall n\geq{N}.
    $$
  donde segue o que quer\'iamos.

  (c) $\Rightarrow)$ Se $\{a_{n}\}$ \'e infint\'esima e eventualmente positiva, dado M positivo, seja N natural tal que 
 $0 < a_{n} < \frac{1}{M}$ para quaisquer $n\geq{N}.$ Logo, $\frac{1}{a_{n}} > M$ para todos os $n\geq{N},$ mostrando que 
 $\frac{1}{a_{n}}\overbracket[0pt]{\longrightarrow}^{n\to \infty}\infty.$

 Reciprocamente, se $\frac{1}{a_{n}}\overbracket[0pt]{\longrightarrow}^{n\to \infty}\infty,$ dado $\epsilon > 0,$ seja N natural tal que
 $\frac{1}{a_{n}} > \frac{1}{\epsilon}$ para todo $n\geq{N}.$ Desta forma, $0 < a_{n} < \epsilon $ para todo $n\geq{N},$ provando o 
 resultado.

 (d.1) $\Rightarrow)$ De fato, se $\liminf_{n\to\infty}a_{n} = r > 0,$ existe $N_{1}$ natural tal que $a_{n}\geq{\frac{r}{2}}$ para todo $n\geq{N_{1}.}$
Dado M positivo, seja $N_{2}$ outro natural tal que $0 < a_{n} < \frac{r}{2M}$ para todo $n\geq{N_{2}}.$ Logo, $\frac{a_{n}}{b_{n}}>\frac{r}{2}\frac{2M}{r} = M$
para todo $n\geq{N} = \max\{N_{1}, N_{2}\}$, tal que $\frac{a_{n}}{b_{n}}\overbracket[0pt]{\longrightarrow}^{n\to \infty}\infty$.

  (d.2) $\Rightarrow)$ Seja $L > 0 $ tal que $|a_{n}|\leq{L}$ para todo n natural. Como $b_{n}\overbracket[0pt]{\longrightarrow}^{n\to \infty}\infty,$
  dado $\epsilon > 0,$ existe N natural tal que $b_{n} > \frac{L}{\epsilon}(\frac{1}{b_{n}} < \frac{\epsilon}{L})$ para todo $n\geq{N}.$ Logo,
  $$
    \biggl|\frac{a_{n}}{b_{n}} - 0\biggr| < L \cdot \frac{\epsilon}{L} = \epsilon,\quad \forall n\geq{N},
  $$
  mostrando que $\frac{a_{n}}{b_{n}}\overbracket[0pt]{\longrightarrow}^{n\to \infty}0.$\qedsymbol
\end{proof*}
  \'E importante notar que, se $a_{n}\overbracket[0pt]{\longrightarrow}^{n\to \infty}\infty$ e $b_{n}\overbracket[0pt]{\longrightarrow}^{n\to \infty}-\infty$, nada
  podemos afirmar de $\lim_{n\to\infty}(a_{n}+b_{n}).$ Neste caso, tudo pode ocorrer! $\{a_{n} + b_{n}\} $ pode convergir para qualquer n\'umero real,
  divergir para $+\infty, -\infty$ ou pode oscilar. Vamos ilustrar a situa\c c\~ao.
 \begin{example}
   Se $a_{n} = \sqrt{n+1}, b_{n} = -\sqrt{n},$ para todo n natural, \'e f\'acil de ver que $a_{n}\overbracket[0pt]{\longrightarrow}^{n\to \infty}\infty$ e $b_{n}\overbracket[0pt]{\longrightarrow}^{n\to \infty}-\infty$.
   Para ver o que ocorre com a sequ\^encia $\{a_{n}+b_{n}\} $, observe que 
     $$
     \sqrt{n+1} - \sqrt{n} = \frac{(\sqrt{n+1}-\sqrt{n})(\sqrt{n+1}+\sqrt{n})}{\sqrt{n+1}+\sqrt{n}} = \frac{1}{\sqrt{n+1}+\sqrt{n}}.
     $$ 
     Segue de (d.2) que $\{a_{n}+b_{n}\} $ \'e infinit\'esima.
 \end{example}
\begin{example}
  Se $a > 1,$ ent\~ao a sequ\^encia $\{a_{n}\}$ dada por $a_{n} = \frac{a^{n}}{n}$ diverge para $+\infty.$ De fato, basta ver que
a = 1 + h, com h positivo, e escrever 
  $$
    \frac{a^{n}}{n} = \frac{(1+h)^{n}}{n} = \frac{1}{n} + h + (n-1)\frac{h^{2}}{2!} + s_{n}.
  $$
  O resultado segue aplicando (a).
\end{example}
 \begin{example}
   Se $a > 1$, ent\~ao a sequ\^encia $\{a_{n}\}$ com $a_{n} = \frac{n!}{a^{n}}$ diverge para $+\infty.$ Com efeito, basta escolher
 $n_{0}$ tal que $\frac{n_{0}}{a} > 2$ e escrever, para $n\geq{n_{0}}, a_{n} = \frac{n_{0}!}{a^{n_{0}}}\frac{n!}{n_{0}!}\frac{1}{a^{n-n_{0}}}.$
 Se $r = \frac{n_{0}!}{a^{n_{0}}}\frac{n!}{n_{0}!},$ temos 
   $$
   a_{n} = r \frac{n(n-1)\cdots(n_{0}+1)}{a^{n-n_{0}}} = r2^{n-n_{0}} + s_{n} = r(n+1 - n_{0}) + \tilde{s}_{n}
   $$
   Novamente, o resultado segue aplicando (a).
 \end{example}
  Por fim, algumas outras propriedades que n\~ao foram inclusas no teorema s\~ao deixadas como exerc\'icio ao leitor.
 \begin{itemize}
  \item[a)] Se $a_{n}\overbracket[0pt]{\longrightarrow}^{n\to \infty}-\infty$ e $\{b_{n}\}$ \'e limitada superiormente, ent\~ao
  $a_{n}+b_{n}\overbracket[0pt]{\longrightarrow}^{n\to \infty}-\infty.$
  \item[b)] Se $a_{n}\overbracket[0pt]{\longrightarrow}^{n\to \infty}-\infty$ e $\liminf_{n\to\infty}b_{n} > 0,$ ent\~ao $\lim_{n\to\infty}a_{n}b_{m} = -\infty.$
  \item[c)] Seja $\{a_{n}\}$ uma sequ\^encia com $a_{n}\neq 0$ para todo n natural. $\{a_{n}\}$ \'e eventualmente negativa e $a_{n}\overbracket[0pt]{\longrightarrow}^{n\to \infty}0$
    se, e somente se, $\frac{1}{a_{n}}\overbracket[0pt]{\longrightarrow}^{n\to \infty}-\infty.$
  \item[d)] Sejam $\{a_{n}\},\{b_{n}\}$ duas sequ\^encias eventualmente negativas com $b_{n}\neq0$ para todo n natural.
    \item[d.1)] Se $\liminf_{n\to\infty}a_{n} < 0$ e $b_{n}\overbracket[0pt]{\longrightarrow}^{n\to \infty}0$, ent\~ao $\lim_{n\to\infty}\frac{a_{n}}{b_{n}} = +\infty.$
    \item[d.2)] Se $\{a_{n}\}$ \'e limitada e $b_{n}\overbracket[0pt]{\longrightarrow}^{n\to \infty}-\infty$, ent\~ao $\frac{a_{n}}{b_{n}}\overbracket[0pt]{\longrightarrow}^{n\to \infty}0.$
    \item[e)] No item (d), analise a situa\c c\~ao em que $\{a_{n}\}$ \'e eventualmente positiva e $\{b_{n}\}$ \'e eventualmente negativa.
 \end{itemize}
\newpage

\section{Aula 10 - 10/04/2023}
\subsection{Motiva\c c\~oes}
 \begin{itemize}
 \item S\'eries;
 \item S\'eries de n\'umeros positivos;
 \item Testes da ra\'iz e da raz\~ao.
 \end{itemize}

\subsection{S\'eries de N\'umeros Reais}
  S\'eries s\~ao tipos particulares de sequ\^encias que possuem suas pr\'oprias propriedades extras. Vamos come\c car com um exemplo, chamado s\'erie harm\^onica (nome derivado da m\'edia harm\^onica).
 \begin{example}
   Considere $s_{n} = 1 + \frac{1}{2} + \cdots + \frac{1}{n}.$ Provemos que $\{s_{n}\}\overbracket[0pt]{\longrightarrow}^{n\to \infty}\infty.$
   Com efeito, como $\{s_{n}\} $ \'e crescente, basta mostrar que ela \'e ilimitada: 
  \begin{align*}
    &s_{1} = 1\\
    &s_{2} = 1 + \frac{1}{2} = 3\frac{1}{2}\\
    &s_{3} = 1 + \frac{1}{2} + \frac{1}{3}\\
    &s_{4} = 1 + \frac{1}{2} + \frac{1}{3} + \frac{1}{4} > 1 + \frac{1}{2} + \frac{1}{2} = 4\frac{1}{2}\\
    &s_{8} = 1 + \frac{1}{2} + \underbrace{\frac{1}{3} + \frac{1}{4}}_{> \frac{1}{2}} + \underbrace{\frac{1}{5} + \frac{1}{6} + \frac{1}{7} + \frac{1}{8}}_{4\frac{1}{8}} > 1 + \frac{1}{2} + \frac{1}{2} + \frac{1}{2} = 5\frac{1}{2}\\
    &s_{16} = s_{8} + \frac{1}{9} +\cdots + \frac{1}{16} > 5\frac{1}{2} + \frac{1}{2} = 6\frac{1}{2}.\\
    \vdots
  \end{align*}
  Analisando o padr\~ao de repeti\c c\~ao desta s\'erie, chegamos em 
    $$
      s_{2^{n}} > (n+2)\frac{1}{2} \Rightarrow s_{2^{n+1}} > s_{2^n} + \frac{1}{2^{n}+1} + \cdots + \frac{1}{2^{n}+ 2^{n}} > (n+2)\frac{1}{2} + \frac{2^{n}}{2^{n+1}} = ((n+1)+2)\frac{1}{2}.
    $$
    Logo, por indu\c c\~ao, $s_{2^{k}} > (k+2)\frac{1}{2}$ para todo k natural, mostrando que $\{s_{n}\} $ n\~ao \'e limitada superiormente.
  Portanto, $\lim_{n\to\infty}s_{n} = \infty.$ \qedsymbol
 \end{example}
 Consideremos a sequ\^encia $\{a_{n}\}$. A partir da sequ\^encia $\{a_{n}\} $, vamos construir a sequ\^encia das somas parciais da seguinte forma: 
\begin{align*}
  &s_{1} = a_{1}\\
  &s_{2} = a_{1} + a_{2}\\
  &s_{3} = a_{1} + a_{2} + a_{3}\\
  &\vdots\\
  &s_{n} = a_{1} + a_{2} + a_{3} + \cdots + a_{n}\\
  &\cdots.
\end{align*}
 \begin{def*}
  A sequ\^encia $\{s_{n}\}$ das somas paricais \'e chamada s\'erie associada a $\{a_{n}\}.$ Cada $s_{n}$ \'e chamado soma parcial da s\'erie, e cada
 $a_{n}$ leva o nome de termo da s\'erie. As nota\c c\~oes s\~ao: 
   $$
   \sum\limits_{n\geq{1}}^{}a_{n},\quad \sum\limits_{}^{}a_{n},\quad \sum\limits_{n=1}^{\infty}a_{n},\quad\text{ ou } a_{1}+a_{2}+\cdots+a_{n}+\cdots.
   $$
 \end{def*}
 Observe que, \`as vezes, consideraremos s\'eries que comecem em algum termo $n_{0}$ ao inv\'es do termo 1. Neste caso, escreveremos 
 $\sum\limits_{n=n_{0}}^{\infty}a_{n}.$
\begin{example}
 \begin{align*}
   &\{a_{n}\} = \{(-1)^{n+1}\}.\text{ A sequ\^encia das somas parciais ser\'a:}\\
   &s_{1} = a_{1} = 1\\
   &s_{2} = a_{1} + a_{2} = 1 - 1 = 0\\
   &\vdots\\
   &s_{n} = 0\text{ ou } 1.
 \end{align*} 
\end{example}
\begin{example}
  Construimos a s\'erie $\{a_{n}\} = \biggl\{\frac{1}{n}\biggr\}$ no exemplo anterior.
\end{example}
\begin{example}
  $\{a_{n}\} = \biggl\{\frac{6}{10^{n}}\biggr\}$ \'e uma s\'erie que, quando constru\'ida, converge para $\frac{2}{3}.$
\end{example}
\begin{def*}
  Diremos que uma s\'erie \'e convergente se a sequ\^encia das somas parciais \'e convergente. Caso contr\'ario, ser\'a dita divergente.
  Se a sequ\^encia $\{s_{n}\} $ \'e convergente para S, dizemos que a s\'erie $\sum\limits_{1}^{\infty}a_{n}$ \'e convergente com soma S. $\square$
\end{def*}
  A soma, multiplica\c c\~ao, etc. De s\'eries \'e definida como no caso das sequ\^encias. Denotaremos s\'eries convergentes por 
    $$
      \sum\limits_{n=1}^{\infty}a_{n} = S = \lim_{n\to\infty}s_{n} = \lim_{n\to\infty}\biggl(\sum\limits_{k=1}^{n}a_{k}\biggr)
    $$
   \begin{example}
     A s\'erie telesc\'opica \'e dada por 
       $$
         \sum\limits_{n=1}^{\infty}\frac{1}{n(n+1)} = 1.
       $$
     Observe que 
     \begin{align*}
       s_{n} &= \frac{1}{1.2} + \frac{1}{2.3} + \cdots + \frac{1}{n(n+1)}\\
             &= (1-\frac{1}{2}) + (\frac{1}{2}-\frac{1}{3}) + \cdots + (\frac{1}{n} - \frac{1}{n+1})\\
             &= 1 - \frac{1}{n+1}.
      \end{align*}
    Portanto, $\lim_{n\to\infty}s_{n} = \lim_{n\to\infty}(1 - \frac{1}{n+1}) = 1 - 0 = 1.$ \qedsymbol
   \end{example}
  \begin{example}
   \begin{itemize}
     \item[i)]$\sum\limits_{1}^{\infty}(-1)^{n}$ \'e divergente, visto que, para termos \'impares da soma parcial, ela vale 0, mas para termos pares,
    ela vale 1.
     \item[ii)]$\sum\limits_{1}^{\infty}2^{n}$ diverge, pois a soma parcial dos termos dela n\~ao \'e limitada.
     \item[iii)]$\sum\limits_{n=1}^{\infty}\frac{1}{n}$ diverge, e recebe o nome de s\'erie harm\^onica. Aqui, $s_{n} = 1 + \frac{1}{2}+\cdots+\frac{1}{n}.$
   \end{itemize}
  \end{example}
  Algumas s\'eries s\~ao importantes, pois podem ser utilizadas para obter informa\c c\~oes sobre outras atrav\'es da compara\c c\~ao, tais como a 
  s\'erie telesc\'opica e a harm\^onica. Outra importante \'e a s\'erie geom\'etrica:
 \begin{example}
   Definimos ela como $\sum\limits_{n\geq{1}}^{}ar^{n-1} = a + ar + ar^{2} + \cdots$. Afirmamos que ela \'e convergente se, e s\'o se,
  $|r| < 1$, convergindo para $\frac{a}{1-r}$. Assim, 
    $$
      a + ar + ar^{2} + \cdots + ar^{n} + \cdots = \frac{a}{1-r}, \quad |r| < 1.
    $$
    Com efeito, se r = 1, ent\~ao $s_{n} = a + a + a + \cdots + a = na,$ que tende a infinito ou menos infinito, dependendo do sinal de a, mostrando sua diverg\^encia.
Agora, se r for diferente de 1, temos: 
  $$
    s_{n} = a + ar + ar^{2} + \cdots + ar^{n-1},\quad rs_{n} = ar + ar^{2} + ar^{3} + \cdots ar^{n}.
  $$
  Subtraindo membro a membro, obtemos 
    $$
      s_{n}(1-r) = a - ar^{n} = a(1-r^{n}).
    $$
    Portanto, $s_{n} = a \frac{1-r^{n}}{1-r} = \frac{a}{1-r} - \frac{a}{1-r}r^{n}.$ Se $|r|<1$, j\'a vimos anteriormente que
  $r^{n}\overbracket[0pt]{\longrightarrow}^{n\to \infty}0,$ tal que 
    $$
      \lim_{n\to\infty}s_{n} = \lim_{n\to\infty}\biggl(\frac{a}{1-r} - \frac{a}{1-r}r^{n}\biggr) = \frac{a}{1-r}.
    $$
    Caso $|r| > 1,$ ou r = -1, vimos anteriormente que $r^{n}$ \'e divergente, provando o que quer\'iamos. \qedsymbol
 \end{example}
\begin{theorem*}
  Se $\sum\limits_{}^{}a_{n}$ \'e uma s\'erie convergente, ent\~ao $\lim_{n\to\infty}a_{n} = 0.$ A rec\'iproca \'e falsa.
\end{theorem*}
\begin{proof*}
  Note que, se $\{s_{n}\}$ a sequ\^encia das somas parciais dos an's \'e convergente, temos 
    $$
      a_{n} = s_{n} - s_{n-1}\overbracket[0pt]{\longrightarrow}^{n\to \infty} s-s= 0,
    $$
    provando o resultado. Um contraexemplo da volta \'e a s\'erie harm\^onica. \qedsymbol
\end{proof*}
 \begin{example}
   Exerc\'icio: $\sum\limits_{n=1}^{\infty}a_{n}$ \'e uma s\'erie convergente se, e somente se, $\sum\limits_{n=n_{0}}^{\infty}a_{n}$ \'e convergente.
 \end{example}
\begin{theorem*}
  Se $a_{n}\geq{0}$ para todo n natural, $\sum\limits_{}^{}a_{n}$ \'e uma s\'erie convergente se, e somente se, a sequ\^encia das somas paricais \'e limitada.
\end{theorem*}
\begin{theorem*}
  Sejam $\sum\limits_{}^{}a_{n}, \sum\limits_{}^{}b_{n}$ s\'eries de termos positivos. Se existem c positivo e $n_{0}$ natural tais que
 $a_{n} \leq{cb_{n}}$ para todo $n\geq{n_{0}},$ ent\~ao 
\begin{itemize}
  \item[i)] Se $\sum\limits_{}^{}b_{n}$ \'e convergente, ent\~ao $\sum\limits_{}^{}a_{n}$ \'e convergente.
    \item[ii)] Se $\sum\limits_{}^{}a_{n}$ \'e divergente, ent\~ao $\sum\limits_{}^{}b_{n}$ \'e divergente. 
\end{itemize}
\end{theorem*}
\begin{example}
  Se $r > 1, \sum\limits_{}^{}\frac{1}{n^{r}}$ \'e convergente. De fato, note que 
  \begin{align*}
    s_{2^{n}-1} &< 1 + \biggl(\frac{1}{2^{r}}+\frac{1}{3^{r}}\biggr) + \biggl(\frac{1}{4^{r}} + \frac{1}{5^{r}} + \frac{1}{6^{r}}+\frac{1}{7^{r}}\biggr) + \cdots\\
                &+\Biggl(\frac{1}{2^{n}-2^{n-1}}+\cdots+\frac{1}{(2^{n}-2)^{r}} + \frac{1}{(2^{n}-1)^{r}}\Biggr) \\
                &< 1 + \frac{2}{2^{r}} + \frac{4}{4^{r}} + \cdots + \frac{2^{n-1}}{2^{(n-1)r}} \\
                &= 1 + \frac{1}{2^{r-1}} + \frac{1}{4^{r-1}} + \cdots + \frac{1}{2^{(n-1)(r-1)}}\leq{\frac{1}{1-2^{-r+1}}}.
  \end{align*}
\end{example}
Como uma sequ\^encia convergente \'e equivalente a uma de Cauchy, vale o seguinte resultado:
\begin{theorem*}
  A s\'erie $\sum\limits_{}^{}a_{n}$ \'e convergente se, e somente se, dado $\epsilon > 0$, existe um natural N tal que 
    $$
      |s_{n+p} - s_{n}| = |a_{n} + \cdots + a_{n+p}| < \epsilon
    $$
    para todo $n\geq{N}$ e todo natural p.
\end{theorem*}
 \begin{def*}
   Uma s\'erie $\sum\limits_{}^{}a_{n}$ \'e absolutamente convergente quando $\sum\limits_{}^{}|a_{n}|$ \'e convergente.$\quad\square.$
 \end{def*}
 Segue do crit\'erio de Cauchy que 
\begin{theorem*}
  Se $\sum\limits_{}^{}a_{n}$ \'e absolutamente convergente, ent\~ao $\sum\limits_{}^{}a_{n}$ \'e convergente. A volta n\~ao vale.
\end{theorem*}
  Com rela\c c\~ao a volta, considere a s\'erie $\sum\limits_{}^{}\frac{(-1)^{n+1}}{n}$ n\~ao \'e absolutamente convergente, mas converge. 
  De fato, note que 
    $$
      s_{2} = 1 - \frac{1}{2}, s_{4} = (1-\frac{1}{2}) + (\frac{1}{3}-\frac{1}{4}), s_{6} = (1-\frac{1}{2}) + (\frac{1}{3}-\frac{1}{4}) + (\frac{1}{5}-\frac{1}{6}), \cdots
    $$
    Assim, $s_{2}<s_{4}<\cdots<s_{2n}<\cdots$ e $\{s_{2n}\}$ \'e crescente e limitada por 1, tal que ela converge. Por outro lado, 
      $$
        s_{1} = 1, s_{3} = 1 + (-\frac{1}{2} + \frac{1}{3}), s_{5} = 1 + (-\frac{1}{2}+\frac{1}{3}) + (-\frac{1}{4}+\frac{1}{5}), \cdots
      $$ 
      Deste forma, $s_{1} > s_{3} > s_{5} > \cdots > s_{2(n-1)} > \cdots$, ou seja, $\{s_{2(n-1)}\}$ \'e decrescente limitada, portanto convergente.

      Por outro lado, $s_{2n+1}-s_{n} = \frac{1}{2n+1}\overbracket[0pt]{\longrightarrow}^{n\to \infty}0,$ mostrando a converg\^encia da s\'erie.
 \begin{example}
   Exerc\'icio: Seja $\{a_{n}\}$ uma sequ\^encia infinit\'esima de termos n\~ao negativos que \'e decrescente. Mostre que $\sum\limits_{}^{}(-1)^{n}a_{n}$
  \'e convergente.
 \end{example}
\begin{example}
  Exerc\'icios: Seja $\sum\limits_{}^{}b_{n}$ uma s\'erie convergente de termos n\~ao-negativos. Se existem k positivo e $n_{0}$
natural tais que $|a_{n}|\leq{kb_{n}}$ para todo $n\geq{n_{0}},$ ent\~ao a s\'erie $\sum\limits_{}^{}a_{n}$ \'e absolutamente convergente.

  Se existem $c\in(0, 1), k > 0$ e $n_{0}$ natural tais que $|a_{n}|\leq{kc^{n}}$ para todo $n\geq{n_{0}},$ ent\~ao a s\'erie 
 $\sum\limits_{}^{}a_{n}$ \'e absolutamente convergente.
\end{example}

\subsection{Testes da Raz\~ao e da Ra\'iz}
 \begin{theorem*}
   Se $\{a_{n}\}$ \'e uma sequ\^encia limitada e $\limsup_{n\to\infty}|a_{n}|^{\frac{1}{n}} = c < 1$, ent\~ao $\sum\limits_{}^{}a_{n}$
  \'e absolutamente convergente.
 \end{theorem*}
\begin{proof*}
  Existe N natural tal que $\sup_{k\geq{n}}|a_{k}|^{\frac{1}{k}} < r = \frac{c+1}{2} < 1$ para todo natural n. Logo, $|a_{n}|<r^{n}$ para
todo $n\geq{N}.$ Segue do Teorema da Compara\c c\~ao que $\sum\limits_{}^{}|a_{n}|$ \'e convergente, ou seja, $\sum\limits_{}^{}a_{n}$ \'e absolutamente convergente. \qedsymbol
\end{proof*}
\begin{example}
  Se p \'e natural, ent\~ao $\sum\limits_{}^{}n^{p}a^{n}$ \'e convergente para $|a|< 1$ e divergente para $|a_{n}| \geq{1}.$ De fato,
basta ver que $\limsup_{n\to\infty}|n^{p}a^{n}|^{\frac{1}{n}} = |a| < 1$ e aplicar o teste da ra\'iz. Para ver que a s\'erie \'e divergente 
quando o m\'odulo de a \'e maior ou igual a 1, basta notar que a sequ\^encia dos termos da s\'erie n\~ao converge para zero neste caso. \qedsymbol
\end{example}
\newpage

\section{Aula 11 - 12/04/2023}
\subsection{Motiva\c c\~oes}
\begin{itemize}
  \item Teste da Raz\~ao;
  \item Teoremas de Dirichlet, Abel e Leibniz;
  \item S\'eries Rearranjadas.
\end{itemize}
\subsection{Teste da Ra\'iz e da Raz\~ao.}
  Comecemos revisitando o Teste da Ra\'iz.
\begin{theorem*}
   Se $\{a_{n}\}$ \'e uma sequ\^encia limitada e $\limsup_{n\to\infty}|a_{n}|^{\frac{1}{n}} = c < 1$, ent\~ao $\sum\limits_{}^{}a_{n}$
  \'e absolutamente convergente.
 \end{theorem*}
\begin{proof*}
  Existe N natural tal que $\sup_{k\geq{n}}|a_{k}|^{\frac{1}{k}} < r = \frac{c+1}{2} < 1$ para todo natural n. Logo, $|a_{n}|<r^{n}$ para
todo $n\geq{N}.$ Segue do Teorema da Compara\c c\~ao que $\sum\limits_{}^{}|a_{n}|$ \'e convergente, ou seja, $\sum\limits_{}^{}a_{n}$ \'e absolutamente convergente. \qedsymbol
\end{proof*}
\begin{example}
  Se p \'e natural, ent\~ao $\sum\limits_{}^{}n^{p}a^{n}$ \'e convergente para $|a|< 1$ e divergente para $|a_{n}| \geq{1}.$ De fato,
basta ver que $\limsup_{n\to\infty}|n^{p}a^{n}|^{\frac{1}{n}} = |a| < 1$ e aplicar o teste da ra\'iz. Para ver que a s\'erie \'e divergente 
quando o m\'odulo de a \'e maior ou igual a 1, basta notar que a sequ\^encia dos termos da s\'erie n\~ao converge para zero neste caso. \qedsymbol
\end{example}
  Como esperado, o resultado seguinte \'e o que conhecemos como Teste da Raz\~ao, e costuma ser mais simples de aplicar na maioria dos casos.
\begin{theorem*}
  Se $\sum\limits_{}^{}b_{n}$ \'e uma s\'erie convergente de termos positivos e $\sum\limits_{}^{}a_{n}$ \'e uma s\'erie de termos
n\~ao-nulos tais que existe $n_{0}\in \mathbb{N}$ satisfazendo 
  $$
    \frac{|a_{n+1}|}{|a_{n}|} \leq{\frac{b_{n+1}}{b_{n}}},\quad \forall n\geq{n_{0}},
  $$
  ent\~ao $\sum\limits_{}^{}a_{n}$ converge absolutamente. Em particular, se $\limsup_{n\to\infty}\frac{|a_{n+1}|}{|a_{n}|} = c < 1,$
tamb\'em vale que $\sum\limits_{}^{}a_{n}$ \'e absolutamente convergente.
\end{theorem*}
\begin{proof*}
  De fato, segue que 
    $$
      \frac{|a_{n_{0}+1}|}{|a_{n_{0}}|} \leq{\frac{b_{n_{0}+1}}{b_{n_{0}}}}, \frac{|a_{n_{0}+2}|}{|a_{n_{0}+1}|} \leq{ \frac{b_{n_{0}+2}}{b_{n_{0}+1}}}, \cdots
    $$
    Logo, $\frac{|a_{n_{0}+p}|}{|a_{n_{0}}|}\leq{\frac{b_{n_{0}+p}}{b_{n_{0}}}}$ e o resultado segue usando o Teorema da Compara\c c\~ao. Por
    fim, o caso particular segue tomando $b_{n} = c^{n}$. \qedsymbol
\end{proof*}
\begin{example}
  Provemos que $\sum\limits_{}^{}\frac{n!}{n^{n}}a^{n}$ \'e convergente para $|a|<e.$ Com efeito, note que, para a n\~ao-nulo, 
    $$
      \frac{\biggl|\frac{(n+1)!}{(n+1)^{(n+1)}}a^{(n+1)}\biggr|}{|\frac{n!}{n^{n}}a^{n}|} = \frac{1}{(1+\frac{1}{n})^{n}}|a|\overbracket[0pt]{\longrightarrow}^{n\to \infty}\frac{|a|}{e}.
    $$ 
    O resultado agora segue do Teste da Raz\~ao.
\end{example}
\begin{example}
  Considere a s\'erie $\sum\limits_{}^{}a_{n}$ com $a_{2n}=2a^{n-1}, a_{2n-1} = a^{2(n-1)}.$ Vamos aplicar o crit\'erio da ra\'iz
e o crit\'erio da raz\~ao para esta s\'erie. Caso n = 2k, temos 
  $$
    \frac{|a_{2k+1}|}{|a_{2k}|} = \frac{|a^{2k}|}{2|a^{2k-1}|} = \frac{|a|}{2}.
  $$
  Agora, se n = 2k-1, 
    $$
      \frac{|a_{2k}|}{|a_{2k-1}|} = \frac{2|a^{2k-1}|}{|a^{2(k-1)}|} = 2|a|.
    $$
    Segue que $\limsup_{n\to\infty}\frac{|a_{n+1}|}{|a_{n}|} = 2|a|.$ Por outro lado, $\sqrt[n]{a_{n}}\overbracket[0pt]{\longrightarrow}^{n\to \infty}|a|.$
Desta forma, o teste da ra\'iz nos d\'a converg\^encia para $|a|<1$, enquanto que o teste da raz\~ao nos d\'a converg\^encia apenas
para $|a|<\frac{1}{2}.$
\end{example}
O exemplo anterior indica uma diferen\c ca no n\'ivel de efic\'acia dos testes. De fato, o Teste da Ra\'iz \'e mais eficiente do que 
o teste da raz\~ao, como segue 
\begin{theorem*}
  Seja $\{a_{n}\}$ uma sequ\^encia limitada de n\'umeros reais n\~ao-nulos. Ent\~ao, 
    $$
      \liminf_{n\to\infty}\frac{|a_{n+1}|}{|a_{n}|}\leq{\liminf_{n\to\infty}}|a_{n}|^{\frac{1}{n}}\leq{\limsup_{n\to\infty}|a_{n}|^{\frac{1}{n}}}\leq{\limsup_{n\to\infty}\frac{|a_{n+1}|}{|a_{n}|}}.
    $$
\end{theorem*}
\begin{proof*}
  Mostremos primeiramente que $\limsup_{n\to\infty}|a_{n}|^{\frac{1}{n}}\leq{\limsup_{n\to\infty}\frac{|a_{n+1}|}{|a_{n}|}}.$
Se n\~ao, seja $c>0$ com $\limsup_{n\to\infty}\frac{|a_{n+1}|}{|a_{n}|}<c<\limsup_{n\to\infty}|a_{n}|^{\frac{1}{n}}.$ Logo,
existe um natural N tal que 
  $$
    \frac{|a_{n+1}|}{|a_{n}|}<c,\quad \forall n\geq{N}.
  $$
  Disto segue que $|a_{N+p}|<|a_{N}|c^{-N}c^{N+p}$ para todo p natural n\~ao-nulo. Sendo assim, $c < \limsup_{n\to\infty}|a_{n}|^{\frac{1}{n}}\leq{c}$,
uma contradi\c c\~ao. 

  Para ver que $\liminf_{n\to\infty}\frac{|a_{n+1}|}{|a_{n}|}\leq{\liminf_{n\to\infty}|a_{n}|^{\frac{1}{n}}},$ procedemos de modo similar.
Comece supondo que existe um $c > 0$ tal que $\liminf_{n\to\infty}\frac{|a_{n+1}|}{|a_{n}|}> c > \liminf_{n\to\infty}|a_{n}|^{\frac{1}{n}}.$
Logo, existe N natural tal que $\frac{|a_{n+1}|}{|a_{n}|}>c$ para todo $n\geq{N},$ e $|a_{N+p}|>|a_{N}|c^{-N}c^{N+p}$ para todo
 $p\in \mathbb{N}^{*}.$ Portanto, temos uma contradi\c c\~ao, pois $c > \liminf_{n\to\infty}|a_{n}|^{\frac{1}{n}} \geq{c}.$ \qedsymbol
\end{proof*}
\begin{crl*}
  Seja $\{a_{n}\}$ uma sequ\^encia limitada de n\'umeros reais n\~ao-nulos. Se existe o limite $\lim_{\to}\frac{|a_{n+1}|}{|a_{n}|}$,
ent\~ao o limite $\lim_{\to}|a_{n}|^{\frac{1}{n}}$ tamb\'em existe e ambos t\^em o mesmo valor.
\end{crl*}
\begin{example}
  Vamos mostrar que 
    $$
      \lim_{\to}\frac{n}{(n!)^{\frac{1}{n}}}=e.
    $$
  De fato, seja $a_{n}=\frac{n^{n}}{n!}$ e note que $(a_{n})^{\frac{1}{n}} = \frac{n}{(n!)^{\frac{1}{n}}}.$ Note tamb\'em que 
    $$
      \frac{a_{n+1}}{a_{n}} = \frac{\frac{(n+1)^{n+1}}{(n+1)!}}{\frac{n^{n}}{n!}} = \frac{(n+1)^{n}}{n^{n}} = (1+\frac{1}{n})^{n}\overbracket[0pt]{\longrightarrow}^{n\to \infty}e.
    $$
\end{example}
  A seguir, apresentaremos tr\^es resultados, o Teorema de Dirichlet sobre o produto de s\'eries, o de Abel sobre o produto de
uma s\'erie por uma sequ\^encia n\~ao-crescente, e o de Leibniz.
\begin{theorem*}
  Seja $\sum\limits_{}^{}a_{n}$ uma s\'erie e $s_{n} = a_{1} + \cdots + a_{n},$ sendo n um natural, as suas somas parciais.
  Se $\{s_{n}\}$ \'e limitada e $\{b_{n}\}$ \'e uma sequ\^encia de n\'umeros reais positivos que \'e n\~ao-crescente e infinit\'esima,
ent\~ao $\sum\limits_{}^{}a_{n}b_{n}$ \'e convergente.
\end{theorem*}
\begin{proof*}
  Segue por indu\c c\~ao que, se $s_{n} = a_{1}+\cdots+a_{n},$
  \begin{align*}
    \sum\limits_{k=1}^{n}a_{k}b_{k} &= a_{1}b_{1} + a_{2}b_{2}+\cdots + a_{n}b_{n}\\
                                    &= s_{1}(b_{1}-b_{2})+s_{2}(b_{2}-b_{3})+\cdots+s_{n-1}(b_{n-1}-b_{n})+s_{n}b_{n}\\
                                    &=\sum\limits_{k=1}^{n-1}s_{k}(b_{k}-b_{k+1}) + s_{n}b_{n}.
  \end{align*}
  Seja $M=\sup{\{|s_{n}|:n\in \mathbb{N}\}}.$ Como $\sum\limits_{}^{}(b_{n}-b_{n+1})$ \'e uma s\'erie convergente de n\'umero reais n\~ao-negativos
e $s_{n}b_{n}\overbracket[0pt]{\longrightarrow}^{n\to \infty}0,$ temos 
  $$
    |s_{k}(b_{k}-b_{k+1})| \leq{M(b_{k}-b_{k+1})},
  $$
  segue que $\sum\limits_{n=1}^{\infty}s_{n}(b_{n}-b_{n+1})$ \'e convergente e que $\sum\limits_{}^{}a_{n}b_{n}$ \'e convergente. \qedsymbol
\end{proof*}
\begin{example}
  Para cada n\'umero real x que n\~ao \'e m\'ultiplo inteiro de $2\pi$, as s\'eries $\sum\limits_{}^{}\frac{\cos{(nx)}}{n}$ e $\sum\limits_{}^{}\frac{\sin{(nx)}}{n}$
  s\~ao convergentes. De fato, vamos utilizar um resultado dos n\'umeros complexos, que afirma que 
    $$
    e^{i\theta} = \cos{(\theta)} + i\sin{(\theta)},\quad i\text{ \'e tal que} i^{2} = -1.
    $$
    Sendo assim, para ver que $\{\sum\limits_{k=1}^{n}\cos{(nx)}\}(\text{ ou }\{\sum\limits_{k=1}^{n}\sin{(nx)}\}$ \'e limitada
  utilizamos que 
    $$
      1 + e^{ix} + e^{i2x} + \cdots + e^{inx} = \frac{1 - e^{i(n+1)x}}{1-e^{ix}},
    $$
    e tomamos parte real e a parte imagin\'aria. O resultado agora segue do Teorema de Dirichlet. \qedsymbol
\end{example}
\begin{theorem*}
  Seja $\sum\limits_{}^{}a_{n}$ uma s\'erie convergente e $\{b_{n}\}$ uma sequ\^encia n\~ao crescente de n\'umeros positivos
(n\~ao necessariamente infinit\'esima). Ent\~ao, a s\'erie $\sum\limits_{}^{}a_{n}b_{n}$ \'e convergente.
\end{theorem*}
\begin{proof*}
  Seja $c = \lim_{n\to\infty}b_{n} = \inf_{n\in \mathbb{N}}b_{n}$ e $s_{n} = a_{1} + \cdots + a_{n}.$ Note que 
    $$
      \sum\limits_{k=1}^{n}a_{k}b_{k} = \sum\limits_{k=1}^{n}a_{k}(b_{k}-c) + c \sum\limits_{k=1}^{n}a_{k}.
    $$
  Agora, pelo Teorema de Dirichlet, $\sum\limits_{n=1}^{\infty}a_{n}(b_{n}-c)$ \'e convergente com soma s. Portanto, 
    $$
    \sum\limits_{}^{}a_{n}b_{n} = s + c \sum\limits_{}^{}a_{n}.\quad\text{\qedsymbol}
    $$
\end{proof*}
\begin{theorem*}
  Seja $\{b_{n}\}$ uma sequ\^encia n\~ao crescente e infinit\'esima. Ent\~ao, a s\'erie $\sum\limits_{}^{}(-1)^{n}b_{n}$ \'e convergente.
\end{theorem*}
\begin{proof*}
  Observe que, no caso em que $a_{n} = (-1)^{n}$ e $s_{n} = a_{1} + \cdots + a_{n},$ ent\~ao $\{s_{n}\}$ \'e limitada. Com isso,
  utilizando o Teorema de Dirichlet, $\sum\limits_{}^{}(-1)^{n}b_{n}$ \'e, portanto, convergente. \qedsymbol
\end{proof*}
\newpage

\section{Aula 12 - 14/04/2023}
\subsection{Motiva\c c\~oes}
\begin{itemize}
  \item Entender s\'eries de pot\^encias;
  \item Estudar s\'eries rearranjadas.
\end{itemize}

\subsection{Testes da Raz\~ao e Ra\'iz Modificados.}
  Come\c camos, antes, com outros exemplos com rela\c c\~ao \`a \'ultima aula. O primeiro \'e 
 \begin{theorem*}
   Seja $\{a_{n}\}$ uma sequ\^encia n\~ao-crescente de n\'umeros reais n\~ao-negativos. A S\'erie $\sum\limits_{}^{}a_{n}$ converge se, e somente se,
  a s\'erie $\sum\limits_{}^{}2^{k}a_{2^{k}}$ \'e convergente.
 \end{theorem*}
 \begin{proof*}
   Sejam $\{s_{n}\}, \{s_{n}^{*}\}$ as sequ\^encias das somas parciais de $\sum\limits_{}^{}a_{n}$ e $\sum\limits_{}^{}2^{k}a_{2^{k}}.$
   Ent\~ao, para todo n natural, 
   \begin{align*}
     s_{n} = a_{1} + a_{2} + \cdots + a_{n} \leq{}s_{2^{n}-1} &= a_{1} + (a_2 + a_{3}) + (a_{4}+a_{5}+a_{6}+a_{7})+\cdots+(a_{2^{n-1}}+\cdots+a_{2^{n}-1})
                                                              &\leq{s_{n-1}^{*}}\leq{s_{n}^{*}.}
   \end{align*}
   Logo, se $\{s_{n}^{*}\}$ \'e limitada, segue que $\{s_{n}\}$ \'e limitada.

   Agora, note que, para todo n natural, 
  \begin{align*}
    s_{2^{n}}=a_{1} + a_{2} + \cdots + a_{2^{n}} &\geq{} \frac{a_{1}}{2} + a_{2} + (a_{3}+a_{4}) + (a_{5}+a_{6}+a_{7}+a_{8})+\cdots+(a_{2^{n-1}+1}+\cdots+a_{2^{n}})\\
                                                 &\geq{}\frac{a_{1}}{2}+a_{2}+2a_{4} + 4a_{8}+\cdots+2^{n-1}a_{2^{n}}=\frac{1}{2}s_{n}^{*}.
  \end{align*}
  Portanto, se $\{s_{n}\}$ \'e limitada, segue que $\{s_{n}^{*}\}$ \'e limitada. \qedsymbol
 \end{proof*}
 \begin{example}
   Do resultado anterior, a s\'erie $\sum\limits_{}^{}\frac{1}{n^{p}}$ \'e convergente se, e s\'o se, a s\'erie $\sum\limits_{}^{}\frac{2^{n}}{2^{np}} = \sum\limits_{}^{}2^{(1-p)n}$ \'e
se, somente se, $p>1.$
 \end{example}
 \begin{example}
   A s\'erie $\sum\limits_{n=2}^{\infty}\frac{1}{n(\log{(n)})}^{p}$ \'e convergente se, e somente se, $p>1.$ Com efeito, segue do resultado anterior que 
     $$
     \sum\limits_{n=1}^{\infty}\frac{2^{n}}{2^{n}(\log{(2^{n})})^{p}} = \sum\limits_{n=1}^{\infty}\frac{1}{\log{2}n^{p}}
     $$
     \'e convergente se, e somente se, $p>1$.
 \end{example}
 \begin{example}
   Vamos provar que e n\~ao \'e racional. De fato, seja $s_{n} = 1 + 1 + \frac{1}{2!} + \frac{1}{3!} + \cdots + \frac{1}{n!}.$ Assim, 
  \begin{align*}
    0 < e - s_{n} &= \frac{1}{(n+1)!}+\frac{1}{(n+2)!} + \frac{1}{(n+3)!} + \cdots \\
                  &< \frac{1}{(n+1)!}[1+\frac{1}{n+1}+\frac{1}{(n+1)^{2}}+ \frac{1}{(n+1)^{3}}+\cdots]
                  &= \frac{1}{(n+1)!}\frac{1}{1-\frac{1}{n+1}} = \frac{1}{nn!}.
  \end{align*}
  Agora, suponha que existem inteiros postivios p e q tais que $e = \frac{p}{q}$. Temos 
    $$
      0 < q!(e-s_{q})<\frac{1}{q}\leq{1}.
    $$
  Por hip\'otese, $q!e$ \'e um inteiro e, como $q!s_{q}$ tamb\'em \'e inteiro, segue que $q!(e-s_{q})$ \'e inteiro em $(0,1)$ e temos
uma contradi\c c\~ao.
 \end{example}
 O exemplo a seguir ilustra um caso em que o teste da raz\~ao n\~ao se aplica, mas o teste da ra\'iz indica converg\^encia.
 \begin{example}
   Considere a s\'erie $\{a_{n}\}$ com $a_{2n} = \frac{1}{2^{n}}$ e $a_{2n-1}= \frac{1}{3^{n}}.$ Note que 
     $$
     \frac{a_{2n}}{a_{2n-1}} = \biggl(\frac{3}{2}\biggr)^{n}\overbracket[0pt]{\longrightarrow}^{n\to \infty}+\infty\quad\text{e}\quad \frac{a_{2n+1}}{a_{2n}} = \frac{1}{3}\biggl(\frac{2}{3}\biggr)^{n}\overbracket[0pt]{\longrightarrow}^{n\to \infty}0.
     $$
  No entanto, fazendo o teste da ra\'iz, 
    $$
    \sqrt[2n]{a_{2n}} = \frac{1}{\sqrt{2}}\overbracket[0pt]{\longrightarrow}^{n\to \infty}\frac{1}{\sqrt[]{2}}\quad\text{e}\quad \sqrt[2n-1]{a_{2n-1}}=\biggl(\frac{1}{3}\biggr)^{\frac{n}{2n-1}}\overbracket[0pt]{\longrightarrow}^{n\to \infty}\frac{1}{\sqrt[]{3}}.
    $$
  Portanto, $\limsup_{n\to\infty}\frac{a_{n+1}}{a_{n}} = +\infty$ e $\limsup_{n\to\infty}\sqrt[n]{a_{n}} = \frac{1}{\sqrt[]{2}}$. Em outras palavras,
o teste da ra\'iz se indica converg\^encia, mas o teste da raz\~ao \'e inconclusivo. \qedsymbol
 \end{example}
  Agora, vamos tornar o teste da ra\'iz mais completo.
  \begin{theorem*}
   Se $\{a_{n}\}$ \'e uma sequ\^encia limitada e $\limsup_{n\to\infty}|a_{n}|^{\frac{1}{n}} = c < 1$, ent\~ao $\sum\limits_{}^{}a_{n}$
  \'e absolutamente convergente. Se $\limsup_{n\to\infty}|a_{n}|^{\frac{1}{n}} = c > 1$, ent\~ao $\sum\limits_{}^{}a_{n}$ \'e divergente. 
  Se c = 1, nada podemos concluir.
 \end{theorem*}
\begin{proof*}
  Existe N natural tal que $\sup_{k\geq{n}}|a_{k}|^{\frac{1}{k}} < r = \frac{c+1}{2} < 1$ para todo natural n. Logo, $|a_{n}|<r^{n}$ para
todo $n\geq{N}.$ Segue do Teorema da Compara\c c\~ao que $\sum\limits_{}^{}|a_{n}|$ \'e convergente, ou seja, $\sum\limits_{}^{}a_{n}$ \'e absolutamente convergente. \qedsymbol
  
  Para mostrar que, se $\limsup_{n\to\infty}|a_{n}|^{\frac{1}{n}} = c > 1$, a s\'erie diverge, considere a fun\c c\~ao $\varphi:\mathbb{N}\rightarrow \mathbb{N}$
  estritamente crescente tal que $\{|a_{\varphi(n)}|^{\frac{1}{\varphi(n)}}\}$ converge para $c>1,$ tal que $\{|a_{\varphi(n)}|\}$
  n\~ao converge para zero. Para ver que nada pode ser dito quando c = 1, tome as s\'eries $\sum\limits_{}^{}\frac{1}{n}, \sum\limits_{}^{}\frac{1}{n^{2}}.$
\end{proof*}
\begin{theorem*}
  Se $\sum\limits_{}^{}a_{n}$ \'e uma s\'erie de termos n\~ao-nulos e $\limsup_{n\to\infty}\frac{|a_{n+1}|}{|a_{n}|} = c < 1,$ ent\~ao
 $\sum\limits_{}^{}a_{n}$ \'e absolutamente convergente. Por\'em, se $\frac{|a_{n+1}|}{|a_{n}|}\geq{1}$ para todo $n\geq{n_{0}}, n_{0}$ algum natural,
 ent\~ao a s\'erie diverge.
\end{theorem*}

\subsection{S\'eries de Pot\^encias.}
\begin{def*}
  Dada uma sequ\^encia $\{a_{n}\}$ de n\'umeros reais, a s\'erie 
    $$
      \sum\limits_{n=0}^{\infty}a_{n}x^{n}
    $$
  \'e chamada de S\'erie de Pot\^encia. Os n\'umeros $a_{n}$ s\~ao chamados de coeficientes da s\'eries e x \'e um n\'umero real.
\end{def*}
  Dependendo da escolha de x, a s\'erie pode convergir ou divergir, como indica o reusltado a seguir.
 \begin{theorem*}
   Ddada a s\'erie de pot\^encias $\sum\limits_{n=0}^{\infty}a_{n}x^{n},$ seja $\alpha = \limsup_{n\to\infty}\sqrt[n]{|a_{n}|}$ e defina 
  \begin{align*}
    & R = \frac{1}{\alpha}\text{ se } 0 < \alpha < \infty\\
    & R = 0 \text{ se } \alpha = \infty\\
    & R = \infty \text{ se } \alpha = 0.
  \end{align*}
 \end{theorem*}
 Ent\~ao, $\sum\limits_{n-0}^{\infty}a_{n}x^{n}$ converge se $|x|<R$, diverge se $|x|>R$ e nada podemos afirmar de $|x|=R.$
\begin{proof*}
  Basta notar que $\limsup_{n\to\infty}\sqrt[n]{|a_{n}x^{n}|} = |x|\limsup_{n\to\infty}\sqrt[n]{|a_{n}|} = |x|\alpha$ e aplicar
o teste da ra\'iz. \qedsymbol
\end{proof*}
 Vamos analisar a converg\^encia das s\'eries de pot\^encias.
\begin{example}
 \begin{itemize}
   \item $\sum\limits_{}^{}n^{n}x^{n}, R = 0. $
   \item $\sum\limits_{}^{}\frac{n^{n}}{n!}x^{n}, R = e^{-1}.$ Com efeito, segue de 
     $$
        \sqrt[n]{\frac{n^{n}}{n!}} = \frac{n}{\sqrt[]{n!}}\overbracket[0pt]{\longrightarrow}^{n\to \infty}e.
     $$
    \item $\sum\limits_{}^{}\frac{x^{n}}{n!}, R = \infty.$
    \item $\sum\limits_{}^{}x^{n}, R=1.$
    \item $\sum\limits_{}^{}\frac{x^{n}}{n^{p}}, p > 0, R = 1.$
 \end{itemize}
\end{example}

\subsection{S\'eries Rearranjadas}
  Seja $\sum\limits_{}^{}a_{n}$ uma s\'erie. Defina as sequ\^encias 
 \begin{itemize}
   \item[+)] $\{a_{n}^{+}\}$ com $a_{n}^{+} = a_{n}$ se $a_{n} > 0$ e $a_{n}^{+}=0$ se $a_{n}\leq{0}.$
   \item[-)] $\{a_{n}^{-}\}$ com $a_{n}^{-} = -a_{n}$ se $a_{n} < 0$ e $a_{n}^{-}=0$ se $a_{n}\geq{0}.$
 \end{itemize}
 As sequ\^encias $\{a_{n}^{+}\}$ e $\{a_{n}^{-}\}$ ser\~ao chamadas de parte positiva e parte negativa de $\{a_{n}\}$. Sendo assim,
 $|a_{n}|=a_{n}^{+}+a_{n}^{-}, a_{n} = a_{n}^{+}-a_{n}^{-}$ e $|a_{n}| = a_{n} + 2a_{n}^{-}.$

 Note que, se $\sum\limits_{}^{}a_{n}$ converge absolutamente, ent\~ao $\sum\limits_{}^{}a_{n}^{+}$ e $\sum\limits_{}^{}a_{n}^{-}$
 s\~ao convergentes. A rec\'iproca tamb\'em vale.

 Al\'em disso, se $\sum\limits_{}^{}a_{n}$ \'e convergente, mas n\~ao absolutamente, segue que ambas $\sum\limits_{}^{}a_{n}^{+}$ e $\sum\limits_{}^{}a_{n}^{-}$ divergem.
\begin{def*}
  Seja $\{a_{n}\}$ a sequ\^encia dos termos da s\'erie $\sum\limits_{}^{}a_{n}, \xi:\mathbb{N}\rightarrow \mathbb{N}$
uma bije\c c\~ao e $b_{n} = a_{\xi(n)}.$ A s\'erie $\sum\limits_{}^{}b_{n}$ \'e chamada uma s\'erie rearranjada de $\sum\limits_{}^{}a_{n}.$
\end{def*}
\begin{example}
  Considere a s\'erie $\sum\limits_{n=1}^{\infty}\frac{(-1)^{n+1}}{n}$. Mostraremos que esta s\'erie converge. Se s \'e a sua soma, temos 
 \begin{align*}
   &s = \sum\limits_{n=1}^{\infty}\frac{(-1)^{n+1}}{n} = 1 - \frac{1}{2} + \frac{1}{3} - \frac{1}{4} + \frac{1}{5} - \frac{1}{6} + \frac{1}{7} -\cdots\\
   &\frac{1}{2}s = \sum\limits_{n=1}^{\infty}\frac{(-1)^{n+1}}{2n} = 0 + \frac{1}{2} + 0 - \frac{1}{4} + 0 + \frac{1}{6} + 0 - \frac{1}{8} + 0 + \cdots\\
   &\frac{3}{2}s = \sum\limits_{n=1}^{\infty}\frac{(-1)^{n+1}}{n} = 1 + 0 + \frac{1}{3} - \frac{1}{2} + \frac{1}{5} + 0 + \frac{1}{7}-\frac{1}{4}+\frac{1}{9}+0+\cdots
 \end{align*}
 Logo, uma s\'erie rearranjada pode ter soma distinta da s\'erie original. 
\end{example}
\begin{theorem*}
  Toda s\'erie rearranjada de uma s\'erie absolutamente convergente \'e convergente com mesma soma.
\end{theorem*}
\begin{proof*}
  Se $a_{n}\geq{}0$, ent\~ao para todo n natural, $\xi:\mathbb{N}\rightarrow \mathbb{N}$ \'e uma bije\c c\~ao e $b_{n} = a_{\xi(n)}$.
  Dado n natural, seja $m_{n} = \max{\{\xi(1), \cdots, \xi(n)\}}$, ent\~ao 
    $$
      \sum\limits_{k=1}^{n}b_{k} \leq{\sum\limits_{k=1}^{m_{n}}a_{k}\leq{\sum\limits_{n=1}^{\infty}a_{n}}}
    $$
  e $\sum\limits_{}^{}b_{n}$ \'e convergente com $\sum\limits_{n=1}^{\infty}b_{n}\leq{\sum\limits_{n=1}^{\infty}a_{n}}$. Por outro lado,
  dado um natural m, seja $n_{m}=\max{\{\xi^{-1}(1), \cdots, \xi^{-1}(m)\}}$. Sendo assim, 
    $$
      \sum\limits_{k=1}^{m}a_{k}\leq{\sum\limits_{k=1}^{n_{m}}b_{k}\leq{\sum\limits_{n=1}^{\infty}b_{k}}}
    $$
  e $\sum\limits_{n=1}^{\infty}a_{n}\leq{\sum\limits_{n=1}^{\infty}b_{n}.}$ Para o caso geral, note que 
    $$
      \sum\limits_{n=1}^{\infty}a_{n} = \sum\limits_{n=1}^{\infty}a_{n}^{+} - \sum\limits_{n=1}^{\infty}a_{n}^{-}
    $$
  de forma que 
    $$
      \sum\limits_{n=1}^{\infty}a_{n} = \sum\limits_{n=1}^{\infty}a_{n}^{+} - \sum\limits_{n=1}^{\infty}a_{n}^{-} = \sum\limits_{n=1}^{\infty}b_{n}^{+}-\sum\limits_{n=1}^{\infty}b_{n}^{-} = \sum\limits_{n=1}^{\infty}b_{n}.
    $$
  Portanto, a s\'erie rearranjada converge para a soma da s\'erie original. \qedsymbol
\end{proof*}
\begin{theorem*}
  Se $\sum\limits_{}^{}a_{n}$ \'e convergente e n\~ao \'e absolutamente convergente, ent\~ao  
 \begin{itemize}
   \item[i)] Dado c real, existe bije\c c\~ao $\xi^{c}:\mathbb{N}\rightarrow \mathbb{N}$ tal que $\sum\limits_{n=1}^{\infty}a_{\xi^{c}(n)} = c$
     \item[ii)] Existem bije\c c\~oes $\xi_{+}$ e $\xi_{-}$ tais que $\sum_{n=1}^{\infty}a_{\xi_{+}(n)}$ diverge para $+\infty$
  e $\sum\limits_{n=1}^{\infty}a_{\xi_{-}(n)}$ diverge para $-\infty.$
 \end{itemize}
\end{theorem*}
\begin{proof*}
  Seja $\{p_{n}\}$ a sequ\^encia dos termos positivos de $\{a_{n}\}$ na ordem em que eles aparecem e $\{q_{n}\}$ a sequ\^encia dos termos 
  n\~ao positivos de $\{a_{n}\}$ na ordem em que eles aparecem. Sabemos que $\sum\limits_{}^{}p_{n}$ e $\sum\limits_{}^{}q_{n}$ divergem.
Dado c real, seja $n_{1}$ o primeiro inteiro tal que 
  $$
    \sum\limits_{n=1}^{n_{1}}p_{n} > c.
  $$
  Em seguida escolha $n_{2}$ o menor inteiro tal que 
    $$
      \sum\limits_{n=1}^{n_{1}}p_{n} + \sum\limits_{n=1}^{n_{2}}q_{n} < c
    $$
  e prossiga com este processo. Desta forma, para todo $k > 1,$ 
    $$
      0 < \sum\limits_{n=1}^{n_{1}}p_{n} - \sum\limits_{n=1}^{n_{2}}q_{n} + \cdots - \sum\limits_{n=1}^{n_{2k-2}}q_{n} + \sum\limits_{n=1}^{n_{2k-1}}p_{n} - c \leq{p_{n_{2k-1}}}
    $$
    e 
      $$
        q_{n_{2k}}\leq{\sum\limits_{n-1}^{n_{1}}pn} - \sum\limits_{n=1}^{n_{2}}q_{n} + \cdots + \sum\limits_{n=1}^{n_{2k-1}}p_{n} - \sum\limits_{n=1}^{n_{2k}}q_{n} - c < 0.
      $$
    Agora, como $q_{n}\overbracket[0pt]{\longrightarrow}^{n\to \infty}0$ e $p_{n}\overbracket[0pt]{\longrightarrow}^{n\to \infty}0,$ o resultado segue.
    Por fim, um processo an\'alogo prova a segunta parte do resultado. \qedsymbol
\end{proof*}
\newpage

\section{Aula 13 - 17/04/2023}
\subsection{Motiva\c c\~oes}
\begin{itemize}
  \item Limites e Continuidade de Fun\c c\~oes;
  \item Limites laterais.
\end{itemize}
\subsection{Intui\c c\~ao e Exemplos Iniciais}
Entender continuidade e limites é fundamental no cálculo. Aqui estão alguns exemplos que podem ajudar a esclarecer esses conceitos:

\begin{enumerate}
    \item Continuidade: \\
    Uma função é contínua em um ponto se o seu valor nesse ponto é igual ao limite da função quando se aproxima desse ponto. Em outras palavras, não há quebras ou saltos no gráfico da função. Um exemplo clássico de uma função contínua é a função quadrática:
    
    \begin{equation*}
        f(x) = x^2
    \end{equation*}
    
    Essa função é contínua para todos os números reais, pois não há quebras ou saltos no gráfico.

    \begin{center}
    \begin{tikzpicture}
    \begin{axis}[
        xlabel={$x$},
        ylabel={$f(x)=x^2$},
        domain=-3:3,
        samples=100,
        axis lines=middle
    ]
    \addplot[blue, thick] {x^2};
    \end{axis}
    \end{tikzpicture}
    \end{center}

    \item Limites: \\
    Um limite é o valor que uma função se aproxima quando a entrada (valor x) se aproxima de um determinado valor. O conceito de limites nos ajuda a entender o comportamento das funções próximas aos pontos onde elas não são definidas ou não são contínuas.
    
    Exemplo 1: Limite simples \\
    Considere a função:
    
    \begin{equation*}
        g(x) = 3x + 2
    \end{equation*}
    
    Encontre o limite quando x se aproxima de 1:
    
    \begin{equation*}
        \lim_{x \to 1} (3x + 2)
    \end{equation*}
    
    Como essa é uma função linear e contínua em todos os lugares, o limite é o mesmo que o valor da função no ponto:
    
    \begin{align*}
        g(1) &= 3(1) + 2 \\
        &= 5
    \end{align*}
    
    Portanto, $\lim_{x \to 1} (3x + 2) = 5$.

    \begin{center}
    \begin{tikzpicture}
    \begin{axis}[
        xlabel={$x$},
        ylabel={$g(x)=3x+2$},
        domain=-2:3,
        samples=100,
        axis lines=middle
    ]
    \addplot[red, thick] {3*x+2};
    \end{axis}
    \end{tikzpicture}
    \end{center}
    
    Exemplo 2: Limite de uma função com descontinuidade \\
    Considere a função:
    
    \begin{equation*}
        h(x) = \frac{x^2 - 1}{x - 1}
    \end{equation*}
    
    Essa função não é definida para x = 1, mas ainda podemos encontrar o limite quando x se aproxima de 1:
    
    \begin{equation*}
        \lim_{x \to 1} \frac{x^2 - 1}{x - 1}
    \end{equation*}
    
    Observe que $(x^2 - 1)$ pode ser fatorado como $(x + 1)(x - 1)$. Então, obtemos:
    
    \begin{equation*}
        \lim_{x \to 1} \frac{(x + 1)(x - 1)}{x - 1}
    \end{equation*}
    
    Agora, podemos cancelar os termos $(x - 1)$:
    
    \begin{equation*}
        \lim_{x \to 1} (x + 1)
    \end{equation*}
    
    Neste ponto, podemos substituir x = 1:

    \begin{align*}
        1 + 1 &= 2
    \end{align*}
    
    Portanto, $\lim_{x \to 1} (h(x)) = 2$, embora a função em si não seja definida para x = 1.

    \begin{center}
    \begin{tikzpicture}
    \begin{axis}[
        xlabel={$x$},
        ylabel={$h(x)=\frac{x^2 - 1}{x - 1}$},
        domain=-3:3,
        ymin=-5, ymax=5,
        samples=100,
        axis lines=middle,
        restrict y to domain=-5:5
    ]
    \addplot[green, thick] {((x^2 - 1)/(x - 1))};
    \addplot[only marks, mark=o, mark options={scale=1.5}] coordinates {(1,2)};
    \end{axis}
    \end{tikzpicture}
    \end{center}
\end{enumerate}

Esses exemplos demonstram como a continuidade e os limites nos ajudam a analisar e entender o comportamento das funções, mesmo em pontos onde elas podem não ser definidas. 
A seguir, formalizaremos essas ideias e provaremos algumas propriedades.

\subsection{Limites de Fun\c c\~oes}
 \begin{def*}
   Seja D um subconjunto de $\mathbb{R}, f:D\rightarrow \mathbb{R}$ uma fun\c c\~ao e p um ponto de acumula\c c\~ao de D. Diremos que
  o limite de f(x) quando x tende a p \'e L se, dado $\varepsilon > 0,$ existe um $\delta > 0$ tal que 
    $$
    x\in D\text{ e } 0 < |x-p| < \delta, \Rightarrow |f(x) - L | < \varepsilon.
    $$
    Em outras palavras, dado $\varepsilon > 0$, existe $\delta = \delta(\varepsilon, p) > 0$ tal que 
      $$
        f(D\cap(p-\delta, p+\delta)/a)\subseteq{(L-\varepsilon, L+\varepsilon).} \square
      $$
 \end{def*}
 Note que, se n\~ao existe um n\'umero real L tal que $\lim_{x\to p}= L,$ ent\~ao diremos que o limite n\~oa existe. Al\'em disso,
o ponto p n\~ao precisa ser um ponto do dom\'inio D e, mesmo que perten\c ca, o valor de f em p n\~ao \'e importante para a defini\c c\~ao.
Apenas os valores de f em pontos arbitrariamente pr\'oximos a p s\~ao importantes para a defini\c c\~ao.
\begin{theorem*}
  Seja $f:D\rightarrow \mathbb{R}$ uma fun\c c\~ao e p um ponto de acumula\c c\~ao de D. O limite de f(x) quando x tende a p,
  caso existe, \'e \'unico, e ser\'a dneotado por 
    $$
    \lim_{x\to p}f(x) = L.
    $$
\end{theorem*}
\begin{proof*}
  De fato, se L e L' s\~ao limites de f(x) quando x tende a p, dado $\varepsilon > 0$, existe $\delta > 0$ tal que 
    $$
      x\in D, 0 < |x-p| < \delta \Rightarrow |f(x)- L| < \varepsilon, |f(x) - L'|< \varepsilon.
    $$
    Logo, dado $\varepsilon > 0$, com a escolha de $\delta$ acima e x em D satisfazendo $0 <|x-p|<\delta$, temos 
      $$
        |L - L'| = |L -f(x) + f(x) - L'| \leq{} |L - f(x)| + |f(x) - L'| < 2\varepsilon.
      $$
    Portanto, como $\varepsilon$ \'e arbitr\'ario, L = L'. \qedsymbol
\end{proof*}
  Quando nos referimos a uma fun\c c\~ao, fica impl\'icito que ela tem um dom\'inio especificado. Dada a fun\c c\~ao
  $f:D\rightarrow \mathbb{R}$ e $D'\subseteq{D},$ denotaremos por $f_{|_{D'}}:D'\rightarrow \mathbb{R}$ a fun\c c\~ao
  definida por $f_{|_{D'}}(x) = f(x)$ para x em D'. Segue dessa defini\c c\~ao que 
 \begin{theorem*}
   Seja D um subconjunto dos reais, $f:D\rightarrow \mathbb{R}$ uma fun\c c\~ao, D' um subconjunto de D e p um ponto de aucmula\c c\~ao de D'.
  Se $\lim_{x\to p}f(x) = L,$ ent\~ao $\lim_{x\to p}f_{|_{D'}}(x) = L.$
 \end{theorem*}
\begin{theorem*}
  Seja $f:D\rightarrow \mathbb{R}$ uma fun\c c\~ao, D' e D'' subconjuntos de D e p um n\'umero real que \'e ponto de acumula\c c\~ao
de D' e de D''. 
 \begin{itemize}
   \item[i)] Se um dos limites $\lim_{x\to p}f_{|_{D'}}(x)$ ou $\lim_{x\to p}f_{|_{D''}}(x)$ n\~ao existe, ou ambos existem,
     mas s\~ao diferentes, ent\~ao o limite $\lim_{x\to p}f(x)$ n\~ao existe.
   \item[ii)] Se $(D'\cup D'')/\{p\} = D/\{p\}$, o limite $\lim_{x\to p}f(x)$ existe se, e somente se, $\lim_{x\to p}f_{|_{D'}}(x)$ e
$\lim_{x\to p}f_{|_{D''}}(x)$ existem e t\^em o mesmo valor.
 \end{itemize}
\end{theorem*}
\begin{proof*}
  A prova da primeira parte segue direto. Para a segunda, existe $\lim_{x\to p }f(x) = L$ se, e somente se, dado $\varepsilon > 0$,
existe $\delta > 0$ tal que 
  $$
    x\in D, 0 < |x-p| < \delta \Rightarrow |f(x)-L| < \varepsilon,
  $$
  o que equivale a dizer que, dado $\varepsilon > 0$, existe $\delta > 0$ tal que 
    $$
      x\in D', 0 < |x-p| < \delta \Rightarrow |f(x) - L|< \varepsilon \Rightarrow |f(x)_{|_{D'}}|<\varepsilon
    $$
  e
  $$
      x\in D'', 0 < |x-p| < \delta \Rightarrow |f(x) - L|< \varepsilon \Rightarrow |f(x)_{|_{D''}}|<\varepsilon
  $$
  Portanto, $\lim_{x\to p}f_{|_{D''}}(x) = \lim_{x\to p}f_{|_{D'}}(x)$. \qedsymbol
\end{proof*}
\begin{def*}
  Se D \'e um subconjunto de $\mathbb{R}$, diremos que p real \'e um ponto de acumula\c c\~ao \`a direita de D se \'e um ponto
  de acumula\c c\~ao de $D_{p}^{+} = D\cap{(p, \infty)}$. Seja $f:D\rightarrow \mathbb{R}$ uma fun\c c\~ao e p um ponto de
  acumula\c c\~ao \`a direita de D. O limite de f(x) quando x tende a p pela direita \'e 
    $$
      \lim_{x\to p^{+}}f(x)\coloneqq \lim_{x\to p}f_{|_{D_{p}^{+}}}(x).
    $$
  Define-se analogamente limite de f(x) quando x tende a pela esquerda e ponto de acumula\c c\~ao \`a esquerda.
\end{def*}
\begin{crl*}
  Seja $f:D\rightarrow \mathbb{R}$ uma fun\c c\~ao e p um ponto de acumula\c c\~ao \`a direita e \`a esquerda de D. Ent\~ao, 
    $$
      \lim_{x\to p}f(x)
    $$
    existe se, e somente se, os limites laterais existem e s\~ao iguais.
\end{crl*}
\begin{theorem*}
  Seja D subconjunto real e $f:D\rightarrow \mathbb{R}$ uma fun\c c\~ao. Tome p como um ponto de acumula\c c\~ao de D. Se
  existe $\lim_{x\to p}f(x) - L$, ent\~ao f \'e limitada em uma vizinha\c ca de p, i.e., existem $M>0$ e $\delta > 0$ tais que 
  se x pertence a D e $0 <|x-p| < \delta,$ ent\~ao $|f(x)|< M$
\end{theorem*}
\begin{proof*}
  Existem $\delta > 0$ tal que se x pertence a D e $0<|x-p|<\delta,$ ent\~ao $|f(x)-L|<1.$ Logo, 
    $$
      |f(x)| \leq{|f(x)-L| + |L|} \leq{1 + |L| = M},
    $$
\end{proof*}
\begin{theorem*}
  Seja $D\subseteq{\mathbb{R}}, f, g, h:D\rightarrow \mathbb{R}$ fun\c c\~oes e p um ponto de acumula\c c\~ao de D. Se
  para todo x em D diferente de p, $f(x)\leq{g(x)}\leq{h(x)}$ e $\lim_{x\to p}f(x) = \lim_{x\to p }h(x) = L,$ ent\~ao $\lim_{x\to p}g(x) = L.$
\end{theorem*}
\begin{proof*}
  Com efeito, dado $\varepsilon > 0$, existe $\delta > 0$ tal que se x pertence a D e $0 < |x-p| < \delta,$ ent\~ao
  $|f(x)-L|<\varepsilon, |h(x)-L|<\varepsilon.$ Logo, 
    $$
      L - \varepsilon < f(x) \leq{g(x)} \leq{h(x)} < L +\varepsilon, \quad \forall x\in D, 0<|x-p|<\delta.
    $$
  Segue que $L-\varepsilon < g(x) < L+\varepsilon,$ para todo x em D que satisfa\c ca a condi\c c\~ao. Em outras palavras, 
    $$
    |g(x)-L| < \varepsilon, \quad \forall x\in D, 0<|x-p|<\delta.\text{ \qedsymbol}
    $$
\end{proof*}
\begin{theorem*}
  Seja D subconjunto de $\mathbb{R}, f:D\rightarrow \mathbb{R}$ e p um ponto de acumula\c c\~ao de D. Se $\lim_{x\to p}f(x) = L > 0,$
  ent\~ao existe $\delta > 0$ tal que $f(x) > 0$ para todo x em D com $0<|x-p|<\delta.$
\end{theorem*}
\begin{proof*}
  Dado $\varepsilon = \frac{L}{2}$ tal que 
  $$
    -\frac{L}{2} < f(x) - L < \frac{L}{2}
  $$
  para todo x em D, $0 < |x-p|<\delta.$ Logo, $0 <\frac{L}{2}<f(x)$ para todo x em D com $0 <|x-p|<\delta.$ \qedsymbol
\end{proof*}
\begin{theorem*}
  Seja $D\subseteq{\mathbb{R}}, f, g:D\rightarrow \mathbb{R}<$ uma fun\c c\~ao e p um ponto de acumula\c c\~ao de D. Se existe $\delta > 0$
  tal que $f(x)\leq{g(x)}$ para todo x em D com $0 <|x-p|<\delta$ e existe $\lim_{x\to p}f(x) = L_{g}$, ent\~ao $L_{f}\leq{L_{g}}.$
\end{theorem*}
\begin{proof*}
  De fato, dado $\varepsilon > 0$, existe $\delta > 0$ tal que x pertence a D e satisfaz a propriedade, ent\~ao 
    $$
      L_{f} - \frac{\varepsilon}{2} \leq{f(x)}\leq{g(x)}\leq{L_{g}+\frac{\varepsilon}{2}.}
    $$
    Portanto, $L_{f}-L_{g}\leq{\varepsilon}$ e, como $\varepsilon > 0$ \'e abritr\'ario, o resultado segue. \qedsymbol
\end{proof*}
\begin{theorem*}
  Seja D subconjunto real, $f:D\rightarrow \mathbb{R}$ e p ponto de acumula\c c\~ao. O limite $\lim_{x\to p}f(x)=L$ se, e s\'o se,
  $\lim_{n\to\infty}f(x_{n})$ existe para toda sequ\^encia $\{x_{n}\}$ em $D/\{p\}$ que converge para p.
\end{theorem*}
\begin{proof*}
  Se $\lim_{x\to p }f(x) = L$ e $\{x_{n}\}$ \'e uma sequ\^encia em $D/\{p\}$ com $x_{n}\overbracket[0pt]{\longrightarrow}^{n\to \infty}p$, dado
  $\varepsilon > 0$, podemos encontrar $\delta > 0$ tal que $|f(x)-L|<\varepsilon$ para quaisquer $x\in D, 0<|x-p|<\delta.$

  Seja N natural tal que $|x_{n}-p|<\delta$ para todo $n\geq{N}.$ Logo, $|f(x_{n})-L| < \varepsilon$ para todo $n\geq{N}.$
  Portanto, $\lim_{n\to\infty}f(x_{n})=L.$ 

  Por outro lado, note que, se $\lim_{n\to\infty}f(x_{n})$ existe para toda sequ\^encia $\{x_{n}\}$ em $D/\{p\}$ que converge
  para p, todas as sequ\^encias $\{f(x_{n})\}$ t\^em o mesmo limite, visto que se elas n\~ao tivessem, construir\'iamos uma sequ\^encia
  $\{x_{n}'\}$ em $D/\{p\}$ que converge para p, mas $\{f(x_{n}')\}$ n\~ao convergiria.

  Agora, se $\lim_{x\to p}f(x)$ n\~ao \'e L, existe $\varepsilon > 0$ tal que para todo n natural n\~ao nulo, $x_{n}\in D, 0 <|x_{n}-p|<\frac{1}{n}$ tal 
  que $|f(x_{n})-L|\geq{\varepsilon}$. Logo, $\lim_{n\to\infty}f(x_{n})$ n\~ao \'e L.\qedsymbol
\end{proof*}
\begin{theorem*}
  Seja $D\subseteq{\mathbb{R}}, f, g:D\rightarrow \mathbb{R}$ fun\c c\~oes e p um ponto de acumula\c c\~ao de D e $\lambda\in \mathbb{R}.$
 \begin{itemize}
   \item[i)] Se existem M e $\delta$ positivos tais que $|f(x)|\leq{M}$ para todo x em D, $0 <|x-p|<\delta$ e $\lim_{x\to p}g(x)=0,$
     ent\~ao $\lim_{x\to p}(f \cdot g)(x) = 0.$
     \item[ii)] Se $\lim_{x\to p}f(x) = L_{f}$ e $\lim_{x\to p}g(x) = L_{g},$ ent\~ao $\lim_{x\to p}(f + \lambda g)(x) = L_{f} + \lambda L_{g}.$
     e $\lim_{x\to p}(f \cdot g)(x) = L_{f}L_{g},$. Al\'em disso, se $L_{g}\neq0,$ ent\~ao $\lim_{x\to p}\frac{f}{g}(x) = \frac{L_{f}}{L_{g}}.$
 \end{itemize}
\end{theorem*}
\begin{theorem*}
  Seja $D\subseteq{\mathbb{R}}, f, g:D\rightarrow \mathbb{R}$ fun\c c\~oes e p um ponto de acumula\c c\~ao de D. O limite
  $\lim_{x\to p}f(x)$ existe se, e somente se, f \'e de Cauchy em p, i.e., dado $\varepsilon > 0$, existe $\delta > 0$ tal que
  para x, y em D, $0<|x-p|<\delta$ e $0<|y-p|<\delta$ implica em $|f(x)-f(y)|<\varepsilon.$
\end{theorem*}
\begin{proof*}
  \'E claro que se $\lim_{x\to p}f(x) = L$, ent\~ao f \'e de Cauchy em p. De fato, dado $\varepsilon > 0$, existe $\delta > 0$ 
  tal que se x \'e um elemento de D para o qual $0<|x-p|<\delta,$ ent\~ao $|f(x)-L|<\frac{\varepsilon}{2}.$ Assim, dados
  x, y em D satisfazendo $0<|x-p|<\delta, 0<|y-p|<\delta,$ ent\~ao 
    $$
      |f(x)-f(y)| = |f(x)-L+L-f(y)|\leq{|f(x)-L| + |f(y)-L|}< \frac{\varepsilon}{2}+\frac{\varepsilon}{2} = \varepsilon.
    $$

  Reciprocamente, se f \'e de Cauchy em p e $\{x_{n}\}$ \'e uma sequ\^encia em $D/\{p\}$ que converge para p, 
  $\{f(x_{n})\}$ \'e de Cauchy e portanto convergente. Detalhando, dado $\varepsilon > 0$, existe $\delta > 0$ tal que 
  se x, y s\~ao elementos de D, $0 < |x-p| <\delta$ e $0 < |x-p| <\delta$ implicar em $|f(x)-f(y)|<\varepsilon$,
seja $\{x_{n}\}$ uma sequ\^encia em $D/\{p\}, x_{n}\overbracket[0pt]{\longrightarrow}^{n\to \infty}p$. Dado $\varepsilon > 0$
tome $\delta$ da defini\c c\~ao de ``f \'e de Cauchy em p'' e N tal que $0<|x_{n}-p|<\delta$ para todo $n\geq{N}.$ Ent\~ao, 
  $$
  |f(x_{n})-f(x_{m})| < \varepsilon\quad \forall n, m\geq{N}.
  $$
  Portanto, $\{f(x_{n})\}$ converge. \qedsymbol
\end{proof*}
\begin{def*}
  Seja D um subconjunto ilimitado superiormente de $\mathbb{R}$ e $f:D\rightarrow \mathbb{R}$ uma fun\c c\~ao. Diremos que 
  o limite de f(x) quando x tende para infinito \'e L em $\mathbb{R}$ se, dado $\varepsilon > 0$, existe $M = M(\varepsilon) > 0$
  tal que 
    $$
      x\in D, x > M \Rightarrow |f(x)-L|<\varepsilon.
    $$
  Escreveremos 
    $$
      \lim_{x\to\infty}f(x) = L.
    $$
  Analogamente, quando D \'e ilimitado inferiormente, definimos 
    $$
      \lim_{x\to-\infty}f(x) = L.\square
    $$
\end{def*}
  O limite da sequ\^encia \'e um caso particular de limite infinito, especificamente o caso $D = \mathbb{N}.$
 \begin{def*}
   Seja D um subconjunto de $\mathbb{R}$ e $f:D\rightarrow \mathbb{R}$ uma fun\c c\~ao. Se p \'e um ponto de acumula\c c\~ao de D,
   diremos que f diverge para $+\infty$ quando x tende para p se, dado M positivo, existe $\varepsilon = \varepsilon(M) > 0$
   tal que 
     $$
       x\in D, 0<|x-p|<\varepsilon \Rightarrow f(x) > M.
     $$
     Escreveremos $\lim_{x\to p}f(x) = +\infty$. Analogamente, define-se $\lim_{x\to p}f(x) =-\infty$. 

     Se D \'e ilimitado superiormente(inferiormente), definimos tamb\'em 
       $$
         \lim_{x\to+\infty}f(x) = \pm\infty\quad(\lim_{x\to-\infty}f(x)=\pm\infty.). \square
       $$
 \end{def*}
\begin{def*}
  Seja D um subconjunto real e $f:D\rightarrow \mathbb{R}$. Se p \'e um ponto de acumula\c c\~ao de D, suponha que existe $\delta_{0} > 0$
  tal que 
    $$
    \sup{\{f(x):x\in D, 0<|x-p|<\delta_{0}\}} < \infty.
    $$
    Ent\~ao, existe (ou diverge para $-\infty$) o limite 
      $$
      \limsup_{x\to p}f(x)\coloneqq \limsup_{\delta\to0}\{f(x):x\in D, 0<|x-p|<\delta\}.
      $$
      Escrevemos $\limsup_{x\to p}f(x) = +\infty$ quando f n\~ao \'e limitada superiormente em nenhuma vizinhan\c ca de p.

      De maneira an\'aloga, se 
        $$
        \inf{\{f(x):x\in D, 0<|x-p|<\delta\}} > -\infty,
        $$
      definimos (podendo ser $+\infty$) 
        $$
        \liminf_{x\to p}f(x)\coloneqq \liminf_{\delta\to0}\{f(x):x\in D, 0<|x-p|<\delta\}.
        $$
      Escrevemos $\liminf_{x\to p}f(x) = -\infty$ quando f n\~ao for limitada inferiormente em uma vizinhan\c ca de p. $\square$
\end{def*}
\newpage

\section{Aula 14 - 19/04/2023}
\subsection{Motiva\c c\~oes}
\begin{itemize}
  \item Valor de ader\^encia - valores para os quais alguma sequ\^encia converge sob a imagem de uma fun\c c\~ao
  \item Continuidade;
  \item Resultados sobre Continuidade.
\end{itemize}
\subsection{Limites Superior e Inferior}
\begin{def*}
  Dizemos que um n\'umero real y \'e um valor de ader\^encia de f no ponto p se existe uma sequ\^encia $\{x_{n}\}$
em $D/\{p\}, x_{n}\overbracket[0pt]{\longrightarrow}^{n\to \infty}p$ e $\lim_{n\to\infty}f(x_{n}) = y.\square$
\end{def*}
\begin{theorem*}
  Seja $D\subseteq{\mathbb{R}}, f:D\rightarrow \mathbb{R}$ uma fun\c c\~ao e p um ponto de acumula\c c\~ao de D.
 \begin{itemize}
   \item[1)] Se l \'e um valor de ader\^encia de f em p, ent\~ao $\liminf_{x\to p}f(x) \leq{l}\leq{\limsup_{x\to p}f(x)}$
    \item[2)] Se f \'e limitada em uma em uma vizinhan\c ca de p, ent\~ao $\limsup_{x\to p}f(x)$ e
    $\liminf_{x\to p}f(x)$ s\~ao valores de ader\^encia de f.
    \item[3)] $\lim_{x\to p}f(x)$ existe se, e somente se, f \'e limitada em uma vizinhan\c ca de p e o conjunto dos
    valores de ader\^encia de f em p \'e unit\'ario.
    \item[4)] Se f \'e limitada em uma vizinhan\c ca de p, dado $\varepsilon > 0$, existe $\delta > 0$ tal que 
      $\liminf_{x\to p}f(x) - \varepsilon < f(x) < \limsup_{x\to p} + \varepsilon$ para todo x em D com $0<|x-p|<\delta.$ 
 \end{itemize}
\end{theorem*}
Antes de prov\'a-lo, observe que, se $L_{\delta} = \sup\{f(x):x\in D, 0 < |x-p| < \delta\},$ ent\~ao 
 $$
    L = \limsup_{x\to p}f(x) = \lim_{\delta\to 0^{+}}L_{\delta}.
 $$
 Assim, dado $\varepsilon > 0$, existe $\delta_{\varepsilon}>0$ tal que $0 < \delta < \delta_{\varepsilon}$ implica 
   $$
     |L_{\delta}-L|< \varepsilon.
   $$
  Isso funciona para provar que um valor \'e o lim sup de uma fun\c c\~ao.
\begin{proof*}
  Prova de 1): Se $\liminf_{x\to p}f(x) = l$ e $\limsup_{x\to p}f(x) = L,$ dado $\varepsilon > 0,$ existe $\delta_{\varepsilon}>0$
  tal que 
    $$
    l - \varepsilon < \inf{\{f(x):x\in D, 0<|x-p|<\delta\}} < l + \varepsilon
    $$
  e 
    $$
      L - \varepsilon < \sup{\{f(x):x\in D, 0<|x-p|<\delta\}} < L + \varepsilon.
    $$
  Para todo $0<\delta<\delta_{\varepsilon}.$ Escolha $\delta_{0} < \delta_{\varepsilon}$. Se r \'e um valor de ader\^encia
  de f em p, existe $x_{n}\in D/\{p\}, x_{n}\overbracket[0pt]{\longrightarrow}^{n\to \infty}p,$ com $f(x_{n})\overbracket[0pt]{\longrightarrow}^{n\to \infty}l$
  Seja N natural tal que $|x_{n}-p| < \delta_{0}$ para todo $n\geq{N}.$ Logo, para todo $n\geq{N},$
 \begin{align*}
   l - \varepsilon &< \inf\{f(x):x\in D, 0 < |x-p| < \delta_{0}\}\\
   &\leq{}f(x_{n})\leq{}\sup\{f(x):x\in D, 0 <|x-p|<\delta_{0}\} < \varepsilon
 \end{align*}
 Segue que $l-\varepsilon\leq{r}\leq{L+\varepsilon}$ para todo $\varepsilon > 0$. Portanto, $l\leq{r}\leq{L}.$

 Prova de 2): Note que, para algum $\delta_{0} > 0$, temos 
   $$
   -\infty < \inf{\{f(x):x\in D, 0 < |x-p|< \delta_{0}\}}\leq{}\sup\{f(x):x\in D, 0 < |x-p| <\delta_{0}\} < \infty.
   $$
   Como 
  \begin{itemize}
    \item $(0, \delta_{0})\ni\delta\mapsto\inf\{f(x):x\in D, 0<|x-p|<\delta\}$ \'e n\~ao-decrescente e 
    \item $(0, \delta_{0})\ni\delta\mapsto\sup\{f(x):x\in D, 0<|x-p|<\delta\}$ \'e n\~ao-crescente,
  \end{itemize}
  existem os limites 
    $$
    \lim_{\delta\to 0^{+}}\inf\{f(x):x\in D, 0 < |x-p| < \delta\} = l, \quad\lim_{\delta\to 0^{+}}\sup\{f(x):x\in D, 0 < |x-p| < \delta\} = L
    $$
    Como p \'e um ponto de acumula\c c\~ao de D, seja $\{x_{n}'\}$ e $\{x_{n}^{L}\}$ sequ\^encias em D tais que $0<\max\{|x_{n}'-p|, |x_{n}^{L}-p|\}<\frac{\delta_{0}}{n}$
    e 
      $$
      \inf{\{f(x):x\in D, 0 < |x-p| < \frac{\delta_{0}}{n}\}}\leq{f(x_{n}')}\leq{}\inf{\{f(x):x \in D, 0 <|x-p|<\frac{\delta_{0}}{n}\}} + \frac{1}{n}
      $$
    e
    $$
      \sup{\{f(x):x\in D, 0 < |x-p| < \frac{\delta_{0}}{n}\}}-\frac{\delta_{0}}{n}\leq{f(x_{n}^{L})}\leq{}\sup{\{f(x):x \in D, 0 <|x-p|<\frac{\delta_{0}}{n}\}} + \frac{1}{n}.
    $$
    O resultado agora segue tomando o limite nas express\~oes acima. 

  Prova de 3): Se o limite existe, f \'e limitada em uma vizinhan\c ca de p e todos os valores de ader\^encia coincidem e,
em particular, o $\limsup_{x\to p}f(x) = \liminf_{x\to p}f(x).$ Por outro lado, se f \'e limitada em uma vizinhan\c ca de um ponto p,
e o conjunto dos valores de ader\^encia \'e unit\'ario, $\liminf_{x\to p}f(x)=\limsup_{x\to p}f(x)$ e todos os valores de ader\^encia
coincidem. Portanto, o limite existe 
  
  Prova de 4): Se $\liminf_{x\to p}f(x) = l$ e $\limsup_{x\to p}f(x) = L,$ dado $\varepsilon > 0,$ existe $\delta_{\varepsilon}>0$
  tal que 
    $$
    l - \varepsilon < \inf{\{f(x):x\in D, 0<|x-p|<\delta\}} < l + \varepsilon
    $$
  e 
    $$
      L - \varepsilon < \sup{\{f(x):x\in D, 0<|x-p|<\delta\}} < L + \varepsilon.
    $$
  Para todo $0<\delta<\delta_{\varepsilon}.$  Logo, para todo $\delta<\delta_{\varepsilon}$ e x em D,$0<|x-p|<\delta$
 \begin{align*}
   l - \varepsilon &< \inf\{f(x):x\in D, 0 < |x-p| < \delta_{0}\}\\
   &\leq{}f(x_{n})\leq{}\sup\{f(x):x\in D, 0 <|x-p|<\delta_{0}\} < \varepsilon
 \end{align*}
 Segue que $l-\varepsilon\leq{r}\leq{L+\varepsilon}$ para todo $\varepsilon > 0$. Portanto, $l\leq{r}\leq{L}.$
\end{proof*}

\subsection{Fun\c c\~oes Cont\'inuas}
 \begin{def*}
   Seja $f:D_{f}\rightarrow \mathbb{R}$ uma fun\c c\~ao e p um ponto de $D_{f}.$ Diremos que $f(x)$ \'e cont\'inua em p
   se, dado $\varepsilon > 0$, existe um $\delta > 0$ tal que 
     $$
       x\in D_{f}, |x-p|<\delta \Rightarrow |f(x)-f(p)|<\varepsilon. 
     $$
    Se isto ocorre para todos os pontos p em $D_{f},$ diremos apenas que f \'e cont\'inua. $\square$
 \end{def*}
 Note que, se p \'e um ponto de $D_{f}$ que \'e de acumula\c c\~ao, ent\~ao f \'e cont\'inua em p se, e somente se,
 $\lim_{x\to p}f(x)=f(p)$ e, se p \'e um ponto isolado de $D_{f},$ ent\~ao f \'e cont\'inua em p.
 \begin{example}
   As seguintes fun\c c\~oes s\~ao cont\'inuas em x=p para todo p real:
  \begin{itemize}
    \item[i)]$f(x) = k$
    \item[ii)]$f(x) = x$
    \item[iii)]$f(x) = x + 1$
    \item[iv)]$f(x) = x^{2}$
  \end{itemize}
 \end{example}
 No entanto, a fun\c c\~ao
 $$
 f(x) = \left\{\begin{array}{ll}
                \frac{x^{2}-1}{x-1},\quad x\neq 1\\
                0,\quad x=1
              \end{array}\right.
 $$
 n\~ao \'e cont\'inua em x=1, pois $\lim_{x\to 1}f(x) = 2\neq 0 = f(1).$
 As fun\c c\~oes cont\'inuas t\^em a seguinte propriedade:
\begin{theorem*}
  Sejam $f_{i}:d_{i}\rightarrow \mathbb{R},i=1,2$ fun\c c\~oes. Suponha que p seja um ponto de acumula\c c\~ao de 
  $D_{f_{1}}\cap{D_{f_{2}}}$ e que $\lim_{x\to p}f_{i}(x) = f_{i}(p), i=1, 2.$ Ent\~ao,
\begin{align*}
 &1) \lim_{x\to p}(f_{1}+f_{2})(x) = \lim_{x\to p}f_{1}(x) + \lim_{x\to p}f_{2}(x) = f_{1}(p)+f_{2}(p)\\
 &2) \lim_{x\to p}kf_{1}(x) = kf_{1}(p)\\
 &3) \lim_{x\to p}f_{1}(x)f_{2}(x) = \lim_{x\to p}f_{1}(x)\lim_{x\to p}f_{2}(x) = f_{1}(p)f_{2}(p)\\
 &4) \text{Se } f_{2}(x)\neq0, \lim_{x\to p}\frac{f_{1}(x)}{f_{2}(x)} = \frac{f_{1}(p)}{f_{2}(p)}.
\end{align*}
\end{theorem*}
\begin{theorem*}
  Se p \'e um ponto de $D_{f}\cap D_{g}$ e $f(x)\leq{g(x)}$ sempre que $x\in{(D_{f}\cap D_{g})}$ e os limites de f e
  quando x tendem a p existem, ent\~ao 
    $$
      \lim_{x\to p}f(x) = f(p)\leq{g(p)} = \lim_{x\to p}g(x).
    $$
\end{theorem*}
Vale tamb\'em a an\'aloga do Teorema do Confronto, mas para fun\c c\~oes cont\'inuas. O resultado a seguir precisava
de continuidade para ser provado:
\begin{theorem*}
  Sejam $f:D_{f}\rightarrow \mathbb{R}, g:D_{g}\rightarrow \mathbb{R}$ fun\c c\~oes tais que a imagem de g est\'a contida
  no dom\'inio de f e $L\in D_{f}.$ Se p \'e um ponto de acumula\c c\~ao de $D_{g}, \lim_{x\to p}g(x) = L$ e f \'e cont\'inua
  em L, ent\~ao 
    $$
      \lim_{x\to p}f(g(x)) = f\biggl(\lim_{x\to p}g(x)\biggr) = f(L).
    $$
\end{theorem*}
\begin{proof*}
  Como f \'e cont\'inua em L, dado $\varepsilon > 0$, exsite $\delta_{f}>0$ tal que 
    $$
      y\in D_{f},\quad |y-L|<\delta_{f} \Rightarrow |f(y)-f(L)|<\varepsilon.
    $$
    Como $\lim_{x\to p}g(x) = L,$ dado $\delta_{f}>0$, existe $\delta_{g}>0$ tal que 
      $$
        x\in D_{g},\quad 0 <|x-p|<\delta_{g} \Rightarrow |g(x)-L|<\delta_{f}.
      $$
    Deta forma, como $Im(g)\subseteq{D_{f}}, D_{f\circ{g}}=D_{g}$ e 
      $$
        x\in D_{g}=D_{f\circ{g}}, \quad 0<|x-p|<\delta_{g}\Rightarrow|g(x)-L|<\delta_{f}\Rightarrow|f(g(x))-L|<\varepsilon.
      $$
      Portanto, $\lim_{x\to p}f(g(x))=f(L).$ \qedsymbol
\end{proof*}
  Em suma, a soma, a multiplica\c c\~ao, a divis\~ao e a composta de fun\c c\~oes cont\'inuas \'e uma fun\c c\~ao cont\'inua. Fun\c c\~oes
  racionais e trigonom\'etricas s\~ao cont\'inuas.
  
\subsection{Resultados Avan\c cados de Continuidade - Parte 1.}
Come\c camos apresentando o resultado conhecido como Teorema da Conserva\c c\~ao do Sinal
 \begin{theorem*}
   Seja $f:D_{f}\rightarrow \mathbb{R}$ uma fun\c c\~ao cont\'inua e $\overline{x}\in D_{f}$ tal que $f(\overline{x})>0.$ 
   Ent\~ao, existe $\delta>0$ tal que $f(x)>0$ para todo $x\in D_{f}$ e $x\in(\overline{x}-\delta, \overline{x}+\delta).$
 \end{theorem*}
 \begin{proof*}
   Como f \'e cont\'inua em $\overline{x},$ dado $\varepsilon = f(\overline{x}) >0$, exsite $\delta > 0$ tal que 
     $$
       x\in D_{f}, x\in(\overline{x}-\delta, \overline{x}+\delta) \Rightarrow f(x)\in (f(\overline{x})-\varepsilon, f(\overline{x})+\varepsilon)) =
  (0, 2f(\overline{x})).
     $$
     Isto prova o resultado. \qedsymbol
 \end{proof*}
 O pr\'oximo \'e o Teorema do Anulamento.
 \begin{theorem*}
   Se $f:[a,b]\rightarrow \mathbb{R}$ \'e cont\'inua e $f(a)<0<f(b),$ ent\~ao existe $\overline{x}\in(a, b)$ tal que $f(\overline{x}) =0.$
 \end{theorem*}
\begin{proof*}
  Faremos apenas o caso $f(a)<0<g(b).$ Seja 
    $$
    A = \{x\in[a,b]:f(s)>0 \forall s\in[x, b].\}
    $$
    Note que $A\neq\emptyset$ e $A\subseteq{[a, b]}.$ Seja $z =\inf{A}$. Pelo Teorema da Conserva\c c\~ao de Sinal,
    $z\in(a, b)$ e $z\not\in A.$ Destarte, $f(z)\leq{0}.$ Por outro lado, do Teorema da Compara\c c\~ao, $f(z)=
    \lim_{x\to z^{+}}f(x)\geq{0}$, pois como $x > z,$ x \'e um elemento de A, o que torna $f(x)>0$. Portanto, f(z)=0.
\end{proof*}
 A seguir, veremos o Teorema do Valor Intermedi\'ario.
\begin{theorem*}
  Seja $f:[a, b]\rightarrow \mathbb{R}$ uma fun\c c\~ao cont\'inua e tal que  $f(a)<f(b)$. Se $f(a)<k<f(b),$
  ent\~ao existe $\overline{x}\in(a, b)$ tal que $f(\overline{x})=k.$
\end{theorem*}
\begin{proof*}
  Considere a fun\c c\~ao $g(x)=f(x)-k.$ Ent\~ao, $g:[a, b]\rightarrow \mathbb{R}$ \'e cont\'inua, $g(a)<0, g(b)>0$
  e do Teorema do Anulamento, existe $\overline{x}\in[a, b]$ tal que $g(\overline{x})=0.$ Portanto, $f(\overline{x})=k.$\qedsymbol
\end{proof*}
Todos eles possuem vers\~oes para trocas de sinais.
\newpage

\section{Aula 15 - 24/04/2023}
\subsection{O que esperar}
\begin{itemize}
  \item O Teorema de Weierstrass;
  \item Topologia da Reta;
  \item Teorema de Borel-Lebesgue.
\end{itemize}
\subsection{Resultados Avan\c cados de Continuidade - Parte 2.}
  Um resultado que vale ser mencionado \'e a caracteriza\c c\~ao de continuidade via limites:
 \begin{theorem*}
   Seja D um subconjunto dos reais, $f:D\rightarrow \mathbb{R}$ e p um ponto de D. A fun\c c\~ao f \'e cont\'inua em p se, e somente
   se, para toda sequ\^encia $\{x_{n}\}$ em D com $x_{n}\overbracket[0pt]{\longrightarrow}^{n\to \infty}p,$ existe o limite $\lim_{n\to \infty}f(x_{n}).$
 \end{theorem*}
\begin{crl*}
  Seja D um subconjunto de $\mathbb{R}, f:D\rightarrow \mathbb{R} $ e p um ponto de D. A fun\c c\~ao f \'e cont\'inua se, e somente
  se, para toda sequ\^encia $\{x_{n}\}$ em D com $x_{n}\overbracket[0pt]{\longrightarrow}^{n\to \infty}p$ temos $f(x_{n})\overbracket[0pt]{\longrightarrow}^{n\to \infty}f(p).$
\end{crl*}
Vamos relembrar os resultados vistos na aula anterior.
 \begin{theorem*}
   Seja $f:D_{f}\rightarrow \mathbb{R}$ uma fun\c c\~ao cont\'inua e $\overline{x}\in D_{f}$ tal que $f(\overline{x})>0.$ 
   Ent\~ao, existe $\delta>0$ tal que $f(x)>0$ para todo $x\in D_{f}$ e $x\in(\overline{x}-\delta, \overline{x}+\delta).$
 \end{theorem*}
 \begin{proof*}
   Como f \'e cont\'inua em $\overline{x},$ dado $\varepsilon = f(\overline{x}) >0$, exsite $\delta > 0$ tal que 
     $$
       x\in D_{f}, x\in(\overline{x}-\delta, \overline{x}+\delta) \Rightarrow f(x)\in (f(\overline{x})-\varepsilon, f(\overline{x})+\varepsilon)) =
  (0, 2f(\overline{x})).
     $$
     Isto prova o resultado. \qedsymbol
 \end{proof*}
 O pr\'oximo \'e o Teorema do Anulamento.
 \begin{theorem*}
   Se $f:[a,b]\rightarrow \mathbb{R}$ \'e cont\'inua e $f(a)<0<f(b),$ ent\~ao existe $\overline{x}\in(a, b)$ tal que $f(\overline{x}) =0.$
 \end{theorem*}
\begin{proof*}
  Faremos apenas o caso $f(a)<0<g(b).$ Seja 
    $$
    A = \{x\in[a,b]:f(s)>0 \forall s\in[x, b].\}
    $$
    Note que $A\neq\emptyset$ e $A\subseteq{[a, b]}.$ Seja $z =\inf{A}$. Pelo Teorema da Conserva\c c\~ao de Sinal,
    $z\in(a, b)$ e $z\not\in A.$ Destarte, $f(z)\leq{0}.$ Por outro lado, do Teorema da Compara\c c\~ao, $f(z)=
    \lim_{x\to z^{+}}f(x)\geq{0}$, pois como $x > z,$ x \'e um elemento de A, o que torna $f(x)>0$. Portanto, f(z)=0.
\end{proof*}
 A seguir, veremos o Teorema do Valor Intermedi\'ario.
\begin{theorem*}
  Seja $f:[a, b]\rightarrow \mathbb{R}$ uma fun\c c\~ao cont\'inua e tal que  $f(a)<f(b)$. Se $f(a)<k<f(b),$
  ent\~ao existe $\overline{x}\in(a, b)$ tal que $f(\overline{x})=k.$
\end{theorem*}
\begin{proof*}
  Considere a fun\c c\~ao $g(x)=f(x)-k.$ Ent\~ao, $g:[a, b]\rightarrow \mathbb{R}$ \'e cont\'inua, $g(a)<0, g(b)>0$
  e do Teorema do Anulamento, existe $\overline{x}\in[a, b]$ tal que $g(\overline{x})=0.$ Portanto, $f(\overline{x})=k.$\qedsymbol
\end{proof*}
As vers\~oes an\'alogas ficam como exerc\'icio para o estudante.
A seguir, um dos resultados fundamentais da an\'alise \'e apresentado, o Teorema de Weierstrass
\begin{theorem*}
  Se $f:[a, b]\rightarrow \mathbb{R}$ for cont\'inua, existir\~ao p, q em $[a, b]$ tais que 
    $$
    f(p)\leq f(x)\leq f(q),\quad \forall x\in[a,b]
    $$
\end{theorem*}
\begin{proof*}
  Vamos mostrar que a imagem da f \'e limitada, incialmente. Se ela n\~ao fosse limitada, dado n natural, existiria $x_{n}\in[a, b]$
  tal que para $x_{0}\in[a, b]$ e $|f(x_{n})| > \max{\{n, |f(x_{n-1})|\}}, n\in \mathbb{N}^{\times}.$ Seja $A=\{x_{n}:n\in \mathbb{N}\}$.
  Com aquele $x_{0},$ obtemos $x_{1}\in[a, b]$ tal que $|f(x_{1})| > 1 + |f(x_{0})|, |f(x_{2})| > 1 + |f(x_{1})| > 2 + |f(x_{0})|.$
  Indutivamente, $|f(x_{n})| > n + |f(x_{0})|.$ Segue que $A\subseteq{[a, b]}$ tem um ponto de acumula\c c\~ao $r\in[a, b].$

  Como f \'e con\t'inua em r, existe $\delta > 0 tal que,$ 
    $$
    x\in(r-\delta , r+\delta )\cap[a, b] = B \Rightarrow |f(x)-f(r)|<1.
    $$

    Obtemos que f(B) \'e limitado e cont\'em infinitos pontos de f(A), uma contradi\c c\~ao. Logo, a imagem de f Im(f) \'e limitada.

    Seja $m=\inf{\{f(x):x\in[a, b]\}}$. Ent\~ao, $f(x)\geq m$ para todo x em $[a, b].$ Se f n\~ao \'e constante, m ser\'a ponto de
    acumula\c c\~ao do conjunto $\{y\in{Im(f)}: y > m\}$. Tome $x_{0}\in[a, b]$ tal que $0 < f(x_{0}) - m$ e $x_{k}\in[a, b]$
    tal que $0 < f(x_{0}) - m < \min{\{\frac{1}{k}, f(x_{k-1})-m\}}$ para k em $\mathbb{N}^{\times}.$

    O conjunto $A = \{x_{k}: k\in \mathbb{N}\}$ \'e infinito e limitado, portanto A tem um ponto de acumula\c c\~ao p. Como f \'e cont\'inua
    em p, para cada n em $\mathbb{N}^{\times},$ existe $\delta_{n} > 0$ tal que 
      $$
      x\in[a, b], |x-p|<\delta_{n} \Rightarrow |f(x)-f(p)|<\frac{1}{n}.
      $$
      Em particular, escolha $x_{k}$ em A com $k>n$ e tal que $|x_{k}-p|<\delta_{n},$ 
        $$
          m-\frac{1}{n} < f(x_{k})-\frac{1}{n}<f(p)<f(x_{k}) + \frac{1}{n} < m + \frac{1}{n} + \frac{1}{k} < m + \frac{2}{n}.
        $$
        Assim, f(p) = m. A afirmativa segue de $-\inf{\{Im(-f)\}} = \sup\{Im(f)\}$ \qedsymbol
\end{proof*}
\begin{crl*}
  Seja $f:[a, b]\rightarrow \mathbb{R}$ cont\'inua. Se 
   $$
    m=\min{\{f(x):x\in[a, b]\}},\quad M = \max{\{f(x):x\in[a, b]\}},
   $$
   ent\~ao 
   $$
     Im(f) = f([a, b]) = [m, M].
   $$
\end{crl*}

\subsection{Topologia da Reta}
\subsubsection{Texto de Motiva\c c\~ao}
A topologia da reta estuda a estrutura dos conjuntos de pontos na reta real, analisando as propriedades topológicas dos mesmos. A reta real é um espaço topológico particularmente interessante devido à sua simplicidade e conexão com outras áreas da matemática, como análise e álgebra.

Considere a sucessão de números reais $(x_n)$. Quando é que podemos dizer que essa sucessão converge para um número real $x$? A topologia da reta nos ajuda a entender a convergência e a continuidade de funções reais.

\begin{center}
\begin{tikzpicture}
\draw[->] (-0.5,0) -- (5,0);
\draw[->] (0,-0.5) -- (0,2);
\node at (5.2,0) {$x$};
\node at (0,2.2) {$y$};

\foreach \x in {1,2,3,4}
\draw (\x,1pt) -- (\x,-1pt) node[anchor=north] {$x_\x$};

\draw[dashed] (1,0) -- (1,1);
\draw[dashed] (2,0) -- (2,0.5);
\draw[dashed] (3,0) -- (3,0.75);
\draw[dashed] (4,0) -- (4,0.875);

\filldraw (1,1) circle (1.5pt);
\filldraw (2,0.5) circle (1.5pt);
\filldraw (3,0.75) circle (1.5pt);
\filldraw (4,0.875) circle (1.5pt);
\end{tikzpicture}
\end{center}

A topologia da reta também lida com conceitos como limites e fronteiras de conjuntos. Por exemplo, qual é o limite superior de um conjunto $A \subseteq \mathbb{R}$? E qual é a fronteira de um conjunto aberto?
Uma ilustra\c c\~ao de conjunto aberto \'e o intervalo (a, b):

\begin{center}
\begin{tikzpicture}
  \draw[->] (-0.5,0) -- (5,0);
\node at (5.2,0) {$x$};

\draw (1,0.1) rectangle (4,0.3);
\draw[thick] (0.5,0) -- (0.5,-0.1) node[below] {$a$};
\draw[thick] (4.5,0) -- (4.5,-0.1) node[below] {$b$};

\draw[<-, thick] (1,0) -- (0.5,0);
\draw[->, thick] (4,0) -- (4.5,0);

\fill[white] (1,0) circle (1.5pt);
\fill[white] (4,0) circle (1.5pt);
\filldraw (1,0) circle (1.5pt);
\filldraw (4,0) circle (1.5pt);
\end{tikzpicture}
\end{center}

Um espaço topológico é um conjunto $X$ juntamente com uma coleção de subconjuntos $\tau$ que satisfazem certas propriedades. Estes subconjuntos são chamados de conjuntos abertos. Na reta real, podemos definir a topologia usual usando intervalos abertos.
Aqui est\'a o c\'odigo representando um conjunto aberto da reta real, tal como seus pontos de fronteira. A barra representa o lugar em que 
os pontos est\~ao, as setas saindo de a e b indicam que os pontos n\~ao est\~ao inclusos no conjunto, tornando-o aberto.

\begin{center}
\begin{tikzpicture}
\draw[->] (-0.5,0) -- (5,0);
\node at (5.2,0) {$x$};

\draw[pattern=north east lines, pattern color=green] (1,0.1) rectangle (4,0.3);
\draw[thick] (0.5,0) -- (0.5,-0.1) node[below] {$a$};
\draw[thick] (4.5,0) -- (4.5,-0.1) node[below] {$b$};

\draw[<-, thick] (1,0) -- (0.5,0);
\draw[->, thick] (4,0) -- (4.5,0);

\filldraw (1,0) circle (1.5pt);
\filldraw (4,0) circle (1.5pt);
\end{tikzpicture}
\end{center}

Na reta real, um conjunto é aberto se, para cada ponto $x$ no conjunto, existe um intervalo aberto contendo $x$ e inteiramente contido no conjunto. Já um conjunto fechado é aquele cujo complementar é aberto.
Outras ideias importantes s\~ao de conectividade e compacidade.

Um conjunto é dito conexo se não puder ser dividido em dois conjuntos abertos disjuntos. Na reta real, intervalos são exemplos de conjuntos conexos.

Um conjunto é compacto se todo subconjunto aberto que o cobre possui um subconjunto finito que também o cobre. Na reta real, intervalos fechados e limitados são exemplos de conjuntos compactos.
O Teorema de Bolzano-Weierstrass afirma que toda sequência infinita de pontos em um conjunto compacto possui uma subsequência convergente. Como esperado,
ele possui aplica\c c\~oes diretas na converg\^encia de sequ\^encias e fun\c c\~oes juntando com os resultados que vimos.

Suponha que você tenha uma sequência infinita de pontos dentro de um intervalo fechado e limitado na reta real. O Teorema de Bolzano-Weierstrass garante que sempre é possível encontrar uma subsequência desses pontos que converge para algum ponto dentro desse intervalo. Isso significa que, por mais ``espalhados" que esses pontos possam estar, sempre haverá algum padrão de convergência.

Portanto, a topologia da reta fornece uma base sólida para o estudo de propriedades topológicas e estruturas dos conjuntos de pontos na reta real. Aprofundemo-nos
em seu estudo.
\subsubsection{Abertos, Fechados, Compactos e Conexos}
\begin{def*}
  Seja A um subconjunto de $\mathbb{R}.$
 \begin{itemize}
   \item[1)] Um ponto a de A \'e dito interior a A se existe um r positivo tal que $(a-r, a+r)\subseteq{A}.$
   \item[2)] O conjunto A \'e aberto se todos os seus pontos s\~ao interiores.
   \item[3)] O conjunto A \'e fechado em $\mathbb{R}$ se seu complementar \'e aberto.
   \item[4)] O interior de A, $A^{\circ},$ \'e o conjunto dos pontos interiores a A.
   \item[5)] O fecho $\overline{A}$ de A \'e a interse\c c\~ao de todos os fechados que cont\'em A.
 \end{itemize}
\end{def*}
\'E poss\'ivel mostrar que intervalos abertos s\~ao conjuntos abertos (tome $r=\min{\{b-x, x-a\}}, x\in(a, b)$), intervalos fechados s\~ao conjuntos fechados, etc. 
\begin{theorem*}
  Seja $\Lambda $ um conjunto e $\{A_{\lambda }\}_{\lambda \in\Lambda } $ uma fam\'ilia de conjuntos.
 \begin{itemize}
  \item[1)] Se cada $A_{\lambda }$ \'e aberto, $\bigcup_{\lambda\in\Lambda  }^{}{A_{\lambda }}$ \'e aberto.
  \item[2)] Se cada $A_{\lambda }$ \'e aberto e $\Lambda $ \'e finito, $\bigcap_{\lambda \in\Lambda }^{}{A_{\lambda }}$
    \'e aberto.
  \item[3)] Se cada $A_{\lambda }$ \'e fechado, $\bigcap_{\lambda \in\Lambda }^{}{A_{\lambda }}$ \'e fechado.
  \item[4)] Se cada $A_{\lambda }$ \'e fechado e $\Lambda $ \'e um conjunto finito, $\bigcup_{\lambda \in\Lambda }^{}{A_{\lambda }}$
    \'e fechado.
 \end{itemize}
 (1) Como cada $A_{\lambda }$ \'e composto de todos os pontos interiores, $\bigcup_{\lambda \in\Lambda }^{}{A_{\lambda }}$ tamb\'em tem
 todos os seus pontos como pontos interiores. Portanto, $\bigcup_{\lambda \in\Lambda }^{}{A_{\lambda }}$ \'e aberto.

 (2) Tome x em $\bigcap_{\lambda \in\Lambda }^{}{A_{\lambda }}.$ Ent\~ao, existe um $r_{\lambda }$ tal que $(x-r_{\lambda }, x+r_{\lambda })\subseteq{A_{\lambda }}.$
 Como temos um n\'umero finito de abertos, tome $r\coloneqq \min{\{r_{\lambda }:\lambda \in\Lambda \}}$. Ent\~ao, 
   $$
   (x-r, x+r)\subseteq{(x-r_{\lambda },x+r_{\lambda })}\subseteq{A}.
   $$
   Portanto, $\bigcap_{\lambda \in\Lambda }^{}{A_{\lambda }}$ \'e aberto.

  Os itens (3) e (4) seguem utilizando os itens (1) e (2) e as leis de De Morgan (complementar da uni\~ao \'e a interse\c c\~ao dos complementares, etc.). \qedsymbol
\end{theorem*}
\begin{example}
  Mostramos que intervalos abertos s\~ao abertos. N\~ao apenas isso, mas o interior de [a, b] \'e (a, b).
\end{example}
\begin{example}
  O interior $A^{\circ}$ de A \'e o maior aberto que cont\'em A. O fecho $\overline{A}$ de A \'e o menor fechado que cont\'em A.
\end{example}
\begin{theorem*}
  Todo subconjunto aberto A de $\mathbb{R}$ se exprime, de maneira \'unica, como uni\~ao enumer\'avel de intervalos abertos disjuntos.
\end{theorem*}
\begin{proof*}
  Primeiramente, note que se $\Lambda $ \'e um conjunto, para cada $\lambda, I_{\lambda }=(a_{\lambda }, b_{\lambda })$ \'e um intervalo
  e $p\in\bigcap_{\lambda \in\Lambda }^{}{I_{\lambda }}$, ent\~ao $\bigcup_{\lambda \in\Lambda }^{}{I_{\lambda }}=(a, b),$ sendo
  $a = \inf_{\lambda \in\Lambda }a_{\lambda }, b = \sup_{\lambda \in\Lambda }b_{\lambda }.$ De fato, \'e claro que $\bigcup_{\lambda \in\Lambda }^{}{I_{\lambda }}\subseteq{(a, b)}$.
  Para provar a outra inclus\~ao, note que $p\in(a, b)$ e se $x\in(a, b), x\leq p ou x > p.$ Agora, 
 \begin{align*}
   &(1) \text{ Se }  x\leq p, a = \inf_{\lambda \in\Lambda }a_{\lambda } < x \Rightarrow \exists \mu_{1}\in\Lambda: a_{\mu_{1}}<x\leq p < b_{\mu_{1}}\\
   &(2) \text{ Se } x > p, b=\sup_{\lambda \in\Lambda }b_{\lambda } < x \Rightarrow \exists \mu_{2}\in\Lambda: a_{\mu_{2}} < p < x < b_{\mu_{2}}.
 \end{align*}
  Em qualquer caso,  $x\in \bigcup_{\lambda \in\Lambda }^{}{I_{\lambda }}.$

  Para terminar, dado x em A, seja $I_{x}$ a uni\~ao de todos os abertos contidos em A que cont\'em x. Segue que 
 \begin{itemize}
   \item[1)] $I_{x} = (a_{x}, b_{x})\subseteq{A};$
   \item[2)] Se x, y s\~ao pontos de A, ou $I_{x}\cap I_{y}=I_{x} = I_{y},$ ou $I_{x}\cap I_{y} = \emptyset$;
   \item[3)] $\bigcup_{x\in A}^{}{I_{x}} = A.$
 \end{itemize}
 Tomando para cada intervalo da decomposi\c c\~ao um \'unico racional, vemos que A pode ser escrito como uni\~ao enumer\'avel de
 intervalos disjuntos. Para ver que esta decomposi\c c\~ao \'e \'unica, basta notar que cada intervalo aberto de uma tal decomposi\c c\~ao
 est\'a contido em algum dos $I_{x}$ e n\~ao pode ser distinto de $I_{x}.$ \qedsymbol
\end{proof*}
\begin{crl*}
  Se I \'e um intervalo aberto e $I = A\cup B,$ em que A e B s\~ao conjuntos abertos e disjuntos, ent\~ao um desses conjuntos \'e vazio.
\end{crl*}
\newpage

\section{Aula 16 - 03/05/2023}
\subsection{O que esperar}
\begin{itemize}
  \item Continua\c c\~ao de Topologia da Reta - Compactos e Conexos.
\end{itemize}
\subsection{Topologia da Reta - Parte II}
  Para ver que a intersec\c c\~ao infinita de intervalos abertos n\~ao \'e necessariamente um conjunto aberto, note que 
    $$
      \bigcap_{n\in \mathbb{N}^{\times}}^{}{(a-\frac{1}{n}, b+\frac{1}{n})} = [a, b].
    $$
  No caso da uni\~ao infinita de conjuntos fechados, observe o seguinte: 
    $$
     \bigcup_{n\in \mathbb{N}^{\times}}^{}{[a+\frac{1}{n}, b - \frac{1}{n}]}=(a, b)
    $$
  Retomemos o teorema passado.
  \begin{theorem*}
    Todo subconjunto aberto A de $\mathbb{R}$ se exprime, de maneira \'unica, como uni\~ao enumer\'avel de intervalos abertos disjuntos.
  \end{theorem*}
  \begin{proof*}
    Primeiramente, note que se $\Lambda $ \'e um conjunto, para cada $\lambda, I_{\lambda }=(a_{\lambda }, b_{\lambda })$ \'e um intervalo
    e $p\in\bigcap_{\lambda \in\Lambda }^{}{I_{\lambda }}$, ent\~ao $\bigcup_{\lambda \in\Lambda }^{}{I_{\lambda }}=(a, b),$ sendo
    $a = \inf_{\lambda \in\Lambda }a_{\lambda }, b = \sup_{\lambda \in\Lambda }b_{\lambda }.$ De fato, \'e claro que $\bigcup_{\lambda \in\Lambda }^{}{I_{\lambda }}\subseteq{(a, b)}$.
    Para provar a outra inclus\~ao, note que $p\in(a, b)$ e se $x\in(a, b), x\leq p ou x > p.$ Agora, 
   \begin{align*}
     &(1) \text{ Se }  x\leq p, a = \inf_{\lambda \in\Lambda }a_{\lambda } < x \Rightarrow \exists \mu_{1}\in\Lambda: a_{\mu_{1}}<x\leq p < b_{\mu_{1}}\\
     &(2) \text{ Se } x > p, b=\sup_{\lambda \in\Lambda }b_{\lambda } < x \Rightarrow \exists \mu_{2}\in\Lambda: a_{\mu_{2}} < p < x < b_{\mu_{2}}.
   \end{align*}
    Em qualquer caso,  $x\in \bigcup_{\lambda \in\Lambda }^{}{I_{\lambda }}.$

    Para terminar, dado x em A, seja $I_{x}$ a uni\~ao de todos os abertos contidos em A que cont\'em x. Segue que 
   \begin{itemize}
     \item[1)] $I_{x} = (a_{x}, b_{x})\subseteq{A};$
     \item[2)] Se x, y s\~ao pontos de A, ou $I_{x}\cap I_{y}=I_{x} = I_{y},$ ou $I_{x}\cap I_{y} = \emptyset$;
     \item[3)] $\bigcup_{x\in A}^{}{I_{x}} = A.$
   \end{itemize}
   Tomando para cada intervalo da decomposi\c c\~ao um \'unico racional, vemos que A pode ser escrito como uni\~ao enumer\'avel de
   intervalos disjuntos. Para ver que esta decomposi\c c\~ao \'e \'unica, basta notar que cada intervalo aberto de uma tal decomposi\c c\~ao
   est\'a contido em algum dos $I_{x}$ e n\~ao pode ser distinto de $I_{x}.$ \qedsymbol
  \end{proof*}
  \begin{crl*}
    Se I \'e um intervalo aberto e $I = A\cup B,$ em que A e B s\~ao conjuntos abertos e disjuntos, ent\~ao um desses conjuntos \'e vazio.
  \end{crl*}
 \begin{def*}
   Um conjunto I em $\mathbb{R}$ \'e dito conexo se ele s\'o ``tem um peda\c co'', ou seja, n\~ao existem conjuntos abertos e disjuntos A, B
   n\~ao-vazios tais que $I = A\cup{B}.\square$
 \end{def*}
 \begin{def*}
   Seja $A\subseteq{\mathbb{R}}$. Um ponto p real \'e aderente a A se existir uma sequ\^encia $\{x_{n}\}$ em A tal que $x_{n}\overbracket[0pt]{\longrightarrow}^{n\to \infty}p.\square$ 
 \end{def*}
  Sabemos que se p \'e um ponto de acumula\c c\~ao de A, ent\~ao p \'e aderente a A. Se p \'e aderente a A e p n\~ao \'e ponto de acumula\c c\~ao de A, ent\~ao
  p pertence a A. Todo ponto interior a A \'e aderente a A e \'e um ponto de acumula\c c\~ao de A.
 \begin{theorem*}
   Um ponto p real \'e aderente a A se, e s\'o se, $A\cap{(p-\varepsilon , p+\varepsilon )}\neq\emptyset$ para todo $\varepsilon  > 0.$
 \end{theorem*}
 \begin{crl*}
   Se $A\subseteq{\mathbb{R}}$ \'e limitado superiormente(inferiormente), ent\~o $\sup{A}(\inf{A})$ \'e aderente a A.
 \end{crl*}
\begin{theorem*}
  O fecho $A^{-}$ de $A\subseteq{\mathbb{R}}$ \'e o conjunto $A'$ dos pontos aderentes de A.
\end{theorem*}
\begin{proof*}
  De fato, se $x\not\in A',$ existe $\varepsilon >0$ tal que $(x-\varepsilon ,x+\varepsilon )\cap{A}\neq\emptyset$. Segue que
  $A'$ \'e fechado e $A^{-}\subseteq{A'}$. Se $x\not\in A^{-}$, existe $\varepsilon >0$ tal que $(x-\varepsilon ,x+\varepsilon )\cap{A}=\emptyset$
  e $x\not\in A'.$ \qedsymbol
\end{proof*}
\begin{def*}
  Sejam A e B subconjuntos de $\mathbb{R}$ com $A\subseteq{B}.$ Diremos que A \'e denso em B se $B\subseteq{A^{-}}\square$.
\end{def*}
\begin{theorem*}
  Sejam A e B subconjuntos de $\mathbb{R}$ com $A\subseteq{B}$. S\~ao equivalentes:
 \begin{itemize}
  \item Todo ponto de B \'e aderente a A.
  \item Todo ponto de B \'e limite de uma sequ\^encia de pontos de A.
  \item Para todo $\varepsilon >0$ e b de B, $(b-\varepsilon , b+\varepsilon )\cap{A}\neq\emptyset.$
 \end{itemize}
\end{theorem*}
\begin{theorem*}
  Todo subconjunto B de $\mathbb{R}$ cont\'em um subconjunto A que \'e enumer\'avel e denso em B.
\end{theorem*}
\begin{proof*}
  Dado $n\in \mathbb{N}^{\times}$, temos $\mathbb{R} = \displaystyle\bigcup_{p\in \mathbb{Z}}^{}{\biggl[\frac{p}{n}, \frac{p+1}{n}\biggr)}$. Para
  cada p inteiro e $n\in \mathbb{N}^{\times}$, escolha $x_{np}\in\biggl[\frac{p}{n}, \frac{p+1}{n}\biggr)\cap{B}$ quando esta intersec\c c\~ao
  for n\~ao-vazia. O conjunto A desses pontos \'e denso em B, j\'a que, para cada $n\in \mathbb{N}^{\times}$ e b em B, existe a em A
  tal que $|a-b|<\frac{1}{n}$. Al\'em disso, A \'e enumer\'avel, visto que a cole\c c\~ao de intervalos $\biggl\{\biggl[\frac{p}{n}, \frac{p+1}{n}\biggr): p\in \mathbb{Z}, n\in \mathbb{N}^{\times}\biggr\}$
  \'e enumer\'avel.\qedsymbol
\end{proof*}
\begin{theorem*}
  Seja $F\subseteq{\mathbb{R}}$ um conjunto fechado, n\~ao-vazio e sem pontos isolados. Ent\~ao F \'e n\~ao enumer\'avel.
\end{theorem*}
\begin{proof*}
  Sejam x, y em F pontos distintos. Coloque $r = \frac{|x-y|}{2}$ e $F_{y}^{\sim} = F\cap{(x-r, x+r)}.$ Segue que $F_{y}^{\sim}$
\'e n\~ao-vazio e n\~ao cont\'em pontos isolados. Seja $F_{y}$ a uni\~ao de $F_{y}^{\sim}$ com os pontos de acumula\c c\~ao de 
$F_{y}^{\sim}$ no conjunto $\{x-r,x+r\}.$ Note que $F_{y}$ \'e claramente fechado e n\~ao tem pontos isolados, \'e limitado e
$y\not\in F.$

  Se $\{y_{1},y_{2},\cdots\},$ seja $F_{y_{1}}.$ Tendo escolhido $F_{y_{1}},\cdots, F_{y_{n-1}}$, ent\~ao caso $y_{n}\not\in F_{y_{n-1}},$ 
escolhemos $F_{y_{n}}=F_{y_{n-1}}.$ Se $y_{n}\in F_{y_{n-1}},$ escolhemos $F_{y_{n}}$ fechado e sem pontos isolados tal que 
$y_{n}\not\in F_{y_{n}}\subseteq{F_{y_{n-1}}}.$ Para cada n natural, seja $x_{n}\in F_{y_{n}}.$ A sequ\^encia $\{x_{n}\}$ \'e limitada
e portanto tem uma subsequ\^encia $\{x_{\varphi(n)}\}$ convergente com limite $\overline{x}$. Observe que $\overline{x}\in \bigcap_{n\in \mathbb{N}}^{}{F_{y_{n}}}$ 
e $\overline{x}\neq y_{n}$ para todo natural n. Portanto, F \'e n\~ao-enumer\'avel, pois a sequ\^encia ``n\~ao enche F'' (N\~ao abrange
todos os pontos de F). \qedsymbol
\end{proof*}
\newpage

\section{Aula 17 - 05/04/2023}
\subsection{O que esperar?}
\begin{itemize}
  \item Coberturas e Teorema de Borel-Lebesgue;
  \item Compactos da Reta;
  \item Continuidade Uniforme.
\end{itemize}
\subsection{Motivações}
Um subconjunto de $\mathbb{R}$ é dito compacto se satisfaz uma das condições abaixo - a caracterização por coberturas, por sequências, ou
por fechados.

Um conjunto $K$ é compacto se, para toda cobertura aberta do conjunto, existe uma subcobertura finita que ainda cobre o conjunto.
Ou seja, se temos um conjunto de intervalos abertos que cobrem todo o conjunto, sempre podemos encontrar um número finito desses intervalos
que ainda cobrem o conjunto.

Um conjunto $K$ é compacto se toda sequência de pontos do conjunto tem uma subsequência que converge para um ponto ainda dentro
do conjunto. Ou seja, se pegarmos uma lista infinita de pontos no conjunto, sempre poderemos encontrar uma "sublista" infinita 
desses pontos que se aproxima cada vez mais de um único ponto no conjunto.

Um conjunto $K$ é fechado se contém todos os seus pontos de acumulação, e é limitado se existe um número real que é maior que todos os elementos do conjunto.

Além dos compactos, estaremos abordando outro tipo de continuidade. A continuidade uniforme é uma noção mais forte do que a continuidade regular. Seja $f: X \rightarrow Y$ uma função entre espaços métricos. Uma função é dita continua se, para todo $\varepsilon > 0$, existe um $\delta > 0$ tal que, para todos $x, y \in X$, se $|x - y| < \delta$ então $|f(x) - f(y)| < \varepsilon$.

A continuidade uniforme fortalece essa condição ao exigir que a escolha de $\delta$ não dependa de $x$. Em outras palavras,
para toda $\varepsilon > 0$, existe um $\delta > 0$ tal que, para todos $x, y \in X$, se $|x - y| < \delta$ então 
$|f(x) - f(y)| < \varepsilon$. Isso implica que podemos encontrar um único "tamanho de passo" $\delta$ que funciona para todo
o domínio da função, em vez de precisar de um passo diferente para cada ponto.
Isso tem implicações importantes quando estamos trabalhando com conjuntos compactos, pois toda função contínua em um
conjunto compacto é uniformemente contínua. Isso é útil porque as funções uniformemente contínuas têm propriedades muito agradáveis,
como a capacidade de serem estendidas de um modo contínuo a partir de um conjunto denso.

A diferença entre continuidade e continuidade uniforme é sutil. No entanto, essa diferença é importante em muitos contextos. 
Por exemplo, a função $f(x) = \frac{1}{x}$ é contínua em todos os pontos do seu domínio (todos os números reais não nulos), 
mas não é uniformemente contínua. Isso ocorre porque, conforme nos aproximamos de zero, precisamos escolher $\delta$ cada vez
menor para garantir que $d(f(x), f(x')) < \varepsilon$ para um $\varepsilon$ fixo. Não podemos encontrar um único $\delta$ que
funcione para todos os pontos do domínio.

Outro fato importante é que toda função contínua em um conjunto compacto é uniformemente contínua. Isso se deve às propriedades
dos conjuntos compactos, que garantem que podemos encontrar um $\delta$ que funcione para todo o conjunto. Isso é uma
propriedade muito útil, pois a continuidade uniforme garante algumas propriedades agradáveis para as funções, como a
capacidade de serem estendidas de maneira contínua a partir de um conjunto denso.

A seguir, vamos elaborar cada um destes conceitos, a fim de obter um entendimento mais profundo sobre o estudo de funções e da 
topologia da reta.

\subsection{Compactos e Coberturas}
 \begin{def*}
   Seja $A\subseteq{\mathbb{R}}$ e $\Lambda $ um conjunto. Uma coleção $\{A_{\lambda }\}_{\lambda \in\Lambda }$ de conjuntos
   é chamada uma cobertura de A se $A\subseteq{\bigcup_{\lambda \in\Lambda }^{}{A_{\lambda }}}.\square$
 \end{def*}
\begin{def*}
  Se $\{A_{\lambda }\}_{\lambda \in\Lambda }$ é uma cobertura, $\Lambda '\subseteq{\Lambda }$ e $A\subseteq{\bigcup_{\lambda '\in\Lambda '}^{}{A_{\lambda '}}},
  \{A_{\lambda '}\}_{\lambda '\in\Lambda '}$ é dita uma subcobertura da cobertura $\{A_{\lambda }\}_{\lambda \in\Lambda }.$
  Se os conjuntos da cobertura são todos abertos, ela é dita uma cobertura aberta. $\square$
\end{def*}
\begin{theorem*}
  Este resultado é conhecido como Teorema de Borel-Lebesgue. Dada uma cobertura $\{I_{\lambda }\}_{\lambda \in\Lambda }$ de [a, b]
  em que cada $I_{\lambda }$ é um intervalo aberto, existe $\Lambda '\subseteq{\Lambda }$ finito tal que $[a, b]\subseteq{\bigcup_{\lambda '\in\Lambda '}^{}{I_{\lambda '}.}}$
\end{theorem*}
\begin{proof*}
  Seja $A=\{x\in[a,b]: \exists \Lambda'\subseteq{\Lambda }\text{ finito }\&[a,x]\subseteq{\bigcup_{\lambda '\in\Lambda '}^{}{I_{\lambda '}}}\}$.
  Note que $A \neq\emptyset.$ Seja $s=\sup{A}.$ Segue que $s\in[a, b]$ e que existe $\lambda _{s}$ tal que $s\in I_{\lambda _{s}}.$
  Como $I_{\lambda _{s}}$ é aberto, $I_{\lambda _{s}}\cap A \neq\emptyset$. Portanto, s = b e $[a, b]$ está contido em uma
  união finita de $I_{\lambda }'s.$ \qedsymbol
\end{proof*}
\begin{crl*}
  Dada uma cobertura aberta $\{A_{\lambda }\}_{\lambda \in\Lambda }$ de [a, b], existe $\Lambda '\subseteq{\Lambda }$ finito
  tal que $[a,b]\subseteq{\bigcup_{\lambda '\in\Lambda '}^{}{A_{\lambda'}}}.$
\end{crl*}
  Basta lembrar que cada aberto da cobertura pode ser escrito como a união enumerável de intervalos abertos disjuntos.
 \begin{crl*}
   Dada uma cobertura aberta $\{A_{\lambda }\}_{\lambda \in\Lambda }$ de um conjunto fechado e limitado F existe $\Lambda '\subseteq{\Lambda }$
   finito tal que $F\subseteq{\bigcup_{\lambda '\in\Lambda '}^{}{A_{\lambda '}}}$
 \end{crl*}
 \begin{proof*}
   Como F é fechado e limitado, $F\subseteq{[a,b]}$. Como $A=[a,b]^{c}$ é aberto e $[a,b]\subseteq{\bigcup_{\lambda \in\Lambda }^{}{A_{\lambda }}\cup{A}}$, temos
   a existência de $\Lambda '\subseteq{\Lambda }$ finito tal que 
     $$
     [a,b] \subseteq{\bigcup_{\lambda '\in\Lambda '}^{}{A_{\lambda '}}\cup{A}}.
     $$
     Portanto, $F\subseteq{\bigcup_{\lambda '\in\Lambda '}^{}{A_{\lambda '}}}.$ \qedsymbol
 \end{proof*}
 \begin{theorem*}
   Dado $K\subseteq{\mathbb{R}}$ são equivalente: 
  \begin{itemize}
    \item[1)] K é fechado e limitado;
    \item[2)] Toda cobertura aberta de K possui uma subcobertura finita;
    \item[3)] Todo subconjunto infinito de K possui um ponto de acumulação pertencente a K;
    \item[4)] Toda sequência de pontos de K possui uma subsequência que converge para um ponto de K.
  \end{itemize}
 \end{theorem*}
\begin{proof*}
  $1) \Rightarrow 2):$ Segue diretamente do corolário anterior.

  $2) \Rightarrow 3):$ Se $A\subseteq{K}$ é infinito e não tem pontos de acumulação em K, para cada k em K existe $I_{k}=r_{k} > 0$
  tal que $(k-r_{k}, k+r_{k})\cap{A} = \{k\}$ ou  $I_{k}\cap A = \emptyset.$ Segue que $\bigcup_{k\in K}^{}{I_{a}}\supseteq{K}$ é
  uma cobertura aberta sem subcobertura finita.

  $3) \Rightarrow 4):$ Dada uma sequência de pontos $\{k_{n}\}$ em K, ela pode ter um número finito ou infinito de valores.
  Em qualquer um dos casos, possui uma subsequência convergente.

  $4) \Rightarrow 1):$ É claro que K é limitado, visto que, caso contrário, existiria uma sequência $\{x_{n}\}$ em K com $x_{0}\in K$
  e $|x_{n}|\geq |x_{n-1}| + 1$ e esta não teria subsequência convergente. Para ver que K é fechado, basta notar que, se $x\in K^{-}$,
  existe sequência $x_{n}\in K$ tal que $x_{n}\overbracket[0pt]{\longrightarrow}^{n\to \infty}x$. Portanto, pela hipótese de 4, $x\in K.$\qedsymbol
\end{proof*}
\begin{def*}
  Um conjunto é compacto se satisfaz qualquer uma das condições do teorema anterior. $\square$
\end{def*}
  Como corolário do teorema passado, temos o Teorema de Bolzano-Weierstrass!
\begin{crl*}
  Todo conjunto infinito e limitado de número reais possui um ponto de acumulação. 
\end{crl*}
\begin{crl*}
  Toda sequência decrescente de compactos não vazio têm interseção não vazia.
\end{crl*}
\begin{def*}
  Seja $A\subseteq{\mathbb{R}}$ e $B\subseteq{A}$, diremos que B é aberto em A se, para cada $b\in B$, existe um $r_{b}>0$ tal 
  que $A\cap (b-r_{b}, b + r_{b})\subseteq{B}.\square$
\end{def*}
  Note que todo conjunto é aberto nele mesmo. Além disso, se $A\subseteq{\mathbb{R}}$ e $B\subseteq{A}$, então B é aberto em A
  se, e somente se, existe um aberto $\mathcal{O}_{B}$ de $\mathbb{R}$ tal que $B=\mathcal{O}_{B}\cap A.$ Além disso,
  se B é aberto, $A\subseteq{B}$ é aberto em B se, e somente se, B é aberto em $\mathbb{R}.$

  Lembre-se que, se $f:D\rightarrow \mathbb{R}$ é uma função, $f^{-1}(\mathcal{O})=\{d\in D: f(d)\in \mathcal{O}\}.$

 \begin{theorem*}
   Seja $D\subseteq{\mathbb{R}}$ e $f:D\rightarrow \mathbb{R}.$ A função f é contínua se, e somente se, para todo
   $\mathcal{O}$ de $\mathbb{R}, f^{-1}(\mathcal{O})$ é aberto em D.
 \end{theorem*}
 \begin{proof*}
   Se $f:D\rightarrow \mathbb{R}$ é contínua, $\mathcal{O}$ é um aberto de $\mathbb{R}$ e $d\in f^{-1}(\mathcal{O}),$
   então $f(d)\in \mathcal{O}$ e dado $\varepsilon  > 0$ tal que $(f(d)-\varepsilon , f(d)+\varepsilon )\subseteq{\mathcal{O}},$
   existe $\delta >0$ tal que $f((d-\delta , d+\delta )\cap D)\subseteq{(f(d)-\varepsilon , f(d)+\varepsilon ).}$ Isot mostra que
   $(d-\delta , d+\delta )\cap D\subseteq{f^{-1}(\mathcal{O})}$ e que $f^{-1}(\mathcal{O})$ é aberto em D.

   Por outro lado, de $f^{-1}(\mathcal{O})$ ser aberto em D para cada $\mathcal{O}$ aberto em $\mathbb{R},$ se d é um elemento
   de D, dado $\varepsilon  > 0$, seja $\mathcal{O} = (f(d)-\varepsilon , f(d)+\varepsilon ).$ Como $d\in f^{-1}((f(d)-\varepsilon , f(d)+\varepsilon ))$
   é aberto em D, existe $\delta >0$ tal que $(d-\lambda , d+\delta )\cap D\subseteq{f^{-1}((f(d)-\varepsilon , f(d)+\varepsilon ))},$
   ou seja, 
     $$
       x\in D, |x-d| < \delta \Rightarrow |f(x)-f(d)| < \varepsilon 
     $$
    e f é contínua em d. \qedsymbol
 \end{proof*}
 \begin{theorem*}
   Se $I\subseteq{\mathbb{R}}$ é um intervalo e $f:I\rightarrow \mathbb{R}$ é uma função contínua, então f(I) é um intervalo.
 \end{theorem*}
\begin{proof*}
  Basta notar que, dados dois pontos $f(a)\neq f(b)$ em f(I) com $a < b,$ tomando a restrição de f ao intervalo [a, b],
  segue do teorema do valor intermediário que, para todo k entre f(a) e f(b), existe um c em (a, b) tal que f(c) = k. Portanto,
  f(I) é um intervalo. \qedsymbol
\end{proof*}
\begin{theorem*}
  Se $K\subseteq{\mathbb{R}}$ é um conjunto compacto e $f:K\rightarrow \mathbb{R}$ é uma função contínua, então f(K) é compacto.
\end{theorem*}
\begin{proof*}
  Seja $\{\mathcal{O}_{\lambda }:\lambda \in\Lambda \}$ uma cobertura aberta de f(K). Como, para cada $\lambda \in \Lambda, f^{-1}(\mathcal{O}_{\lambda })$ 
  é aberto em K, existe $U_{\lambda }$ aberto em $\mathbb{R}$ tal que $U_{\lambda }\cap K = f^{-1}(\mathcal{O}_{\lambda }).$
  Assim, $\{U_{\lambda }:\lambda \in\Lambda \}$ é uma cobertura aberta de K. Como K é compacto, existe $\lambda '\subseteq{\Lambda }$
  finito tal que $\bigcup_{\lambda '\in\Lambda '}^{}{U_{\lambda '}}\supseteq{K}.$ Segue que $\{\mathcal{O}_{\lambda '}: \lambda '\in\Lambda '\}$
  é uma subcobertura finita da cobertura $\{\mathcal{O}_{\lambda }:\lambda \in \Lambda \}$ de f(K). Portanto, f(K) é compacto. \qedsymbol
\end{proof*}
\begin{proof*}
  Uma prova alternativa faz uso de sequências: Seja $\{y_{n}\}$ uma sequência em f(K). Então, existe uma sequência
  $\{x_{n}\}$ em K tal que $y_{n}=f(x_{n}).$ Como K é compacto, $\{x_{n}\}$ tem uma subsequência $\{x_{\varphi (n)}\}(\varphi :\mathbb{N}\rightarrow \mathbb{N}$
  estritamente crescente) convergente com limite $\overline{x}\in K.$ Como $x_{\varphi (n)}\overbracket[0pt]{\longrightarrow}^{n\to \infty}\overline{x},
  y_{\varphi (n)}=f(x_{\varphi (n)})\overbracket[0pt]{\longrightarrow}^{n\to \infty}f(x)$ e $\{y_{n}\}$ tem uma
  subsequência convergente com limite em f(K). Portanto, f(K) é compacto. \qedsymbol
\end{proof*} 
\begin{theorem*}
  Se $K\subseteq{\mathbb{R}}$ é um conjunto compacto e $f:K\rightarrow \mathbb{R}$ é contínua, existem $x_{1},x_{2}\in K$ tal que
  $f(x_{1})\leq f(x)\leq f(x_{2})$ para todo x em K.
\end{theorem*}
\begin{proof*}
  Como f(K) é compacto, $L=\sup{\{y:y\in f(K)\}}$ e $l=\inf{\{y:y\in f(K)\}}$ pertencem a f(K). Portanto, existem $x_{1},x_{2}\in K$
  tais que $f(x_{1})=l\leq f(x)\leq L = f(x_{2})$ para todo x de K. \qedsymbol
\end{proof*}

\section{Aula 18 - 10/05/2023}
\subsection{O que esperar?}
\begin{itemize}
  \item Semicontinuidades;
  \item Diferenciabilidade.
\end{itemize}
\subsection{Semicontinuidade}
\begin{def*}
  Se \(D\subseteq{\mathbb{R}}\), c é um ponto de acumula\c cão de D e \(f:D\rightarrow \mathbb{R}\) 
  é uma fun\c cão que é limitada em uma vizinhan\c ca de c em \(\mathbb{R}\), definimos 
    \[
      \limsup_{x\to c}f(x)\limsup_{r\to 0^{+}}\{f(x):x\in D, 0 < |x-c|< r\}\quad\text{ e } \liminf_{x\to c}f(x)=\liminf_{r\to 0^{+}}{f(x):x\in D, 0 < |x-c|<r}.
    \]
    Definimos os análogos para pontos em \(c\in D^{-}\) como
      \[
        \overline{\text{Lim}}_{x\to p}f(x)\coloneqq \limsup_{x\to p}f(x)\coloneqq \limsup_{r\to 0^{+}}\{f(x):x\in D, |x-p| < r\}\quad\text{ e } \underline{\text{Lim}}_{x\to p}f(x)\coloneqq \liminf_{r\to 0^{+}}\{f(x):x\in D, |x-p|<r||\}square
      \]
\end{def*}
\begin{def*}
  Seja \(f:D\rightarrow \mathbb{R}\) e c um ponto de D. Então, f é semicontínua
  superiormente em c se 
    \[
      f(c) = \overline{Lim}_{x\to c}f(x) \quad(f(c)\geq \limsup_{x\to p}f(x))
    \]
    e semicontínua inferiormente se 
      \[
        f(c) = \underline{Lim}_{x\to c}f(x) \quad(f(c)\leq \liminf_{x\to p}f(x))
      \]
    Se f é semicontínua superiormente (inferiormente) em todos os pontos de D, dizemos
    que f é semicontínua superiormente (inferiormente). \(\square\)
\end{def*}
 \begin{theorem*}
   Seja \(f:D\rightarrow \mathbb{R}\) uma fun\c cão semicontínua superiormente (inferiormente).
   Se k é um número real, então existe um aberto \(O_{k}\) de \(\mathbb{R}\) tal que 
     \[
       O_{k}\cap D = \underbrace{\{x\in D: f(x)<k\}}_{f^{-1}((-\infty, k))} \quad(O_{k}\cap D=\underbrace{\{x\in D: f(x)>k\}}_{f^{-1}((k, \infty))})
     \]
 \end{theorem*}
\begin{proof*}
  Para c em D com \(f(c) < k,\), da defini\c cão de semicontinuidade, segue que
  existe \(r_{c}>0\) tal que \(f(x) < k\) para todo x em D, \(|x-c| < r_{c}.\) Seja
  \(I_{c} = (c-r_{c}, c+r_{c})\) e defina 
    \[
      O_{k} = \bigcup_{c\in f^{-1}((-\infty, k))}^{}{I_{c}.}
    \] 
    De fato, se c pertence a esta interse\c cão, então \(f(c) < k\). Com isso, 
      \[
        k > f(c) = \overline{Lim}_{x\to c}f(x) = \limsup_{r\to 0^{+}}\{f(s):s\in D, |s-c|< r\}
      \]
      Logo, existe \(r_{c}\) tal que \(\sup\{f(s):s\in D, |s-c|<r_{c}\} < k.\) Assim,
      \(f(s)<k\) para todo s em \((c-r_{c}, c+r_{c})\) e em D. Finalmente, por conta disso,
      \(c\in D\cap(c-r_{c}, c+r_{c})\subseteq{f^{-1}((-\infty, k))}\) e \(\bigcup_{c\in f^{-1}((-\infty, k))}^{}{D\cap (c-r_{c}, c+r_{c})} = f^{-1}(-\infty, k)\)
  É claro que, para todo x em \(O_{i}\cap D = f^{-1}((-\infty, k))\).
\end{proof*}

\subsection{Motiva\c cão para Diferenciabilidade}
 A diferenciabilidade é um conceito fundamental na análise matemática, que
estende o conceito de derivada de uma função em um ponto. Uma função é dita
diferenciável em um ponto se existe um número que pode ser interpretado como
a inclinação da reta tangente à curva representativa da função naquele ponto.

 Seja $f: \mathbb{R} \to \mathbb{R}$ uma função. Dizemos que $f$ é 
diferenciável em $x_0 \in \mathbb{R}$ se o seguinte limite existe:

\begin{equation}
f'(x_0) = \lim_{{h \to 0}} \frac{f(x_0 + h) - f(x_0)}{h}
\end{equation}

O número $f'(x_0)$ é chamado de derivada de $f$ em $x_0$.

A diferenciabilidade tem diversas aplicações em matemática e ciências. 
Ela permite, por exemplo, calcular a taxa de variação de uma quantidade, 
otimizar funções e resolver equações diferenciais. Alguns exemplos incluem

\subsubsection{Função Linear}

Considere a função linear $f(x) = mx + b$, onde $m$ e $b$ são constantes. A derivada de $f$ em qualquer ponto $x_0$ é $f'(x_0) = m$, que é a inclinação da reta.

\subsubsection{Função Quadrática}

Considere a função quadrática $f(x) = ax^2 + bx + c$, onde $a$, $b$ e $c$ são constantes. A derivada de $f$ em qualquer ponto $x_0$ é $f'(x_0) = 2ax_0 + b$.

Pense na diferenciabilidade como uma maneira de medir o quão "suave" é uma função. Se você pode desenhar a curva de uma função sem levantar a caneta do papel, então a função é provavelmente diferenciável. A derivada em um ponto é simplesmente a inclinação da reta que melhor se ajusta à curva naquele ponto.
Vamos entrar mais a fundo nessas ideias

\subsection{Derivadas}
\begin{def*}
  Seja \(f:D_{f}\rightarrow \mathbb{R}\) uma fun\c cão e \(p\in D_{f}\) um ponto de
  acumula\c cão de \(D_{f}\). Se existir o limite 
    \[
      \lim_{x\to p}\frac{f(x)-f(p)}{x-p} = L\in \mathbb{R},
    \]
    diremos que L é a derivada de f em p e escrevemos 
      \[
        f'(p) = L = \lim_{x\to p}\frac{f(x)-f(p)}{x-p}=\lim_{h\to 0}\frac{f(p+h)-f(p)}{h}. \square
      \]
\end{def*}

  Se \(f:D\rightarrow \mathbb{R}\) possui derivada num ponto d de D que é também um
  ponto de acumula\c cão de D, para \(h\in \mathbb{R}\) tal que d+h pertence a D, escrevemos 
    \[
      r(h) = f(d+h)-f(d)-f'(d)h.
    \]
    Nesses pontos, definimos \(r:\{h\in \mathbb{R}: d+h\in D_{f}\}\rightarrow \mathbb{R}\)
    e escrevemos \(f(d+h) = f(d)+f'(d)h + r(h)\) e, fazendo \(\sigma (h) = \frac{r(h)}{h}, h\neq 0,\)
    temos \(\lim_{h\to 0}\sigma (h)=0.\)

    Note que f é diferenciável em d se, e somente se, existe fun\c cão \(\sigma \)
    com \(\lim_{h\to 0}\sigma (h)=0\) tal que \(f(d+h)=f(d)+[f'(d)+\sigma (h)]h\).
    \begin{def*}
      Sejam \(f:D_{f}\rightarrow \mathbb{R}\) uma fun\c cão e \(p\in D_{f}\) um ponto
      de acumula\c cão à direita de \(D_{f}\). Se existir o limite 
        \[
          \lim_{x\to p^{+}}\frac{f(x)-f(p)}{x-p}=L^+\in \mathbb{R},
        \]
        diremos que \(L^+\) é a derivada à direita de f em p e escrevemos 
          \[
            f'(p^+) = L^+=\lim_{x\to p^+}\frac{f(x)-f(p)}{x-p}= \lim_{h\to 0^{+}}\frac{f(p+h)-f(p)}{h}
          \]
          Analogamente, define-se derivada à esquerda. \(\square\)
    \end{def*}
      Já definimos a derivada de \(f:D_{f}\rightarrow \mathbb{R}\) em pontos p de
      \(D_{f}\) que também são pontos de acumula\c cão de \(D_{f}\). Sendo assim, 
      se 
        \[
          D_{f'}= \biggl\{x\in D_{f}:x\text{ é um ponto de acumula\c cão de $D_{f}$ e }\lim_{h\to 0}\frac{f(x+h)-f(x)}{h}\text{ existe.}\biggr\}\subseteq{D_{f}},
        \]
    definimos a fun\c cão \(f':D_{f'}\rightarrow \mathbb{R}\) por 
      \[
        f'(x) = \lim_{h\to 0}\frac{f(x+h)-f(x)}{h}, \quad x\in D_{f'}.
      \]
    A fun\c cão f' é dita a fun\c cão derivada ou apenas derivada de f.
    \hypertarget{diff_cont}{
     \begin{theorem*}
       Se f for diferenciável em p de \(D_{f}\), então f será contínua em p.
     \end{theorem*}
    }

   \begin{proof*}
     Recorde que \(p\in D_{f}\) é um ponto de acumula\c cão de \(D_{f}\).
     Vamos mostrar que \(\lim_{x\to p}f(x)=f(p)\) ou que \(\lim_{x\to p}(f(x)-f(p))=0.\)
     Escrevemos 
       \[
         f(x)-f(p) = \frac{f(x)-f(p)}{x-p}(x-p).
       \]
       Assim, 
         \[
           \lim_{x\to p}(f(x)-f(p)) = \lim_{x\to p}\frac{f(x)-f(p)}{x-p}\lim_{x\to p}(x-p) = f'(p)\cdot 0 = 0.
         \]

         Portanto, f é contínua em p.
   \end{proof*}
   Note que a recíproca não vale. A fun\c cão \(f(x) = |x|\) é contínua em x, mas
   não é diferenciável em x=0.
   \hypertarget{negative_diff}{   \begin{example}
        Se f não é contínua em p, então f não é diferenciável em p.
      \end{example}}
   \begin{example}
     A fun\c cão 
       \[
         f(x) = \left\{\begin{array}{ll}
             x^2,\quad x\geq 2\\
             2,\quad x\leq 2
           \end{array}\right.
       \]
       não é contínua, logo não diferenciável.
   \end{example}
   \begin{def*}
     Seja f derivável em \(D_{f'}\). A fun\c cão \(f':D_{f'}\rightarrow \mathbb{R}\)
     é dita derivada de f ou derivada primeira de f. Definimos a derivada segunda de f
     como 
       \[
         (f')'(x) = \lim_{h\to 0}\frac{f(x+h)-f'(x)}{h}
       \]
       quando o limite existe. Escrevemos \(f''=f^{(2)} = (f')'\) para denotá-la.

       Para \(n \in\mathbb{N}^{\times}\), podemos definir a derivada n-ésima analogamente. \(\square\).
   \end{def*}
   \begin{theorem*}
     Se \(k\in \mathbb{R}\) e \(n\in \mathbb{N}^{\times}\), valem as propriedades
     de derivada vistas em cálculo I. Além disso, valem as regras
     da derivada do produto, quociente, soma e multiplica\c cão por constante.
   \end{theorem*}
   Lembrando: 
     \[
       (f(x)g(x))' = f'(x)g(x) + f(x)g'(x)\quad \biggl(\frac{f}{g}\biggr)' = \frac{f'(x)g(x) - g'(x)f(x)}{g(x)^2}
     \]

     \hypertarget{chain_rule}{
      \begin{theorem*}
        Sejam \(f:D_{f}\rightarrow \mathbb{R}, g:D_{g}\rightarrow \mathbb{R}\) diferenciáveis com 
        \(Im(g)\subseteq{D_{f}.}\) Se g é diferenciável em p, g(p) é ponto de acumula\c cão de \(D_{f}\),
        f é diferenciável em g(p) e \(h=f\circ{g}\), então h é diferenciável em p e 
          \[
            h'(p) = f'(g(p))g'(p).
          \]
      \end{theorem*}
     }
    \begin{proof*}
       Fa\c ca q = g(p). Sejam \(\sigma _{g}\) e \(\sigma _{f}\) definidas em vizinhan\c cas
       de 0 com \(\lim_{h\to 0}\sigma_{g}(h) = 0\) e \(\lim_{k\to 0}\sigma _{f}(k) = 0\) tais que 
      \begin{align*}
        &g(p+h) = g(p) + [g'(p) + \sigma _{g}(h)]h\\
        &f(q+k) = f(q) + [f'(q) + \sigma _{f}(k)]k.
      \end{align*}
      Fazendo \(k=g(p+h)-g(p) = [g'(p)+\sigma _{g}(h)]h\), temos \(g(p+h) = q + k\) e
     \begin{align*}
       f(g(p+h)) = f(q+k) &= f(q) + [f'(q) + \sigma _{f}(k)]k\\
                          &= f(q) + [f'(q) + \sigma _{f}(k)][g'(p)+\sigma _{g}(h)]h\\
                          &=f(g(p))+f'(f(p))g'(p)h + [\sigma _{f}(g(p+h)-g(p))[g'(p)+\sigma _{g}(h)]+f'(q)\sigma _{g}(h)]h.
     \end{align*}
 Agora, se \(\sigma _{f\circ{g}}(h) =[\sigma _{f}(g(p+h)-g(p))[g'(p)+\sigma _{g}(h)]+f'(q)\sigma _{g}(h)]h\),
     temos \(\lim_{h\to 0}\sigma _{f\circ{g}}(h)=0.\) \qedsymbol
    \end{proof*}
    Seja \(f:D_{f}\rightarrow \mathbb{R}\) uma fun\c cão que tem inversa, \(D_{f^{-1}}=Im(f)\)
    e \(f^{-1}:D_{f^{-1}}\rightarrow \mathbb{R}\). Então, para todo x em \(D_{f^{-1}}\), 
      \[
        f(f^{-1}(x)) = x.
      \]
    Vimos que se f é contínua em um compacto, \(f^{-1}\) é contínua. Se, além disso,
    \(f \) e \(f^{-1}\) forem deriváveis, pela \hyperlink{chain_rule}{Regra da Cadeia,}
      \[
        f'(f^{-1}(x))(f^{-1})'(x) = x' = 1.
      \]
    Logo, \(f'(f^{-1}(x))\neq0\) e 
      \[
        (f^{-1})'(x) = \frac{1}{f'(f^{-1}(x))}.
      \]
    Esse estudo sobre a inversa é sumarizado no próximo resultado: 
   \begin{theorem*}
    Seja f injetiva, p um ponto de acumula\c cão de Im(f). Se f for diferenciável em \(q=f^{-1}(p)\) e
    \(f^{-1}\) é contínua em p, então \(f^{-1}\) é diferenciável em p se, e somente se,
    \(f'(f^{-1}(p))\neq 0.\) Neste caso, 
      \[
    (f^{-1})'(x) = \frac{1}{f'(f^{-1}(x))}.
      \]
   \end{theorem*}
   \begin{proof*}
     Se \(f'(f^{-1}(p))\neq0\), como \(f^{-1}\) é contínua em p, \(\lim_{h\to 0}f^{-1}(p+h)=f^{-1}(p)\).
     Usando \(f(f^{-1}(x)), x\in D_{f^{-1}}\), temos 
    \begin{align*}
      (f^{-1})'(p) &= \lim_{h\to 0}\frac{f^{-1}(p+h)-f^{-1}(p)}{h} = \lim_{h\to 0}\frac{1}{\frac{h}{f^{-1}(p+h)-f^{-1}(p)}}\\
                   &= \lim_{h\to 0}\frac{1}{\frac{f(f^{-1}(p+h))-f(f^{-1}(p))}{f^{-1}(p+h)-f^{-1}(p)}} = \frac{1}{f'(f^{-1}(p))}
    \end{align*} 
    Por outro lado, se \(f^{-1}\) é diferenciável em p, da \hyperlink{chain_rule}{regra da cadeia}
    aplicada a \(f\circ{f^{-1}}\), temos \(f'(f^{-1}(p))\cdot (f^{-1})'(p)=1\) e \(f'(f^{-1}(p))\neq0\).
   \end{proof*}
   \begin{example}
     Se \(g(x)=x^{\frac{1}{n}}\), então \(g'(x) = \frac{1}{n}x^{\frac{1}{n}-1}, 2\leq n\in \mathbb{N}\). 
     Lembre-se que, \(x > 0\) se n for par e \(x\neq0\) se n for ímpar.
    
     Com efeito, note que \(g(x) = x^{\frac{1}{n}} = f^{-1}(x)\), em que \(f(u) = u^n.\) Então, 
       \[
         g'(x) = (f^{-1})(x) = \frac{1}{f'(f^{-1}(x))} = \frac{1}{n(x^{\frac{1}{n}})^{n-1}} = \frac{1}{n(x^{\frac{n-1}{n}})} = \frac{1}{n}x^{\frac{1}{n}-1}.
       \]
   \end{example}
\newpage

\section{Aula 19 - 12/05/2023}
\subsection{O que esperar?}
\begin{itemize}
  \item Fun\c cões Deriváveis em Intervalos; 
  \item Teorema de Darboux;
  \item Descontinuidades da Derivada.
\end{itemize}

\subsection{Fun\c cões Deriváveis em Intervalos}
 \begin{example}
   Seja \(f:\mathbb{R}\rightarrow \mathbb{R}\) dada por \(f(x) = x^{3}\). Já 
sabemos que f é contínua. Como 
\begin{align*}
  f(x) - f(y) &=(x-y)(x^{2}+xy+y^{2})\\
              &=(x-y)((x+\frac{y}{2})^{2}+\frac{3}{4}y^{2})\\
              &=(x-y)((\frac{x}{2}+y)^{2}+\frac{3}{4}x^{2}).
\end{align*}
  Segue que \(x\neq y\) implica \(f(x)\neq f(y)\) e f é injetora. Além disso,
\(\lim_{x\to \pm\infty}f(x)=\pm\infty\)e, pelo teorema do valor intermediário,
f é uma bije\c cão de \(\mathbb{R}\) em \(\mathbb{R}\). Além disso,
\(f^{-1}:\mathbb{R}\rightarrow \mathbb{R}\) é contínua pois ela é contínua em
qualquer intervalo compacto. A inversa de f é denotada \(f^{-1}(x) = x^{\frac{1}{3}} = \sqrt[3]{x}.\)
  Do teorema sobre a derivada da inversa e do fato que \(f'(x) = 3x^{2}\) 
deduzimos que \(f^{-1}:\mathbb{R}\rightarrow \mathbb{R}\) é diferenciável se,
e somente se, \(x\in \mathbb{R}\backslash\{0\}\) e, para estes valores de x. 
  \[
    \overbrace{(\sqrt[3]{x})'}^{(f^{-1})'(x)} = \frac{1}{f'(f^{-1}(x))} = \frac{1}{3(\sqrt[3]{x})^{2}} = \frac{1}{3x^{\frac{2}{3}}}.
  \]
\end{example}
 \begin{def*}
   Seja \(f:D\rightarrow \mathbb{R}\). Dizemos que f tem um máximo (mínimo) local no
   ponto de em D se existe \(\delta >0\) tal que \(f(x)\leq f(d) (f(x)\geq f(d))\)
   para todo x em D com \(|x-d| < \delta .\) Quando a desigualdade é estrita,
   dizemos que f tem um máximo (mínimo) local estrito. Os máximos e mínimos locais
   serão chamados de valores extremos e os pontos em que a fun\c cão assume valores máximos
   ou mínimos são chamados de pontos de máximo ou de mínimo. \(\square\)
 \end{def*}
  Pela defini\c cão de derivada (à direita), segue que 
 \begin{enumerate}
 \item Se \(f:D\rightarrow \mathbb{R}\) é não-decrescente (não-crescente) e é
diferenciável em um ponto d de D, então \(f'(d)\geq 0 (f'(d)\leq 0)\). Vale o mesmo
para fun\c cões diferenciáveis à direita.
 \item Se \(f:D\rightarrow \mathbb{R}\) é derivável à direita (esquerda) em um ponto
\(d\in D\) e \(f'(d^{+}) > 0 (f'(d^{-}) > 0)\) então existe \(\delta >0\) tal que
\(x\in D, x\in (d, d+\delta )(x\in (d-\delta , d)\) implica \(f(x) > f(d) (f(x) < f(d)).\)
 \item Se \(f:D\rightarrow \mathbb{R}\) é derivável à direita (esquerda) em um 
ponto \(d\in D\) e \(f(d^{-}) < 0 (f'(d^{-}) < 0)\), então existe \(\delta >0\) tal que
\(x\in D, x\in(d, d+\delta )(x\in (d-\delta, d))\) implica \(f(x) < f(d) (f(x) > f(d)).\)
  \item Se \(f:D\rightarrow \mathbb{R}\) é derivável em um ponto \(d\in D\),
d sendo ponto de acumula\c cão à direita e à esquerda e \(f'(d) > 0\), existe \(\delta  > 0\)
tal que \(x\in D, d-\delta < x < d < y < d+\delta\) implica \(f(x) < f(d) < f(y).\)
\item \(f:D\rightarrow \mathbb{R}\) é derivável em um ponto d em D, d é um ponto
de acumula\c cão à direita e à esquerda e f tem um valor extremo local em d, então \(f'(d) = 0.\)
 \end{enumerate}
\begin{def*}
  Seja I um intervalo e \(f:I\rightarrow \mathbb{R}\) uma fun\c cão diferenciável.
Se \(f':I\rightarrow \mathbb{R}\) for contínua diremos que f é continuamente diferenciável
em I, ou simplesmente f é de classe \(C^{1}\) em I.
\end{def*}
  Note que existe fun\c cão diferenciável em um intervalo I que não é 
continuamente diferenciável.
 \begin{example}
   Seja \(f:\mathbb{R}\rightarrow \mathbb{R}\) dada por 
     \[
       f(x) = \left\{\begin{array}{ll}
           x^{2}\sin{(\frac{1}{x})},\quad x\neq0\\
           f(0) = 0.
         \end{array}\right.
     \]
     Então, f é diferenciável, mas f' não é contínua em x = 0.
 \end{example}
 \hypertarget{darboux}{
\begin{theorem*}
  Se \(f:[a, b]\rightarrow \mathbb{R}\) é diferenciável com \(f'(a)\neq f'(b)\), então,
  para todo C entre \(f'(a)\) e \(f'(b)\), existe \(c\in (a, b)\) tal que 
  \(f'(c) = C\).
\end{theorem*}}
\begin{proof*}
  Suponha que \(f'(a) < 0 < f'(b)\). Segue que, para x próximo a a em \([a, b], f(x) < f(a)\)
e, para x próximo a b em \([a, b], f(x) < f(b).\) Logo, o ponto de mínimo (que existe pelo
\hyperlink{weierstrass}{Teorema de Weierstrass} c de f ocorre em \((a, b)\) e 
portanto \(f'(c) = 0.\) Para o caso geral, consideramos 
\begin{itemize}
  \item Se \(f'(a) < C < f'(b), g(x) = f(x) - C \cdot x\)
  \item Se \(f'(a) > C > f'(b), g(x) = C \cdot x - f(x).\)
\end{itemize}
\end{proof*}
\begin{theorem*}
  Se I é um intervalo e \(f:I\rightarrow \mathbb{R}\) é diferenciável, então
  f' não tem descontinuidades de primeira espécie.
\end{theorem*}
\begin{proof*}
  Se a é um ponto de acumula\c cão à direita de I e \(L^{+}=\lim_{x\to a^{+}}f(x)\)
existe, mostremos que \(L^{+} = f'(a).\)

  De modo análogo (exercício), se a é um ponto de acumula\c cão à esquerda de I 
e \(L^{-}=\lim_{x\to a^{-}}f'(x)\) existe, mostre que \(f'(a) = L^{-}.\)

  Se \(L^{+} > f'(a)\) e \(C\in(f'(a), L^{+})\), existe \(\delta >0\) tal que
 \(f'(x) > C\) para todo \(x\in(a, a-\delta ).\) Escolhendo \(b\in(a, a+\delta )\),
 temos \(f'(b) > C > f'(a)\), uma contradi\c cão com o \hyperlink{darboux}{Teorema de Darboux},
 pois este implica a existência de \(c\in(a, b)\) tal que \(f'(c) = C.\) Logo,
 \(f'(a)\geq L^{+}.\)

  Se \(f'(a) > L^{+}\) e \(C\in(L^{+}, f'(a))\), existe \(\delta >0\) tal que 
 \(f'(x) < C\) para todo x em \((a, a+\delta )\). Escolhendo b em \((a, a+\delta )\),
 temos \(f'(b) < C < f'(a)\), novamente contradizendo o \hyperlink{darboux}{Teorema de Darboux},
 visto que ele implica que deve existir c em \((a, b)\) tal que \(f'(c) = C.\) Assim,
 \(f'(a)\leq L^{+}.\) Portanto, \(L^{+} = f'(a).\) \qedsymbol
\end{proof*}
\begin{theorem*}
  Se \(f:[a, b)\rightarrow \mathbb{R}\) é contínua e diferenciável à direita com 
derivada à direita, \(D^{+}f:[a, b)\rightarrow \mathbb{R}.\) Se \(D^{+}f(x)\leq 0 (D^{+}f(x)\geq 0)\)
para todo x em \([a, b)\) e \(f(a) = 0\), então \(f(x)\leq 0 (f(x)\geq 0)\) em \([a, b).\)
\end{theorem*}
\begin{proof*}
  Suponha primeiramente que \(D^{+}f(x) < 0\) para todo \(x\in[a, b).\) Se o resultado
é falso, existe ao menos um \(x\in(a, b)\) tal que \(f(x) > 0.\) Seja
\(x_{0}=\inf\{x\in(a, b): f(x) > 0\}.\)

  Da continuidade de f, \(f(x_{0}) = 0\) e da defini\c cão de \(x_{0}\) existe 
uma sequência \(x_{n}\in(x_{0}, b)\) tal que \(x_{n}\overbracket[0pt]{\longrightarrow}^{n\to \infty}x_{0}.\)
Assim, 
  \[
    D^{+}f(x_{0}) = \lim_{n\to \infty}\frac{f(x_{n})-f(x_{0})}{x_{n}-x_{0}}\geq 0,
  \]
uma contradi\c cão. Logo, \(f(x)\leq 0\) para todo x em \([a, b)\).

  Agora, consideramos o caso geral \(D^{+}f(x)\leq 0\) para todo \(x\in[a, b).\)
Neste caso, consideramos a fun\c cão auxiliar \(f_{\varepsilon }(x) = f(x) - \varepsilon (x-a)\) 
e temos \(f_{\varepsilon }(x)\leq 0\) para todo \(x\in[a, b)\) e \(\varepsilon >0.\)
Disto segue que, para todo \(x\in[a, b), f(x)\leq 0\). O outro caso será deixado como
exercício. \qedsymbol
\end{proof*}
\begin{example}
  Encontre uma fun\c cão \(f:\mathbb{R}\rightarrow \mathbb{R}\) que é diferenciável
  à direita, tal que \(D^{+}f(x) < 0\) para todo \(x\neq0, D^{+}f(0)=0,\) f é positiva
  para \(x > 0\) e negativa para \(x < 0 (f(x)\overbracket[0pt]{\longrightarrow}^{x\to \pm\infty}\pm\infty.)\)
\end{example}
 \begin{crl*}
   Se \(f:[a, b)\rightarrow \mathbb{R}\) é contínua e diferenciável à direita com derivada
à direita \(D^{+}f:[a, b)\rightarrow \mathbb{R}.\) Se \(D^{+}f(x)\leq 0\) para todo x em \([a, b)\),
então f é não-crescente em \([a, b).\)
 \end{crl*}
\begin{proof*}
  Se \(a\leq c < d < b,\) seja \(g:[c, b)\rightarrow \mathbb{R}\) definida por 
 \(g(x) = f(x) - f(x)\) e \(D^{+}g(x)\leq 0\) para todo x em \([c, b)\). Segue do 
 teorema que \(g(x)\leq 0\) para todo x em \([c, b).\) Em particular, \(g(d) = f(d) - f(c)\leq 0\). \qedsymbol
\end{proof*}
\begin{crl*}
  Se \(f:[a, b)\rightarrow \mathbb{R}\) é contínua e diferenciável à direita com
derivada à direita \(D^{+}f:[a, b)\rightarrow \mathbb{R}.\) Se \(D^{+}f(x)\geq 0\)
para todo x em \([a, b)\), então f é não-decrescente em \([a, b).\)
\end{crl*}
  A prova deste exercício é deixada como exercício. Além disso, como exercício,
enuncie e prove resultados semelhantes aos anteriores para a derivada à esquerda.
\newpage

\section{Aula 20 - 15/05/2023}
\subsection{O que esperar?}
\begin{itemize}
  \item Teorema do Valor Médio;
  \item Consequências do TVM;
  \item Regra de L'Hospital;
  \item Teorema de Taylor.
\end{itemize}
\subsection{Teorema do Valor Médio}
  O resultado a seguir não pode ser completado até o futuro, no entanto, pode
ser útil já afirmá-lo.
\hypertarget{incomplete_1}{\begin{crl*}
  Se \(f:[a, b)\rightarrow \mathbb{R}\) é contínua e diferenciável à direita com 
  derivada à direita \(D^{+}f:[a, b]\rightarrow \mathbb{R}\) contínua, então f é de classe
  \(C^{1}\).
\end{crl*}}
\begin{proof*}
  Seja \(g = D^{+}f\) e defina 
    \[
      h(t) = f(a) + \int_{a}^{t} g(\tau )d\tau.
    \]
  A fun\c cão h é continuamente diferenciável em \([a, b)\). Se \(\varphi (t) = h(t)-f(t)\),
  então \(\varphi (a) = 0\) e \(D^{+}\varphi (t) = 0\) em [a, b). Do teorema anterior,
  \(\varphi (t)\leq 0\) em \([a, b)\).

  Como \(-\varphi (t)\) também satisfaz as hipóteses do teorema anterior, \(\varphi (t)\geq 0\).
  Portanto, \(\varphi = 0\), ou seja, f = h em \([a, b)\). \qedsymbol
\end{proof*}
  O Teorema do Valor Médio é um dos teoremas fundamentais da análise. Segue sua formula\c cão segundo
  Cauchy
  \hypertarget{cauchy_mvt}{
 \begin{theorem*}
   Se \(f, g:[a, b]\rightarrow \mathbb{R}\) são fun\c cões contínuas que são
   diferenciáveis em \((a, b)\), existe c em \((a, b)\) tal que 
     \[
       [f(b)-f(a)]g'(c) = [g(b)-g(a)]f'(c).
     \]
 \end{theorem*}}
  Um de seus corolários é o teorema de Rolle
  \hypertarget{Rolle}{
 \begin{crl*}
   Se \(f:[a, b]\rightarrow \mathbb{R}\) é contínua em \([a, b]\) e diferenciável
 em \((a, b)\) e \(f(a)=f(b),\) então existe um c em (a, b) tal que \(f'(c) = 0.\)
 \end{crl*}}
\begin{proof*}
  Segue fazendo \(g(x) = 1\) no teorema anterior. \qedsymbol
\end{proof*}
 Outro corolário é o Teorema do Valor Médio em sua versão mais simplificada e mais usada
 \hypertarget{mean_value}{
\begin{crl*}
  Se \(f:[a, b]\rightarrow \mathbb{R}\) é contínua em \([a, b]\) e diferenciável em
  \((a, b)\), então existe \(c\in (a, b)\) tla que 
    \[
      f(b) - f(a) = f'(c)(b-a),
    \]
    ou seja, 
      \[
        f'(c) = \frac{f(b)-f(a)}{b-a}.
      \]
\end{crl*}}
\begin{proof*}
  Basta tomar \(g(x) = x\) no teorema anterior. \qedsymbol
\end{proof*}
  Vamos provar o \hyperlink{cauchy_mvt}{teorema do valor médio de Cauchy}.
 \begin{proof*}
   Se 
     \[
       h(x) = [f(b)-f(a)]g(x) - [g(b)-g(a)]f(x),\quad (a\leq x\leq b),
     \]
  então h é contínua em \([a, b]\), diferenciável em \((a, b)\) e 
    \[
      h(a) = f(b)g(a) - f(a)g(b) = h(b).\quad \text{(para verificar isso, basta abrir as contas)}
    \] 
  Para provar o teorema, basta mostrarmos que \(h'(c) = 0\) para algum c em \((a, b)\).

Se h é constante, isto vale para todo c em \((a, b).\) 
Se\(h(x) > h(a)\) para algum \( x\in (a, b),\) seja c um ponto de \([a, b]\) 
no qual h atinge seu máximo. Como \(h(a) = h(b),\) segue que c pertence a \((a, b)\) e \(h'(c) = 0\).

Se \(h(x) < h(a)\) para algum x em \((a, b)\), escolhemos c em \([a, b]\)
para o qual h atinge seu mínimo. Portanto, assim como antes, \(c\in (a, b)\) e
 \(f'(c) = 0.\) \qedsymbol
 \end{proof*}

  Algumas outras consequências do \hyperlink{mean_value}{TVM} são:
 \begin{crl*}
  \begin{itemize}
    \item[a)] Se \(f'(x)\geq  0,\) então f é não-decrescente em \((a, b)\);
      \item[b)] Se \(f'(x) > 0,\) então f é crescente em \((a, b)\)
        \item[c)]  Se \(f'(x) = 0,\) então f é constante em \((a, b)\)
          \item[d)] Se \(f'(x)\leq  0,\) então f é não-crescente em \((a, b)\)
            \item[e)] Se \(f'(x)\leq  0,\) então f é decrescente em \((a, b)\)
  \end{itemize}
 \end{crl*}
 Uma observa\c cão, também, é o Teorema da Fun\c cão Inversa
\begin{theorem*}
  Se \(I\subseteq{\mathbb{R}}\) é um intervalo aberto, \(x_{0}\) um ponto de I,
  \(f:I\rightarrow \mathbb{R}\) uma fun\c cão \(C^{1}\) e \(f'(x_{0})\neq 0\), então
  existe \(\delta >0\) tal que \(f:(x_{0}-\delta, x_{0}+\delta )\rightarrow \mathbb{R}\) 
  é injetora, \(f((x_{0}-\delta ,x_{0}+\delta )) = J\) é um intervalo e \(f^{-1}:J\rightarrow I\)
  é continuamente diferenciável com 
    \[
      (f^{-1})'(x) = \frac{1}{f'(f^{-1}(x))}.
    \]
\end{theorem*}
  O resultado a seguir é conhecido como Regra de L'Hospital, um dos resultados
  mais úteis no cálculo de limites indeterminados de fun\c cões.
\hypertarget{Lhospital}{ 
  \begin{theorem*}
   Sejam \(f, g\) diferenciáveis em \((a, b)\) e \(g'(x)\neq0\) para todo
   x em \((a, b)\), em que \(-\infty\leq a < b\leq +\infty\) e 
     \[
       \frac{f'(x)}{g'(x)}\overbracket[0pt]{\longrightarrow}^{x\to a}A.
     \]
  Se 
    \[
      f(x)\overbracket[0pt]{\longrightarrow}^{x\to a}0\text{ e }g(x)\overbracket[0pt]{\longrightarrow}^{x\to a}0
    \]
  ou se 
    \[
      g(x)\overbracket[0pt]{\longrightarrow}^{x\to a}+\infty,
    \]
  então 
    \[
      \frac{f(x)}{g(x)}\overbracket[0pt]{\longrightarrow}^{x\to a}A.
    \]
  O resultado permanece válido se \(x\overbracket[0pt]{\longrightarrow}b\), ou se
  \(g(x)\longrightarrow+\infty.\)
 \end{theorem*}
}
 \begin{proof*}
   Primeiramente, consideramos o caso \(-\infty\leq A < +\infty\). Se \(q > r > A\),
   existe c em \((a, b)\) tal que, se \(a < x < c\), então 
     \[
       \frac{f'(x)}{g'(x)} < r.
     \]
    Se \(a < x < y < c,\) do \hyperlink{cauchy_mvt}{Teorema do Valor Médio de Cauchy}, existe 
    \(t\in (x, y)\) tal que 
      \[
        \frac{f(x) - f(y)}{g(x) - g(y)} = \frac{f'(t)}{g'(t)} < r,
      \] 
    visto que, se \(\lim_{x\to a}\frac{f'(x)}{g'(x)} = A\), fazendo \(q > r > A, q\in \mathbb{R},\)
    deve existir \(c > a\) tal que, para todo x em \((a, c)\), 
    \[
      \frac{f'(x)}{g'(x)} < r.
    \]

    Se a primeira condi\c cão vale, então, fazendo \(x\rightarrow a\) na desigualdade acima, 
      \[
        \frac{f(y)}{g(y)}\leq r < q \quad (a < y < c),
      \]

    Caso a segunda condi\c cão seja verdadeira, mantendo y fixado na equa\c cão em que aplicamos o TVM de Cauchy anterior,
    podemos escolher \(c_{1}\in (a, y)\) tal que \(g(x) > g(y) \) e \(g(x) > 0\) se 
    \(a < x < c_{1}.\) Multiplicando ela por \(\frac{g(x)-g(y)}{g(x)},\) obtemos 
      \[
        \frac{f(x)}{g(x)} < r - r\frac{g(y)}{g(x)} + \frac{f(y)}{g(x)}\quad (a < x < c_{1})
      \]
    Fazendo x tender para a nesta desigualdade, segue que existe \(c_{2}\in (a, c_{1})\)
    tal que 
      \[
        \frac{f(x)}{g(x)} < q\quad (a < x < c_{2}).
      \]
    Assim, temos para qualquer \(q > A\), existe \(c_{2}\) tal que \(\frac{f(x)}{g(x)} < q\) se \(a < x < c_{2}.\)
      \[
        \lim_{x\to a^{+}} \frac{f(x)}{g(x)}\leq q \quad\forall q > A.
      \]
      Do mesmo modo, se \(-\infty < A\leq +\infty\) e \(p < A,\) podemos encontrar
    \(c_{3}\) tal que 
      \[
        p < \frac{f(x)}{g(x)}\quad (a < x < c_{3}).
      \]
      Disto segue o resultado. \qedsymbol
 \end{proof*}
  O próximo resultado é conhecido como Teorema de Taylor, e permite-nos quebrar
  uma fun\c cão com várias derivadas em um polinômio com derivadas mais simples.
  Ele pode ser visto como uma generaliza\c cão do \hyperlink{mean_value}{TVM}.
  \hypertarget{taylor}{
 \begin{theorem*}
   Se \(n\in \mathbb{N}^{\times}, f:[a, b]\rightarrow \mathbb{R}\) é uma fun\c cão
   n-1 vezes diferenciável em \([a, b]\) e n-vezes diferenciável em \((a, b)\) com
   \(f^{(n-1)}:[a, b]\rightarrow \mathbb{R}\) contínua. Sejam \(\alpha , \beta \in[a, b]\) diferentes
   e 
     \[
       P(t) = \sum\limits_{k=0}^{n-1}\frac{f^{(k)}(\alpha )}{k!}(t-\alpha )^{k}.
     \]
    Então, existe \(\xi\) entre \(\alpha  \) e \(\beta \) tal que 
      \[
        f(\beta ) = P(\beta ) + \frac{f^{(n)}(\xi)}{n!}(\beta -\alpha )^{n}.
      \]
    Para n=1, este é o \hyperlink{mean_value}{TVM}. Em geral, o teorema mostra
    como aproximar f por polinômios e fornece uma maneira de estimar o erro
    se conhecermos limita\c cões para \(|f^{(n)}(\xi)|.\)
 \end{theorem*}}
 \begin{proof*}
   Seja M o número definido por 
     \[
       f(\beta ) = P(\beta ) + M(\beta -\alpha )^{n}.
     \]
    Fazendo 
      \[
        g(t) = f(t) - P(t) - M(t-\alpha )^{n}\quad (a\leq t\leq b).
      \]
    Precisamos mostrar que \(n!M = f^{(n)}(\xi)\) para algum \(\xi\) entre 
  \(\alpha \) e \(\beta \). Segue que 
    \[
      g^{(n)}(t) = f^{(n)}(t) - n!M\quad (a < t < b).
    \]
    Resta mostrar que \(g^{(n)}(\xi) = 0\) para algum \(\alpha \leq \xi\leq \beta .\) 
    Como \(P^{(k)}(\alpha ) = f^{(k)}(\alpha ), k=0, \cdots, n-1,\) temos 
      \[
        g(\alpha ) = g'(\alpha ) = \cdots = g^{(n-1)}(\alpha ) = 0.
      \]
    Nossa escolha de M implica que \(g(\beta ) = 0\) e, do \hyperlink{mean_value}{Teorema
    do Valor Médio,} \(g'(x_{1}) = 0\) para algum \(x_{1}\) entre \(\alpha \) e \(\beta \).
    Como \(g'(\alpha ) = 0\), de modo semelhante, \(g''(x_{2}) = 0\) para algum
    \(x_{2}\) entre \(\alpha \) e \(x_{1} \). Portanto, depois de n, chegamos à conclusão
    que \(g^{(n)}(x_{n}) = 0\) para algum \(x_{n}\) entre \(\alpha \) e \(x_{n-1}\), isto é,
    entre \(\alpha \) e \(\beta \). \qedsymbol
 \end{proof*}
 \subsection{Fun\c cões Convexas}
\begin{def*}
  Seja I um intervalo. Uma fun\c cão \(f:I\rightarrow \mathbb{R}\) é convexa quando,
  dados \(a < x < b\) em I, o ponto \((x, f(x))\) fica abaixo da reta que
  liga os pontos \((a, f(a))\) e \((b, f(b))\). \(\square\)
\end{def*}
A equa\c cão da reta é 
  \[
    y = \frac{f(b)-f(a)}{b-a}(x-a)+f(a) \text{ ou } y = \frac{f(b)-f(a)}{b-a}(x-b)+f(b).
  \]
  Logo, \(f:I\rightarrow \mathbb{R}\) é convexa se uma das desigualdades 
    \[
      \frac{f(x) - f(a)}{x-a}\leq \frac{f(b) - f(a)}{b-a}\leq \frac{f(b)-f(x)}{b-x}
    \]
    está sempre satisfeita sempre que \(a < x < b\) em I. 
\hypertarget{convex_characterization}{ \begin{theorem*}
   Seja \(I\subseteq{\mathbb{R}}\) um intervalo e \(f:I\rightarrow \mathbb{R}\) duas vezes diferenciável.
   Então, f é convexa se, e somente se, \(f''(x)\geq 0\) para todo x em I.
 \end{theorem*}}
\newpage
\section{Aula 21 - 17/05/2023}
\subsection{O que esperar?}
\begin{itemize}
  \item Fun\c cões convexas;
  \item Séries de Taylor;
  \item Fun\c cões analíticas;
  \item Fun\c cões de varia\c cão limitada.
\end{itemize}
\subsection{Fun\c cões Convexas}
  Vamos relembrar o que foi visto na última aula sobre fun\c cões convexas antes de mais nada.
  \begin{def*}
  Seja I um intervalo. Uma fun\c cão \(f:I\rightarrow \mathbb{R}\) é convexa quando,
  dados \(a < x < b\) em I, o ponto \((x, f(x))\) fica abaixo da reta que
  liga os pontos \((a, f(a))\) e \((b, f(b))\). \(\square\)
\end{def*}
A equa\c cão da reta é 
  \[
    y = \frac{f(b)-f(a)}{b-a}(x-a)+f(a) \text{ ou } y = \frac{f(b)-f(a)}{b-a}(x-b)+f(b).
  \]
  Logo, \(f:I\rightarrow \mathbb{R}\) é convexa se uma das desigualdades 
    \[
      \frac{f(x) - f(a)}{x-a}\leq \frac{f(b) - f(a)}{b-a}\leq \frac{f(b)-f(x)}{b-x}
    \]
    está sempre satisfeita sempre que \(a < x < b\) em I. 
\begin{theorem*}
   Seja \(I\subseteq{\mathbb{R}}\) um intervalo e \(f:I\rightarrow \mathbb{R}\) duas vezes diferenciável.
   Então, f é convexa se, e somente se, \(f''(x)\geq 0\) para todo x em I.
 \end{theorem*}
 Come\c camos agora por provar esse resultado.
 \begin{proof*}
   Se \(f''(x)\geq 0\) para todo x em I, então dados a, a+h em I, existe
c entre a e a+h tal que \(f(a+h) = f(a) + f'(a)\cdot h + \frac{f''(c)}{2}\cdot h^{2}.\)
Como \(f''(c)\geq 0, f(a+h)\geq f(a) + f'(a)\cdot h.\) Disto segue que 
 \(\frac{f(a+h)-f(a)}{h}\leq f'(a)\) se \(h < 0\) e \(\frac{f(a+h)-f(a)}{h}\geq f'(a)\)
 se \(h > 0.\) Em outras palavras, se \(a < x < b\) em I, então \(\frac{f(x)-f(a)}{x-a}\leq \frac{f(b)-f(x)}{b-x}\)
 ou seja, 
   \[
     (f(x)-f(a))(b-x)\leq (f(b)-f(x))(x-a).
   \]
   Deste modo, 
  \begin{align*}
    (f(x)-f(a))(b-a-(x-a)) &\leq (f(b)-f(a)-(f(x)-f(a)))(x-a)\text{ e }\\
                           &(f(x)-f(a))(b-a)\leq (f(b)-f(a))(x-a).
  \end{align*}
  Com isso, provamos a convexidade de f.

  Reciprocamente, se f é convexa, dados \(a < x < b\) em I, temos 
    \[
      \frac{f(x)-f(a)}{x-a}\leq \frac{f(b)-f(a)}{b-a}\leq \frac{f(x) - f(b)}{x-b}.
    \]
  Fazendo \(x\rightarrow a\) e \(x\rightarrow b\), 
    \[
      f'(a)\leq \frac{f(b)-f(a)}{b-a}\leq f'(b)
    \]
  e f é não-decrescente em I. Portanto, \(f''(x)\geq 0\) para todo x em I. \qedsymbol
 \end{proof*}
  Observe que, se f é diferenciável, então \(f'\) é crescente se, e somente se,
f é convexa. Além disso, podemos mostrar analogamente que, se \(f''(x) > 0\) para
todo x em I, então f é estritamente convexa em I. A recíproca é falsa, no entanto.
Basta tomar \(f(x) = x^{4}\), a qual é estritamente convexa em \(\mathbb{R}\), mas
 \(f''(0) = 0.\)

  Seja \(f:I\rightarrow \mathbb{R}\) de classe \(C^{\infty}.\) Se \(a, x\in I^{o},\) então
podemos escrever, para todo \(k\in \mathbb{N},\) 
  \[
  f(x) = f(a) - f'(a)(x-a) + \frac{f''(a)}{2!}(x-a)^{2} + \cdots + \frac{f^{(n-1)}(a)}{(n-1)!}(x-a)^{n-1} + r_{n}((x-a)),
  \]
  em que \(r_{n}((x-a))=\frac{f^{(n)}((1-\theta_{n})a+\theta_{n}x)}{n!}(x-a)^{n},\) com
 \(0 < \theta_{n} < 1.\) A série 
   \[
     \sum\limits_{n=0}^{\infty}\frac{f^{(n)}(a)}{n!}(x-a)^{n}
   \]
chama-se a série de Taylor da fun\c cão f em torno do ponto a. Esta série
pode ou não convergir e, mesmo que convirja, sua soma pode ser diferente
de f(x).
\begin{example}
  Seja \(f:\mathbb{R}\rightarrow \mathbb{R}\) definida por \(f(0)=0\) e 
 \(f(x)=e^{-\frac{1}{x^{2}}}\) se \(x\neq0.\) Mostre que f é \(C^{\infty}, f^{(n)}(0)=0\)
 para todo n natural e portanto a série de Taylor de f em \(x=0\) é convergente para
 \(f(0)\), mas não coincide com f para nenhum \(x\neq0.\)
\end{example}
\begin{def*}
  Se \(I\subseteq{\mathbb{R}}\) é um intervalo aberto e \(f:I\rightarrow \mathbb{R}\)
é uma fun\c cão, dizemos que f é analítica em I se, para cada \(a\in I\), existe
 \(\varepsilon >0\) tal que a série de Taylor \(\sum\limits_{n=0}^{\infty}\frac{f^{(n)(a)}}{n!}\cdot (x-a)^{n}\)
 é convergente com soma \(f(x)\) para todo x em \((a-\varepsilon, a+\varepsilon ).\square\)
\end{def*}
  É claro que a série de Taylor \(\sum\limits_{}^{}\frac{f^{(n)}(a)}{n!}\cdot (x-a)^{n}\) converge
  para f(x) se, e somente se, \(\lim_{n\to a}r_{n}((x-a))=0.\)
 \begin{example}
   A série de Taylor da fun\c cão seno é dada por 
     \[
       \sin{(x)} = \sum\limits_{n=0}^{\infty}\frac{(-1)^{n}}{(2n+1)!}x^{2n+1}.
     \]
 \end{example}
  Veremos mais tarde que, se a série de potências \(\sum\limits_{n=1}^{\infty}a_{n}(x-a)^{n}\)
tem raio de convergência \(R>0\), então a fun\c cão definida por 
  \[
    f(x) = \sum\limits_{n=0}^{\infty}a_{n}(x-a)^{n}, \quad x\in(a-R, a+R)
  \]
  é analítica.
\subsection{Fun\c cões de Varia\c cão Limitada (BV)}
  Se r é um número real, colocamos \(r^{+}=\max\{r, 0\}, r^{-}=\max\{-r, 0\}\).
 \begin{def*}
   Uma cole\c cão \(\{a_{0}, \cdots, a_{k}\}\) de pontos em \([a, b]\) é chamada
uma parti\c cão do intervalo \([a, b]\) se \(a=a_{0} < a_{1} < a_{2} < \cdots<a_{k}=b.\)
Seja \(f:[a, b]\rightarrow \mathbb{R}\) e \(\{a_{0},\cdots,a_{k}\}\) uma parti\c cão
de \([a, b]\). Escrevemos 
 \begin{align*}
   &p=\sum\limits_{n=1}^{k}[f(a_{i})-f(a_{i-1})]^{+},\\
   &n=\sum\limits_{i=1}^{k}[f(a_{i})-f(a_{i-1})]^{-},\\
   &t=\sum\limits_{i-1}^{k}[f(a_{i})-f(a_{i-1})]=p+n,\\
   &f(b)-f(a)=p-n.
 \end{align*}
 \end{def*}
 Sejam 
\begin{align*}
  &P_{a}^{b}=\sup\{p:k\in \mathbb{N}, a=a_{0} < a_{1} < \cdots < a_{k}=b\text{ parti\c cão de [a, b]}\}
  &N_{a}^{b}=\sup\{n:k\in \mathbb{N}, a=a_{0} < a_{1} < \cdots < a_{k}=b\text{ parti\c cão de [a, b]}\}
  &T_{a}^{b}=\sup\{t:k\in \mathbb{N}, a=a_{0} < a_{1} < \cdots < a_{k}=b\text{ parti\c cão de [a, b]}\}
\end{align*}
  Dizemos que \(P_{a}^{b}, N_{a}^{b}\) e \(T_{a}^{b}\) são as varia\c cões positiva,
  negativa e total de f. É claro que 
    \[
      \max\{P_{a}^{b}, N_{a}^{b}\}\leq T_{a}^{b}\leq P_{a}^{b} + N_{a}^{b}\quad\text{ e } f(b)-f(a) = P_{a}^{b} - N_{a}^{b}.
    \]
\begin{def*}
  A fun\c cão \(f:[a, b]\rightarrow \mathbb{R}\) é de varia\c cão limitada se \(T_{a}^{b} < \infty.\) Denotamos
isso por \(f\in BV([a, b]). \square\)
\end{def*}
\begin{theorem*}
 \begin{itemize}
\item[1)] Se \(f:[a, b]\rightarrow \mathbb{R}\) é Lipschitz contínua, então \(f:[a, b]\rightarrow \mathbb{R}\)
  é de varia\c cão limitada.
\item[2)] Se \(f:[a, b]\rightarrow \mathbb{R}\) é monótona, então \(f:[a, b]\rightarrow \mathbb{R}\) é
  de varia\c cão limitada.
\item[3)] Se \(f:[a, b]\rightarrow \mathbb{R}\) é de varia\c cão limitada, existem
  fun\c cões não-decrescente \(g, h:[a, b]\rightarrow \mathbb{R}\) tais que \(f(x) = g(x)-h(x).\)
 \end{itemize}
\end{theorem*}
\begin{proof*}
  \(1. \Rightarrow )\) Se f é Lipschitz, \(\max\{P_{a}^{b}, N_{a}^{b}\leq T_{a}^{b}\leq L(b-a) < \infty,\}\)
em que \(L > 0\) é a constante de Lipschitz.

  \(2. \Rightarrow )\) Se f é monótona, então \(T_{a}^{b} = [f(b)-f(a)] < \infty.\)

  \(3. \Rightarrow )\) Se \(T_{a}^{b} < \infty,\) defina \(g, h:[a, b]\rightarrow \mathbb{R}\) por 
 \(g(x) = f(a) + P_{a}^{x}\) e \(h(x) = N_{a}^{x}\) para cada x de \([a, b]\). Portanto,
 \(g, h\) são não-decrescentes e \(f(x) = g(x) - h(x).\) \qedsymbol
\end{proof*}
\newpage
\section{Aula 22 - 19/05/2023}
\subsection{O que esperar?}
\begin{itemize}
  \item Medida exterior;
  \item Continuidade Lipschitz;
\end{itemize}
\subsection{Medida Exterior}
\begin{def*}
  Se \(I=(a, b)\), defina \(l(I) = b-a.\) Dado \(A\subseteq{\mathbb{R}}\), existe
  uma família contável de intervalos abertos que cobrem A. Seja \(\mathcal{U}_{A}\)
  a cole\c cão de todas as coberturas contáveis de intervalos abertos de A. Então,
  definimos a medida exterior de A como 
    \[
      m^{*}(A) = \inf\{\sum\limits_{}^{}l(I_{n}):\{I_{n}\}\in \mathcal{U}_{A}\}.\square
    \]
\end{def*}
  É claro que \(m^{*}(\emptyset) = 0, m^{*}((a, b))\leq b-a, m^{*}(\{x\}) = 0\) para todo
  x real e que, se \(A\subseteq{B}, m^{*}(A)\leq m^{*}(B).\)
 \begin{lemma*}
   Temos \(m^{*}[a, b]=m^{*}(a, b]=m^{*}[a, b)=m^{*}(a, b) = b-a.\)
 \end{lemma*}
\begin{proof*}
  Como \([a, b]\subseteq{(a-\frac{\varepsilon }{2}, b+\frac{\varepsilon }{2})}\), para todo \(\varepsilon >0\), temos
  \(m^{*}([a, b])\leq b-a.\) Por outro lado, se \(\{I_{n}\}\in \mathcal{U}_{[a, b]}\) existe
  uma subcole\c cão finita \(\{I_{n_{1}}, \cdots, I_{n_{k}}\}\) de \(\{I_{n}\}\) tal que
  \(\bigcup_{i=1}^{k}{I_{n_{i}}}\supseteq{[a, b]}\) e 
    \[
      \sum\limits_{i=1}^{k}l(I_{n_{i}})\leq \sum\limits_{}^{}l(I_{n}).
    \]
    Podemos tomar uma subcobertura de \(\{I_{n_{1}}, \cdots, I_{n_{k}}\}\) de forma que
    \(a\in I_{n_{j_{1}}} = (y_{1}, x_{1})\) e, recursivamente, se \(x_{j}\leq b, x_{j}\in I_{n_{i_{j+1}}} = (y_{j+1},x_{j+1})\)
    parando para \(j_{0}\leq k\) tal que \(I_{n_{i_{j_{0}}}}\ni b\). Assim, como
    \(y_{j} < x_{j-1} < x_{j},\) 
      \[
        \sum\limits_{j=1}^{j_{0}}l(I_{n_{i_{j}}}) = \sum\limits_{}^{}(x_{j} - y_{j}) > x_{j_{0}} - y_{1} > b - a
      \]
      e \(m^{*}([a, b]) = b-a.\) \qedsymbol
\end{proof*}
 \begin{lemma*}
   Se \(\{A_{n}\}\) é uma família contável de subconjuntos de \(\mathbb{R}\), então 
     \[
       m^{*}\biggl(\bigcup_{}^{}{A_{n}}\biggr)\leq \sum\limits_{}^{}m^{*}(A_{n}).
     \]
 \end{lemma*}
\begin{proof*}
  Só temos que considerar o caso em que \(\sum\limits_{}^{}m^{*}(A_{n}) < \infty.\) Como
  \(m^{*}(A_{n})\) é finita, dado \(\varepsilon  > 0\), seja \(\{I_{n, i}\}_{i}\in \mathcal{U}_{A_{n}}\) tal que
  \(A_{n}\subseteq{\bigcup_{i}^{}{I_{n, i}}}\) e \(\sum\limits_{i}^{}l(I_{n, i}) < m^{*}(A_{n}) + 2^{-n}\varepsilon\)
  Logo, \(\{I_{n, i}\}_{n, i}\in \mathcal{U}_{\bigcup_{}^{}{A_{n}}}\) e 
    \[
      m^{*}\biggl(\bigcup_{n}^{}{A_{n}}\biggr)\leq \sum\limits_{n, i}^{}l(I_{n, i}) = \sum\limits_{n}^{}\sum\limits_{i}^{}l(I_{n, i}) < \sum\limits_{n}^{}(m^{*}(A_{n}) + \varepsilon 2^{-n}).
    \]
  Portanto, \(m^{*}\biggl(\bigcup_{n}^{}{A_{n}}\biggr) = \sum\limits_{}^{}m^{*}(A_{n}) + \varepsilon\) e, como \(\varepsilon \) é arbitrário, o resultado segue.
\end{proof*}
\begin{crl*}
 \begin{itemize}
   \item[1)] Se \(A\subseteq{\mathbb{R}}\) é enumerável, \(m^{*}(A) = 0.\)
     \item[2)] Se \(\{A_{n}\}\) é uma família contável de subconjuntos de \(\mathbb{R}\) e
       \(m^{*}(A_{n}) = 0\) para todo n, então \(m^{*}(\bigcup_{n}^{}A_{n} = 0.\)
 \end{itemize}
\end{crl*}
\begin{example}
  Seja \(I_{j} = [a_{j}, b_{j}], 1\leq j\leq n\) com \(b_{j} < a_{j+1}, 1\leq j\leq n-1.\) Mostre que 
    \[
      m^{*}\biggl(\bigcup_{j=1}^{n}{I_{j}}\biggr) = \sum\limits_{j=1}^{n}(b_{j}-a_{j}).
    \]
\end{example}
  Dado um conjunto E, uma cobertura de Vitali dele é uma cole\c cão de intervalos
  \(\mathcal{I}\) tal que para todo x em E e \(\varepsilon  > 0\), existe \(I\in \mathcal{I}\) tal que
  \(x\in I\) e \(m^{*}(I) < \varepsilon .\) 
\hypertarget{vitali_covering}{
\begin{lemma*}
  Seja \(E\subseteq{[a, b]}\), tal que \(m^{*}(E)\leq b-a.\) Se \(\mathcal{I}\)
  é uma cobertura de E por intervalos não degenerados (Não é um ponto só) e tal que, dados \(x\in E\) e
  \(\varepsilon >0\), existe \(I\in \mathcal{I}\) tal que \(x\in I\) e \(l(I) < \varepsilon .\)
  Então, dado \(\varepsilon  > 0\), exsite uma cole\c cão finita e disjunta \(\{I_{1}, \cdots, I_{N}\}\subseteq{\mathcal{I}}\)
  tal que 
    \[
      m^{*}\biggl(E\biggl\backslash\biggr. \bigcup_{n=1}^{N}{I_{n}}\biggr) < \varepsilon .
    \]
\end{lemma*}}
\begin{proof*}
  Basta considerar o caso com cada intervalo de \(\mathcal{I}\) fechado. Podemos assumir que
  \(\mathcal{I}\ni I \subseteq{\mathcal{O}=(b-1, b+1)}\) e que \(I\cap E \neq\emptyset\) para todo I
  de \(\mathcal{I}\). 

  Escolhemos uma sequência \(\{I_{n}\}\) de intervalos disjuntos de \(\mathcal{I}\) da
  seguinte forma: Seja \(I_{1}\in \mathcal{I}\) qualquer e se \(I_{1}, \cdots, I_{n}\) já foram
  escolhidos, seja \(r_{n}\) o supremo dos comprimentos dos intervalos de \(\mathcal{I}\) que não interseptam
  nenhum dos \(I_{1}, \cdots, I_{n}\). Claramente, \(r_{n} <  l(\mathcal{O}).\) Se 
  \(E\not\subseteq{\bigcup_{i=1}^{n}{I_{i}}},\) encontramos \(I_{n+1}\in \mathcal{I}\)
  disjunto de \(I_{1}, \cdots, I_{n}\) e tal que \(l(I_{n+1} > \frac{1}{2}r_{n}\).

  Assim, \(\{I_{n}\}\) é uma sequência disjunta de intervalos em \(\mathcal{I}\) e, como
  \(\bigcup_{}^{}{I_{n}\subseteq{\mathcal{O}}}, \sum\limits_{}^{}l(I_{n})\leq l(\mathcal{O}).\) Logo,
  existe \(N\in \mathbb{N}\) tal que 
    \[
      \sum\limits_{N+1}^{\infty} l(I_{n}) < \frac{\varepsilon }{5}.
    \]
  Seja 
    \[
      R = E\biggl\backslash\bigcup_{n=1}^{N}{I_{n}}\biggr.
    \]
  Mostraremos que \(m^{*}(R) < \varepsilon .\) Se \(x\in R,\) como \(F = \bigcup_{n=1}^{N}{I_{n}}\) é
  fechado e \(x\not\in F,\) existe \(I\) em \(\mathcal{I}, x\in I\) e \(I\cap F=\emptyset.\)

  Agora, se \(I\cap I_{i} = \emptyset\) para \(i\leq \kappa \), temos \(l(I)\leq r_{\kappa } < 2l(I_{\kappa +1}.\)
  Como \(\lim_{\kappa \to \infty}l(I_{\kappa }) = 0,\) o intervalo I deve intersectar
  pelo menos um dos intervalos \(I_{\kappa }.\)

  Seja n o menor inteiro tal que \(I\cap I_{n} \neq\emptyset\). Claramente \(n > N\), e
  \(l(I)\leq r_{n-1} < 2l(I_{n}).\) Como \(x\in I\) e \(I\cap I_{n} \neq\emptyset\), a distância de
  x ao ponto médio de \(I_{n}\) é no máximo \(l(I) + \frac{1}{2}l(I_{n}) < \frac{5}{2} l(I_{n}).\)

  Logo, x pertence ao intervalo \(J_{n}\) tendo o mesmo ponto médio que \(I_{n}\) e 
  o quíntuplo do comprimento. Desta forma, 
    \[
      R\subseteq{\bigcup_{N+1}^{\infty}{J_{n}}}
    \]
  e, portanto, 
    \[
      m^{*}(R)\leq \sum\limits_{n=N+1}^{\infty}l(J_{n}) = 5\sum\limits_{n=N+1}^{\infty}l(I_{n}) < \frac{5\varepsilon }{5} = \varepsilon .\text{ \qedsymbol}
    \]
\end{proof*}
\begin{lemma*}
  Se \(f:[a, b]\rightarrow \mathbb{R}\) é monótona, então f é diferenciável exceto 
  possivelmente em um conjunto \(E\subseteq{[a, b]}\) com \(m^{*}(E) = 0\).
\end{lemma*}
\begin{proof*}
  Faremos apenas o caso f não-decrescente. Considere 
 \begin{align*}
   &\overline{d^{+}}f(x) = \limsup_{h\to 0^{+}} \frac{f(x+h)-f(x)}{h}\text{ e }\overline{d^{-}}f(x) = \limsup_{h\to 0^{+}}\frac{f(x) - f(x-h)}{h}\\
   &\underline{d^{+}}f(x) = \liminf_{h\to 0^{+}} \frac{f(x+h)-f(x)}{h}\text{ e }\underline{d^{-}}f(x) = \liminf_{h\to 0^{+}}\frac{f(x) - f(x-h)}{h}\\
 \end{align*}
 Provemos que o conjunto dos \(x\in[a, b]\) tais que \(\underline{d^{-}}f(x) < \overline{d^{+}}f(x)\) 
 ou \(\overline{d^{-}}f(x) < \underline{d^{+}}f(x)\) tem medida exterior nula.
\end{proof*}
\newpage

\section{Aula 23 - 22/05/2023}
\subsection{O que esperar?}
\begin{itemize}
  \item Continuando medida exterior;
  \item Relacionando medida e derivada;
\end{itemize}
\end{document}
